\section{Proof of Theorem~\ref{thm: u decay}}
\label{app: elliptic}

This appendix proves Theorem~\ref{thm: u decay}: the exponential decay of the Kolmogorov solution $u$ and of its derivatives of orders up to $4$. The decay of $u$ and of its derivatives of orders up to $3$ is proved in \citep[Theorem~1]{ye2026mse}, so essentially only the derivative of order $4$ needs to be estimated.\footnote{The paper \citep{ye2026mse} is an earlier work of the first author, Xuda Ye, and is under review at the IMA Journal of Numerical Analysis. Beyond it, the proofs of this paper use no self-cited results that are not yet peer-reviewed.} The bound \eqref{eq: poisson bound} on $\phi_h$ then follows by summation over the time steps. Appendix~\ref{sec: variations} writes out the variation processes and their pathwise bounds. Appendix~\ref{sec: semigroup} states the estimates on the semigroup $e^{t\mathcal L}$ that the proof uses and proves the short-time bound \eqref{eq: short time}. For a function $g$ with bounded derivatives of orders $1$ to $4$, Appendix~\ref{sec: fourth} bounds $\nabla^4 (e^{t\mathcal L} g)$ by differentiating the derivative formula \eqref{eq: BEL} $3$ times. Appendix~\ref{sec: elliptic conclusion} combines this bound with the decay \eqref{eq: grad decay} and the short-time bound \eqref{eq: short time}.

The proof is written for the potential $U$ and the test function $f$ of Assumptions~\ref{ass: U} and \ref{ass: f}, and it uses the derivatives of $U$ up to order $5$. Throughout this appendix, the constant $C$ depends only on $m$, $M$, $R$ and $d$, unless otherwise specified.

\subsection{Variation processes}
\label{sec: variations}

Let $X_t$ be the solution of \eqref{eq: langevin} started at $X_0 = x$, and write $\partial_w$ for the derivative with respect to the initial state $x$ along a vector $w$. For a function $g$ on $\mathbb R^d$, the derivatives of $e^{t\mathcal L} g$ in the initial state are expressed through the variation processes. The variation processes of orders $1$ to $4$ are
\begin{equation}
    J_t := \partial_x X_t , \qquad
    K_t[w] := \partial_w J_t , \qquad
    \Theta_t[w,w'] := \partial_{w'} K_t[w] , \qquad
    \Xi_t[w,w',w''] := \partial_{w''} \Theta_t[w,w'] ,
    \label{eq: variation processes}
\end{equation}
where the variations of orders $2$ to $4$ are taken along vectors $w, w', w'' \in \mathbb R^d$, so that, like $J_t$, they take values in $\mathbb R^{d \times d}$. From the definition \eqref{eq: variation processes},
\begin{equation*}
\begin{gathered}
    J_t w = \partial_w X_t , \qquad K_t[w] w' = \partial_w \partial_{w'} X_t , \\
    \Theta_t[w,w'] w'' = \partial_w \partial_{w'} \partial_{w''} X_t , \qquad \Xi_t[w,w',w''] w''' = \partial_w \partial_{w'} \partial_{w''} \partial_{w'''} X_t ,
\end{gathered}
\end{equation*}
hence $K_t[w] w'$, $\Theta_t[w,w'] w''$ and $\Xi_t[w,w',w''] w'''$ are symmetric in all their directions. Since the noise of \eqref{eq: langevin} is additive, differentiating the dynamics in its initial state gives ordinary differential equations along the path, with the initial values $J_0 = I_d$ and $K_0 = \Theta_0 = \Xi_0 = 0$:
\begin{align}
    \frac{\D J_t}{\D t} &= -\nabla^2 U(X_t)\, J_t , \label{eq: first variation}\\
    \frac{\D K_t[w]}{\D t} &= -\nabla^2 U(X_t)\, K_t[w] - \nabla^3 U(X_t)[J_t w]\, J_t , \label{eq: second variation}\\
    \frac{\D \Theta_t[w,w']}{\D t} &= -\nabla^2 U(X_t)\, \Theta_t[w,w'] - \nabla^3 U(X_t)[J_t w]\, K_t[w'] - \nabla^3 U(X_t)[J_t w']\, K_t[w] \notag\\
    &\qquad - \nabla^3 U(X_t)\bigl[ K_t[w] w' \bigr]\, J_t - \nabla^4 U(X_t)[J_t w, J_t w']\, J_t . \label{eq: third variation}
\end{align}
These equations are derived in \citep[Equations~(A.1), (A.8) and~(A.9)]{ye2026mse}. Differentiating \eqref{eq: third variation} once more, along $w''$, gives the equation of the variation $\Xi_t$ of order $4$,
\begin{align}
    & \frac{\D \Xi_t[w,w',w'']}{\D t} = -\nabla^2 U(X_t)\, \Xi_t[w,w',w'']
    - \sum_{(2,1)} \nabla^3 U(X_t)[J_t w'']\, \Theta_t[w,w']
    - \sum_{(2,1)} \nabla^3 U(X_t)\bigl[ K_t[w] w' \bigr]\, K_t[w''] \notag \\
    &\qquad - \nabla^3 U(X_t)\bigl[ \Theta_t[w,w'] w'' \bigr]\, J_t
    - \sum_{(2,1)} \nabla^4 U(X_t)[J_t w, J_t w']\, K_t[w''] \notag \\
    &\qquad - \sum_{(2,1)} \nabla^4 U(X_t)\bigl[ K_t[w] w', J_t w'' \bigr]\, J_t
    - \nabla^5 U(X_t)[J_t w, J_t w', J_t w'']\, J_t ,
    \label{eq: variations}
\end{align}
where, for an expression $F$ of the $3$ directions in which the underlined pair forms a block, the sum runs over the partitions of $\{w,w',w''\}$ into a pair and a singleton,
\begin{equation}
    \sum_{(2,1)} F(\underline{w,w'}, w'') := F(\underline{w,w'}, w'') + F(\underline{w,w''}, w') + F(\underline{w',w''}, w).
    \label{eq: sum21}
\end{equation}
In each sum $\sum_{(2,1)}$ of \eqref{eq: variations}, the summand written out has the pair $\{w,w'\}$ and the singleton $w''$.

The bounds of the next lemma are Gr\"onwall estimates, which grow exponentially in $t$. The factor $t$ in the bounds on $K_t$, $\Theta_t$ and $\Xi_t$ is what the proof of Lemma~\ref{lem: fourth from lipschitz} uses at short times.

\begin{lemma}
\label{lem: variation bounds}
    Let Assumption~\ref{ass: U} hold. Then, pathwise, for all $t \Ge 0$ and unit vectors $w, w', w''$,
    \begin{equation}
    \begin{gathered}
        \norm{J_t} \Le e^{M t} , \qquad
        \norm{K_t[w]} \Le M t\, e^{2M t} , \qquad
        \norm{\Theta_t[w,w']} \Le M(1+3M t)\, t\, e^{3M t} , \\
        \norm{\Xi_t[w,w',w'']} \Le C (1+t)^2\, t\, e^{4M t} ,
    \end{gathered}
    \label{eq: variation bounds}
    \end{equation}
    with $C$ depending only on $M$.
\end{lemma}

\begin{proof}
    The bounds on $J_t$, $K_t$ and $\Theta_t$ are \citep[Lemma~6]{ye2026mse}. For $\Xi_t[w,w',w'']$, the matrix $J_t J_s^{-1}$ solves \eqref{eq: first variation} on $[s,t]$ from $I_d$, so $\norm{J_t J_s^{-1}} \Le e^{M(t-s)}$ by Gr\"onwall's inequality, and Duhamel's principle with $\Xi_0 = 0$ gives
    \begin{equation*}
        \Xi_t[w,w',w''] = \int_0^t J_t J_s^{-1} \bigg( \frac{\D \Xi_s[w,w',w'']}{\D s} + \nabla^2 U(X_s)\, \Xi_s[w,w',w''] \bigg) \D s .
    \end{equation*}
    The factor $\frac{\D \Xi_s[w,w',w'']}{\D s} + \nabla^2 U(X_s)\, \Xi_s[w,w',w'']$ of the integrand is the sum of the $14$ terms of \eqref{eq: variations} other than $-\nabla^2 U(X_s)\, \Xi_s[w,w',w'']$. The bounds \eqref{eq: variation bounds} on $J_s$, $K_s$ and $\Theta_s$, with $\norm{\nabla^i U} \Le M$ for $i = 3, 4, 5$, bound these terms group by group: the $5$ summands of the next display bound, in this order, the sum carrying $\nabla^3 U$ and $\Theta_s$, the sum carrying $\nabla^3 U$ and $2$ factors $K_s$, the term carrying $\nabla^3 U$ and $\Theta_s[w,w'] w''$, the $2$ sums carrying $\nabla^4 U$ together, and the term carrying $\nabla^5 U$:
    \begin{equation*}
    \begin{aligned}
        & \quad \biggl\| \frac{\D \Xi_s[w,w',w'']}{\D s} + \nabla^2 U(X_s)\, \Xi_s[w,w',w''] \biggr\| \\
        & \Le 3M^2 (1+3Ms)\, s\, e^{4Ms} + 3M^3 s^2 e^{4Ms} + M^2 (1+3Ms)\, s\, e^{4Ms} + 6M^2 s\, e^{4Ms} + M e^{4Ms} \\
        & = M \bigl( 1 + 10Ms + 15M^2 s^2 \bigr) e^{4Ms} \Le 15M (1+Ms)^2 e^{4Ms} .
    \end{aligned}
    \end{equation*}
    Hence
    \begin{equation*}
        \norm{\Xi_t[w,w',w'']} \Le \int_0^t e^{M(t-s)}\, 15M (1+Ms)^2 e^{4Ms} \,\D s \Le 15M (1+Mt)^2\, t\, e^{4Mt} \Le C (1+t)^2\, t\, e^{4Mt} .
    \end{equation*}
\end{proof}

\subsection{Semigroup estimates}
\label{sec: semigroup}

The estimates of this subsection hold for every $g \in C^4(\mathbb R^d)$ with bounded derivatives of orders $1$ to $4$; the rate $\lambda > 0$ depends only on $m$, $M$, $R$ and $d$. Lemmas~\ref{lem: semigroup}, \ref{lem: BEL} and \ref{lem: variance bound} are proved in \citep[Appendix~A]{ye2026mse} under the bounds of Assumption~\ref{ass: U} of orders up to $4$, and are restated without proof, since the proof of Lemma~\ref{lem: fourth from lipschitz} uses them explicitly. Lemma~\ref{lem: short time} extends the short-time bounds of \citep[Equations~(A.11) and~(A.12)]{ye2026mse} to the derivative of order $4$ and is proved here.

\begin{lemma}
\label{lem: semigroup}
    Let Assumption~\ref{ass: U} hold. Then, for all $t \Ge 0$,
    \begin{equation}
        \bigl\| \nabla (e^{t\mathcal L} g) \bigr\|_\infty \Le C e^{-\lambda t} \norm{\nabla g}_\infty ;
        \label{eq: grad decay}
    \end{equation}
    and for $0 < t \Le 1$,
    \begin{equation}
        \bigl\| \nabla^2 (e^{t\mathcal L} g) \bigr\|_\infty \Le C\, t^{-\frac12}\, \norm{\nabla g}_\infty , \qquad
        \bigl\| \nabla^3 (e^{t\mathcal L} g) \bigr\|_\infty \Le C\, t^{-1}\, \norm{\nabla g}_\infty .
        \label{eq: smoothing 23}
    \end{equation}
\end{lemma}

The decay \eqref{eq: grad decay} is \citep[Equation~(A.7)]{ye2026mse}, obtained from the reflection coupling of the overdamped Langevin dynamics \citep{eberle2016reflection}. The smoothing bounds \eqref{eq: smoothing 23} are \citep[Lemmas~8 and~9]{ye2026mse}, obtained by differentiating the derivative formula \eqref{eq: BEL} below up to $2$ times in the initial state.

The derivative formula \citep{elworthy1994formulae}, in the form of \citep[Lemma~7]{ye2026mse}, expresses the derivative of $e^{2\tau\mathcal L} g$ along a vector $v$ through a stochastic integral of the variation process $J_s$ of order $1$.

\begin{lemma}
\label{lem: BEL}
    Let Assumption~\ref{ass: U} hold. Then, for $\tau > 0$, $\psi := e^{\tau\mathcal L} g$ and $v \in \mathbb R^d$,
    \begin{equation}
        v \cdot \nabla (e^{2\tau\mathcal L} g)(x) = \frac{1}{\sqrt2\,\tau}\, \mathbb E^x \biggl[ \psi(X_\tau) \int_0^\tau (J_s v) \cdot \D B_s \biggr] , \qquad \forall x \in \mathbb R^d .
        \label{eq: BEL}
    \end{equation}
\end{lemma}

The next lemma, \citep[Equation~(A.17)]{ye2026mse} in the proof of \citep[Lemma~8]{ye2026mse}, bounds the variance of $\psi(X_\tau)$.

\begin{lemma}
\label{lem: variance bound}
    Let Assumption~\ref{ass: U} hold. Then, for $\tau > 0$, $\psi := e^{\tau\mathcal L} g$ and $X_0 = x \in \mathbb R^d$,
    \begin{equation}
        \Var\bigl( \psi(X_\tau) \bigr) \Le 2 d\, \tau\, e^{4M\tau}\, \norm{\nabla g}_\infty^2 .
        \label{eq: variance bound}
    \end{equation}
\end{lemma}
The randomness of $\psi(X_\tau)$ comes from the Brownian motion $(B_t)_{t \Ge 0}$ on $[0,\tau]$, which drives $X_t$ from $X_0 = x$.

The last estimate of this subsection bounds the derivatives of $e^{t\mathcal L} g$ over a finite time; it covers the times $t \Le 1$ in the proof of Theorem~\ref{thm: u decay} and supplies the derivative bounds needed in the proof of Lemma~\ref{lem: fourth from lipschitz} and in the case $t \Ge 1$ of that theorem.

\begin{lemma}
\label{lem: short time}
    Let Assumption~\ref{ass: U} hold. There is $C$, depending only on $M$, such that for every $g \in C^4(\mathbb R^d)$ with bounded derivatives of orders $1$ to $4$ and all $t \Ge 0$,
    \begin{equation}
        \bigl\| \nabla^j (e^{t\mathcal L} g) \bigr\|_\infty \Le C (1+t)^3 e^{4M t} \sum_{i=1}^{j} \norm{\nabla^i g}_\infty , \qquad j = 1,2,3,4 .
        \label{eq: short time}
    \end{equation}
\end{lemma}

\begin{proof}
    For $j \Le 3$, \eqref{eq: short time} follows from the bounds \citep[Equations~(A.11) and~(A.12)]{ye2026mse},
    \begin{equation*}
    \begin{gathered}
        \bigl\| \nabla (e^{t\mathcal L} g) \bigr\|_\infty \Le e^{Mt} \norm{\nabla g}_\infty , \qquad
        \bigl\| \nabla^2 (e^{t\mathcal L} g) \bigr\|_\infty \Le e^{2Mt} \norm{\nabla^2 g}_\infty + M t\, e^{2Mt} \norm{\nabla g}_\infty , \\
        \bigl\| \nabla^3 (e^{t\mathcal L} g) \bigr\|_\infty \Le e^{3Mt} \norm{\nabla^3 g}_\infty + 3M t\, e^{3Mt} \norm{\nabla^2 g}_\infty + M (1+3Mt)\, t\, e^{3Mt} \norm{\nabla g}_\infty ,
    \end{gathered}
    \end{equation*}
    whose coefficients are each at most $C (1+t)^3 e^{4Mt}$.

    For $j = 4$, fix $x \in \mathbb R^d$ and unit vectors $w, w', w'', w'''$. By \eqref{eq: variation processes}, the derivatives of $X_t$ in its initial state along these vectors are $J_t w$, $K_t[w] w'$, $\Theta_t[w,w'] w''$ and $\Xi_t[w,w',w''] w'''$, together with the same expressions along the other subsets of $\{w,w',w'',w'''\}$. The chain rule for higher derivatives applied to $g(X_t)$, and the exchange of the derivatives with the expectation, justified by the pathwise bounds \eqref{eq: variation bounds} and the bounded derivatives of $g$, give
    \begin{align}
        & \nabla^4 (e^{t\mathcal L} g)(x)[w,w',w'',w'''] = \mathbb E^x \biggl[  \nabla^4 g(X_t)[J_t w, J_t w', J_t w'', J_t w'''] + \sum_{(2,1,1)} \nabla^3 g(X_t)\bigl[ K_t[w] w', J_t w'', J_t w''' \bigr] \notag \\
        & \quad + \sum_{(3,1)} \nabla^2 g(X_t)\bigl[ \Theta_t[w,w'] w'', J_t w''' \bigr] + \sum_{(2,2)} \nabla^2 g(X_t)\bigl[ K_t[w] w', K_t[w''] w''' \bigr] \notag \\
        & \quad + \nabla g(X_t) \cdot \Xi_t[w,w',w''] w''' \biggr] ,
        \label{eq: fourth chain rule}
    \end{align}
    where, for an expression $F$ of the $4$ directions in which each underlined group of directions forms a block, the sums run over the partitions of $\{w,w',w'',w'''\}$ into blocks of the sizes in the subscript:
    \begin{align*}
        \sum_{(2,1,1)} F(\underline{w,w'}, w'', w''') &:= F(\underline{w,w'}, w'', w''') + F(\underline{w,w''}, w', w''') + F(\underline{w,w'''}, w', w'') \\[-8pt]
        & \quad + F(\underline{w',w''}, w, w''') + F(\underline{w',w'''}, w, w'') + F(\underline{w'',w'''}, w, w') , \\[8pt]
        \sum_{(3,1)} F(\underline{w,w',w''}, w''') &:= F(\underline{w,w',w''}, w''') + F(\underline{w,w',w'''}, w'') + F(\underline{w,w'',w'''}, w') + F(\underline{w',w'',w'''}, w) , \\
        \sum_{(2,2)} F(\underline{w,w'}, \underline{w'',w'''}) &:= F(\underline{w,w'}, \underline{w'',w'''}) + F(\underline{w,w''}, \underline{w',w'''}) + F(\underline{w,w'''}, \underline{w',w''}) .
    \end{align*}

    By \eqref{eq: variation bounds}, the terms of the $5$ groups in \eqref{eq: fourth chain rule} are bounded by
    \begin{equation*}
    \begin{gathered}
        e^{4Mt} \norm{\nabla^4 g}_\infty , \qquad M t\, e^{4Mt} \norm{\nabla^3 g}_\infty , \qquad M(1+3Mt)\, t\, e^{4Mt} \norm{\nabla^2 g}_\infty , \\
        M^2 t^2 e^{4Mt} \norm{\nabla^2 g}_\infty , \qquad C (1+t)^2\, t\, e^{4Mt} \norm{\nabla g}_\infty ,
    \end{gathered}
    \end{equation*}
    each at most $C (1+t)^3 e^{4Mt}$ times the norm of the derivative of $g$ it carries. Summing the $15$ terms and taking the supremum over $x$ and over the unit vectors gives \eqref{eq: short time} for $j = 4$.
\end{proof}

\subsection{The derivative of order $4$}
\label{sec: fourth}

This subsection extends the smoothing bounds \eqref{eq: smoothing 23} to the derivative of order $4$; the proof uses Lemmas~\ref{lem: variation bounds}, \ref{lem: semigroup}, \ref{lem: BEL}, \ref{lem: variance bound} and \ref{lem: short time}.

\begin{lemma}
\label{lem: fourth from lipschitz}
    Let Assumption~\ref{ass: U} hold. For every $g \in C^4(\mathbb R^d)$ with bounded derivatives of orders $1$ to $4$ and every $0 < t \Le 1$,
    \begin{equation}
        \bigl\| \nabla^4 (e^{t\mathcal L} g)(x) \bigr\| \Le C\, t^{-\frac32}\, \norm{\nabla g}_\infty , \qquad \forall x \in \mathbb R^d .
        \label{eq: fourth from lipschitz}
    \end{equation}
\end{lemma}

\begin{proof}
    Set $\tau = t/2 \Le \frac12$ and $\psi = e^{\tau\mathcal L} g$, so that $e^{2\tau\mathcal L} g = e^{t\mathcal L} g$, and fix $x \in \mathbb R^d$ and unit vectors $v, w, w', w''$. By \eqref{eq: short time}, $\psi$ has bounded derivatives of orders $1$ to $4$. Differentiate \eqref{eq: BEL} in Lemma~\ref{lem: BEL} along $w$, $w'$ and $w''$. By \eqref{eq: variation processes}, a derivative that falls on the stochastic integral raises the order of the variation process in its integrand, from $J_s v$ to $K_s[w] v$, to $\Theta_s[w,w'] v$ and to $\Xi_s[w,w',w''] v$. Distributing the $3$ derivatives between $\psi(X_\tau)$ and the stochastic integral gives
    \begin{align}
        & \sqrt2\, \tau\, \nabla^4 (e^{t\mathcal L} g)(x)[v,w,w',w'']  = \mathbb E^x \biggl[ \partial_w \partial_{w'} \partial_{w''} \psi(X_\tau) \int_0^\tau (J_s v) \cdot \D B_s \biggr] \notag \\
        & \quad + \sum_{(2,1)} \mathbb E^x \biggl[ \partial_w \partial_{w'} \psi(X_\tau) \int_0^\tau \bigl( K_s[w''] v \bigr) \cdot \D B_s \biggr] + \sum_{(2,1)} \mathbb E^x \biggl[ \partial_{w''} \psi(X_\tau) \int_0^\tau \bigl( \Theta_s[w,w'] v \bigr) \cdot \D B_s \biggr] \notag \\
        & \quad + \mathbb E^x \biggl[ \Bigl( \psi(X_\tau) - \mathbb E^x \bigl[ \psi(X_\tau) \bigr] \Bigr) \int_0^\tau \bigl( \Xi_s[w,w',w''] v \bigr) \cdot \D B_s \biggr] ,
        \label{eq: BEL third}
    \end{align}
    where $\sum_{(2,1)}$ is the sum over the partitions of $\{w,w',w''\}$ defined in \eqref{eq: sum21}.
    
    In the term of \eqref{eq: BEL third} carrying $\Xi_s$, $\psi(X_\tau)$ is centered at its mean $\mathbb E^x \bigl[ \psi(X_\tau) \bigr] = e^{t\mathcal L} g(x)$. The centering does not change the value of this term, since the stochastic integral has mean $0$, and it makes \eqref{eq: variance bound} applicable. The exchanges of the derivatives with the expectation and with the stochastic integral are justified as in the proof of \citep[Lemma~8]{ye2026mse}. The variation processes are differentiable in the initial state along each path and bounded by \eqref{eq: variation bounds} uniformly in $x$, so the difference quotients of the integrands converge in $L^2$ over $[0,\tau]$ and the paths, and by the It\^o isometry those of the stochastic integrals converge in $L^2$. Since $\psi$ has bounded derivatives of orders $1$ to $4$, the difference quotients of the factors in front of the stochastic integrals also converge in $L^2$, and the Cauchy--Schwarz inequality passes the limits through the expectation.

    The derivatives of $\psi(X_\tau)$ in \eqref{eq: BEL third} follow from the chain rule and \eqref{eq: variation processes}:
    \begin{equation*}
    \begin{gathered}
        \partial_{w''} \psi(X_\tau) = \nabla \psi(X_\tau) \cdot J_\tau w'' , \qquad
        \partial_w \partial_{w'} \psi(X_\tau) = \nabla^2 \psi(X_\tau)[J_\tau w, J_\tau w'] + \nabla \psi(X_\tau) \cdot K_\tau[w] w' , \\
        \partial_w \partial_{w'} \partial_{w''} \psi(X_\tau) = \nabla^3 \psi(X_\tau)[J_\tau w, J_\tau w', J_\tau w''] + \sum_{(2,1)} \nabla^2 \psi(X_\tau)\bigl[ K_\tau[w] w', J_\tau w'' \bigr] + \nabla \psi(X_\tau) \cdot \Theta_\tau[w,w'] w'' .
    \end{gathered}
    \end{equation*}
    Since $\tau \Le \frac12$, \eqref{eq: short time} and \eqref{eq: smoothing 23} at time $\tau$ give
    \begin{equation*}
        \norm{\nabla\psi}_\infty \Le C\norm{\nabla g}_\infty , \qquad
        \norm{\nabla^2\psi}_\infty \Le C\tau^{-\frac12}\norm{\nabla g}_\infty , \qquad
        \norm{\nabla^3 \psi}_\infty \Le C \tau^{-1}\norm{\nabla g}_\infty ,
    \end{equation*}
    and \eqref{eq: variation bounds} gives $\norm{J_s} \Le C$ and $\norm{K_s[w]}, \norm{\Theta_s[w,w']}, \norm{\Xi_s[w,w',w'']} \Le C s$ for $s \Le \tau$. The bounds on $\nabla^j \psi$ and on the variation processes then give
    \begin{equation}
        \bigl| \partial_{w''} \psi(X_\tau) \bigr| \Le C\norm{\nabla g}_\infty , \quad
        \bigl| \partial_w \partial_{w'} \psi(X_\tau) \bigr| \Le C\tau^{-\frac12}\norm{\nabla g}_\infty , \quad
        \bigl| \partial_w \partial_{w'} \partial_{w''} \psi(X_\tau) \bigr| \Le C\tau^{-1}\norm{\nabla g}_\infty ,
        \label{eq: psi derivative bounds}
    \end{equation}
    and \eqref{eq: variance bound} bounds the centered factor of the term carrying $\Xi_s$:
    \begin{equation}
        \Var\bigl( \psi(X_\tau) \bigr)^{\frac12} = \biggl( \mathbb E^x \Bigl| \psi(X_\tau) - \mathbb E^x \bigl[ \psi(X_\tau) \bigr] \Bigr|^2 \biggr)^{\frac12} \Le C \tau^{\frac12} \norm{\nabla g}_\infty .
        \label{eq: centered factor}
    \end{equation}
    By the It\^o isometry and \eqref{eq: variation bounds},
    \begin{equation}
    \begin{gathered}
        \mathbb E^x \biggl| \int_0^\tau (J_s v) \cdot \D B_s \biggr|^2 \Le C\tau , \qquad
        \mathbb E^x \biggl| \int_0^\tau \bigl( K_s[w''] v \bigr) \cdot \D B_s \biggr|^2 \Le C\tau^3 , \\
        \mathbb E^x \biggl| \int_0^\tau \bigl( \Theta_s[w,w'] v \bigr) \cdot \D B_s \biggr|^2 \Le C\tau^3 , \qquad
        \mathbb E^x \biggl| \int_0^\tau \bigl( \Xi_s[w,w',w''] v \bigr) \cdot \D B_s \biggr|^2 \Le C\tau^3 ,
    \end{gathered}
    \label{eq: ito bounds}
    \end{equation}
    and these bounds hold for every term of the sums $\sum_{(2,1)}$ in \eqref{eq: BEL third}, since $w$, $w'$ and $w''$ are arbitrary unit vectors. The Cauchy--Schwarz inequality, with \eqref{eq: psi derivative bounds}, \eqref{eq: ito bounds} and \eqref{eq: centered factor}, bounds the terms of the $4$ groups in \eqref{eq: BEL third} by
    \begin{equation*}
        C \tau^{-1} \cdot \tau^{\frac12} , \qquad C \tau^{-\frac12} \cdot \tau^{\frac32} , \qquad C \cdot \tau^{\frac32} , \qquad C \tau^{\frac12} \cdot \tau^{\frac32}
    \end{equation*}
    times $\norm{\nabla g}_\infty$. For $\tau \Le \frac12$, each is at most $C \tau^{-\frac12} \norm{\nabla g}_\infty$, so dividing \eqref{eq: BEL third} by $\sqrt2\, \tau$ and taking the supremum over the unit vectors gives
    \begin{equation*}
        \bigl\| \nabla^4 (e^{t\mathcal L} g)(x) \bigr\| \Le C \tau^{-\frac32} \norm{\nabla g}_\infty ,
    \end{equation*}
    which is \eqref{eq: fourth from lipschitz} with $\tau = t/2$.
\end{proof}

\subsection{Conclusion of the proof}
\label{sec: elliptic conclusion}

\begin{proof}[Proof of Theorem~\ref{thm: u decay}]
    The bound on $|u(x,t)|$ in \eqref{eq: u bound} and the bounds on $\nabla^j u(x,t)$ for $1 \Le j \Le 3$ are \citep[Theorem~1]{ye2026mse}, whose assumptions on $U$ and $f$ are the bounds of Assumptions~\ref{ass: U} and \ref{ass: f} up to the orders $5$ and $3$, respectively; decreasing $\lambda$ if necessary, these bounds and \eqref{eq: grad decay} hold with the same rate $\lambda$. For the order $4$, since $\pi(f)$ is a constant, $\nabla^4 u(x,t) = \nabla^4 (e^{t\mathcal L} f)(x)$, and we split at $t = 1$. For $t \Le 1$, the short-time bound \eqref{eq: short time} with $g = f$ and Assumption~\ref{ass: f} give
    \begin{equation*}
        \norm{\nabla^4 u(x,t)} \Le C \sum_{i=1}^{4} \norm{\nabla^i f}_\infty \Le C L \Le C L\, e^{\lambda}\, e^{-\lambda t} .
    \end{equation*}
    For $t \Ge 1$, the semigroup property gives $e^{t\mathcal L} f = e^{\mathcal L} \bigl( e^{(t-1)\mathcal L} f \bigr)$, and $e^{(t-1)\mathcal L} f$ has bounded derivatives of orders $1$ to $4$, since \eqref{eq: short time} at time $t - 1$ bounds them. Hence \eqref{eq: fourth from lipschitz} at time $1$, with $e^{(t-1)\mathcal L} f$ in place of $g$, bounds $\norm{\nabla^4 u(x,t)}$ by $C \norm{\nabla (e^{(t-1)\mathcal L} f)}_\infty$, and \eqref{eq: grad decay} at time $t - 1$ bounds this gradient, so that
    \begin{equation*}
        \norm{\nabla^4 u(x,t)} \Le C\, \bigl\| \nabla (e^{(t-1)\mathcal L} f) \bigr\|_\infty \Le C\, e^{-\lambda(t-1)}\, \norm{\nabla f}_\infty \Le C L\, e^{\lambda}\, e^{-\lambda t} .
    \end{equation*}
    Together, the cases $t \Le 1$ and $t \Ge 1$ give the bound on $\nabla^4 u(x,t)$ in \eqref{eq: u bound}.

    For the discrete Poisson solution \eqref{eq: discrete poisson solution}, the geometric bounds \eqref{eq: u bound} on the summands justify differentiating the series term by term. Summing the derivative bound in \eqref{eq: u bound} over $t = kh$ gives, for $\lambda h \Le 1$ and $1 \Le j \Le 4$,
    \begin{equation*}
        \norm{\nabla^j \phi_h(x)} \Le h \sum_{k=0}^{\infty} \norm{\nabla^j u(x,kh)} \Le C L\, h \sum_{k=0}^{\infty} e^{-\lambda k h} = \frac{C L\, h}{1 - e^{-\lambda h}} \Le \frac{2 C L}{\lambda} ,
    \end{equation*}
    which is \eqref{eq: poisson bound} after $2/\lambda$ is absorbed into $C$.
\end{proof}

\section{Proofs for Section~\ref{sec: modified}}
\label{app: modified}

Throughout this appendix, the constant $C$ depends only on $m$, $M$, $R$ and $d$. The next lemma records the dissipativity of $U$ used in the moment estimates, and the Gaussian tail of $\pi$ that follows from it.

\begin{lemma}
\label{lem: dissipative}
    Let Assumption~\ref{ass: U} hold. For all $x \in \mathbb R^d$,
    \begin{equation}
        \langle x, \nabla U(x) \rangle \Ge \frac{m}{4} |x|^2 - C , \qquad U(x) - U(0) \Ge \frac{m}{8} |x|^2 - C ,
        \label{eq: dissipative bounds}
    \end{equation}
    and consequently $\pi(x) \Le C e^{-\frac{m}{8}|x|^2}$ and $\int_{\mathbb R^d} |x|^q\, \pi(\D x) \Le C_q$ for every $q \Ge 1$.
\end{lemma}

\begin{proof}
    By \citep[Lemma~4]{ye2026mse}, with $R_1 := 4R(1+M/m)$,
    \begin{equation}
        \langle \nabla U(x) - \nabla U(y), x-y \rangle \Ge -M |x-y|^2 , \quad
        \langle \nabla U(x) - \nabla U(y), x-y \rangle \Ge \frac{m}{2} |x-y|^2 \ \text{ for } |x-y| \Ge R_1 .
        \label{eq: dissipative}
    \end{equation}
    Applying \eqref{eq: dissipative} to the pair $(x,0)$ and using $|\nabla U(0)| \Le M$,
    \begin{equation*}
        \langle x, \nabla U(x) \rangle = \langle x, \nabla U(x) - \nabla U(0) \rangle + \langle x, \nabla U(0) \rangle \Ge
        \begin{cases}
            \dfrac{m}{2} |x|^2 - M|x| , & |x| \Ge R_1 , \\[1ex]
            -M|x|^2 - M|x| , & |x| < R_1 ,
        \end{cases}
    \end{equation*}
    and both cases are at least $\frac{m}{4}|x|^2 - C$, which is the drift bound in \eqref{eq: dissipative bounds}. Integrating along the segment from $0$ to $x$, with \eqref{eq: dissipative} at the pair $(rx, 0)$ and the integral split at $r|x| = R_1$ when $|x| \Ge R_1$,
    \begin{equation*}
        U(x) - U(0) = \int_0^1 \langle \nabla U(rx), x \rangle \,\D r \Ge \frac{m}{4} \bigl( |x|^2 - R_1^2 \bigr) - \frac{M}{2} R_1^2 - M |x| ,
    \end{equation*}
    which gives the lower bound on $U(x) - U(0)$ in \eqref{eq: dissipative bounds}. The upper bound $U(y) - U(0) \Le M|y| + \frac{M}{2}|y|^2$ and the lower bound on $U(x) - U(0)$ in \eqref{eq: dissipative bounds} give
    \begin{equation*}
        \pi(x) = \frac{e^{-(U(x) - U(0))}}{\int_{\mathbb R^d} e^{-(U(y) - U(0))} \D y} \Le \frac{e^{C - \frac m8 |x|^2}}{\int_{\mathbb R^d} e^{-M|y| - \frac M2 |y|^2} \D y} \Le C e^{-\frac m8 |x|^2} ,
    \end{equation*}
    and hence $\int_{\mathbb R^d} |x|^q\, \pi(\D x) \Le C_q$ for every $q \Ge 1$.
\end{proof}

\subsection{Proof of Lemma~\ref{lem: modified class}: modified potential estimates}
\label{sec: modified class proof}

\begin{proof}[Proof of Lemma~\ref{lem: modified class}]
    The proof bounds the derivatives of $V$ and of $V\theta_h$, which gives \eqref{eq: modified class}, and then chooses $\delta$ so that $\tilde U_h$ satisfies \eqref{eq: modified class bounds} and the convexity condition. Under Assumption~\ref{ass: U}, $|\nabla U(x)| \Le M(1+|x|)$ and $\norm{\nabla^j U(x)} \Le M$ for $2 \Le j \Le 7$. The derivatives of $|\nabla U|^2$ follow from the Leibniz rule,
    \begin{equation}
        \nabla^i |\nabla U|^2 = \sum_{l=0}^{i} \binom{i}{l}\, \mathrm{Sym}\, \bigl\langle \nabla^{l+1} U ,\, \nabla^{i-l+1} U \bigr\rangle , \qquad 0 \Le i \Le 5 ,
        \label{eq: Leibniz}
    \end{equation}
    where $\langle \cdot , \cdot \rangle$ contracts one index of each factor, so that both sides are tensors of order $i$. The symmetrization is needed: $\langle \nabla^{l+1} U, \nabla^{i-l+1} U \rangle$ is symmetric within its first $l$ indices and within its last $i-l$ indices, but in general not across these groups.
    
    For $i \Ge 1$, each term of \eqref{eq: Leibniz} contains the factor $\nabla U$ at most once and is therefore bounded by $C(1+|x|)$; the derivatives $\nabla^i \Delta U$, contractions of $\nabla^{i+2} U$, are bounded by $dM$ for $0 \Le i \Le 5$. Hence $V = |\nabla U|^2 - \Delta U$ of \eqref{eq: V} satisfies
    \begin{equation}
        |V(x)| \Le C \bigl( 1 + |x| \bigr)^2 , \qquad \norm{\nabla^i V(x)} \Le C \bigl( 1 + |x| \bigr) , \quad 1 \Le i \Le 5 , \qquad \forall x \in \mathbb R^d ,
        \label{eq: V growth}
    \end{equation}
    where $C$ depends only on $M$ and $d$.

    Let $\delta \in (0,1]$, fixed by \eqref{eq: delta bound} at the end of the proof, and let $h \Le \delta$ throughout. By \eqref{eq: modified potential}, $\tilde U_h - U = \frac h4\, V \theta_h$ with $\theta_h(x) = \theta(hx/\delta)$. The function $\theta_h$ is supported in $\{|x| \Le 2\delta/h\}$, where
    \begin{equation}
        \frac h\delta \bigl( 1 + |x| \bigr) \Le \frac h\delta + 2 \Le 3 ,
        \label{eq: cutoff support}
    \end{equation}
    and $\nabla^l \theta_h(x) = (h/\delta)^l\, (\nabla^l \theta)(hx/\delta)$ is bounded by $C (h/\delta)^l$ for $0 \Le l \Le 5$, with $C$ depending on the fixed function $\theta$ only. By the Leibniz rule again,
    \begin{equation}
        \nabla^i (V \theta_h) = \sum_{l=0}^{i} \binom{i}{l}\, \mathrm{Sym} \bigl( \nabla^{i-l} V \otimes \nabla^l \theta_h \bigr) , \qquad 1 \Le i \Le 5 .
        \label{eq: Leibniz cutoff}
    \end{equation}
    On the support of $\theta_h$, \eqref{eq: V growth}, $h \Le \delta$ and \eqref{eq: cutoff support} give
    \begin{equation*}
        \norm{\nabla^{i-l} V(x) \otimes \nabla^l \theta_h(x)} \Le
        \begin{cases}
            C \bigl( 1 + |x| \bigr) , & 0 \Le l \Le i - 1 , \\
            C (h/\delta)^i \bigl( 1 + |x| \bigr)^2 \Le 3C \bigl( 1 + |x| \bigr) , & l = i .
        \end{cases}
    \end{equation*}
    Hence, by \eqref{eq: Leibniz cutoff}, the derivatives of $V\theta_h$ of orders $1$ to $5$ grow at most linearly on the support of $\theta_h$,
    \begin{equation}
        \bigl\| \nabla^i (V\theta_h)(x) \bigr\| \Le C \bigl( 1 + |x| \bigr) , \qquad 1 \Le i \Le 5 , \qquad \forall x\in \mathbb R^d: |x| \Le \frac{2\delta}{h} .
        \label{eq: V cutoff growth}
    \end{equation}
    Since $\nabla^i (V\theta_h)$ vanishes for $|x| > 2\delta/h$, \eqref{eq: V cutoff growth} gives \eqref{eq: modified class}, and, by \eqref{eq: cutoff support},
    \begin{equation}
        \norm{\nabla^i \tilde U_h(x) - \nabla^i U(x)} = \frac h4\, \bigl\| \nabla^i (V\theta_h)(x) \bigr\| \Le C \delta , \qquad 1 \Le i \Le 5 , \qquad \forall x \in \mathbb R^d ,
        \label{eq: modified derivative difference}
    \end{equation}
    with $C$ independent of $\delta$ and $h$. Choose $\delta \in (0,1]$ so small that this constant satisfies
    \begin{equation}
    C\delta \Le \min\Bigl\{ \frac m2 ,\, M \Bigr\} ,
    \label{eq: delta bound}
    \end{equation}
    and let $h_0 \Le \delta$. Then \eqref{eq: modified derivative difference} gives
    \begin{equation}
        \nabla^2 \tilde U_h(x) \succeq \nabla^2 U(x) - C\delta\, I_d , \qquad \norm{\nabla^i \tilde U_h(x)} \Le \norm{\nabla^i U(x)} + C\delta , \quad 1 \Le i \Le 5 , \qquad \forall x \in \mathbb R^d ,
        \label{eq: modified derivative perturbation}
    \end{equation}
    which, with Assumption~\ref{ass: U} and \eqref{eq: delta bound}, gives \eqref{eq: modified class bounds} and the convexity condition of Lemma~\ref{lem: modified class}. Finally, $\tilde U_h \in C^5(\mathbb R^d)$, since $\theta_h \in C^\infty(\mathbb R^d)$ and $U \in C^7(\mathbb R^d)$ gives $V \in C^5(\mathbb R^d)$.
\end{proof}

\subsection{Proof of Lemma~\ref{lem: shifted bias}: modified bias estimate}
\label{sec: shifted bias proof}

\begin{proof}[Proof of Lemma~\ref{lem: shifted bias}]
    The proof bounds the distance between $\tilde\pi_h * \mathcal N(0, \frac h2 I_d)$ and $\pi$, weighted by $1 + |x|$, by $C h^2$. It writes $\tilde\pi_h$ as $\pi\, e^{-(\tilde U_h - U)}$ divided by its normalizing constant, expands this constant to the order $h$, where, since $\pi(V) = 0$, the term of order $h$ comes only from the cutoff and is of order $h^2$, and expands $\tilde\pi_h * \mathcal N(0, \frac h2 I_d)$ to the order $4$ along the Gaussian shift. Since $\tilde\pi_h \propto e^{-\tilde U_h} = e^{-U} e^{-(\tilde U_h - U)}$, where $\tilde U_h - U = \frac h4\, V\theta_h$ by \eqref{eq: modified potential}, we have
    \begin{equation}
        \tilde\pi_h(x) = \frac{\pi(x)\, e^{-(\tilde U_h - U)(x)}}{\pi\bigl( e^{-(\tilde U_h - U)} \bigr)} .
        \label{eq: modified target ratio}
    \end{equation}
    By \eqref{eq: V growth}, $|\tilde U_h - U| \Le C h (1 + |x|^2)$, while $\pi(x) \Le C e^{-\frac{m}{8}|x|^2}$ by Lemma~\ref{lem: dissipative}. Shrinking $h_0$ so that $C h_0 \Le \frac{m}{16}$, with $C$ the constant of the bound on $|\tilde U_h - U|$, we have for $h \Le h_0$,
    \begin{equation}
        e^{|\tilde U_h - U|} \Le C e^{\frac{m}{16}|x|^2} , \qquad \pi\bigl( (1+|x|)^5\, e^{|\tilde U_h - U|} \bigr) \Le C .
        \label{eq: weighted moments}
    \end{equation}
    Taylor's formula for the exponential gives
    \begin{equation}
        e^{-(\tilde U_h - U)} = 1 - \bigl( \tilde U_h - U \bigr) + \varrho_2(x) , \quad |\varrho_2(x)| \Le \frac12 \bigl( \tilde U_h - U \bigr)^2 e^{|\tilde U_h - U|} \Le C h^2 \bigl( 1 + |x| \bigr)^4 e^{|\tilde U_h - U|} ,
        \label{eq: exponential taylor}
    \end{equation}
    hence the remainder $\varrho_2(x)$ satisfies
    \begin{equation}
        \pi\bigl( (1+|x|)\, |\varrho_2(x)| \bigr) \Le C h^2\, \pi\bigl( (1+|x|)^{5} e^{|\tilde U_h - U|} \bigr) \Le C h^2 .
        \label{eq: r bound}
    \end{equation}

    Next, $\pi(V) = - \pi(\mathcal L U) = 0$ by \eqref{eq: integration by parts} with $f = 1$ and \eqref{eq: V}.
    Replacing $V$ by $V\theta_h$ changes $\pi(V)$ by an exponentially small amount: since $1 - \theta_h$ vanishes on $\{|x| \Le \delta/h\}$,
    \begin{equation}
        \pi\bigl( (1+|x|)\, |V|\, (1 - \theta_h) \bigr) \Le C\, \pi\bigl( (1+|x|)^3\, \mathbf 1_{\{|x| \Ge \delta/h\}} \bigr) \Le C e^{-\frac{m \delta^2}{16 h^2}} \Le C h^2 ,
        \label{eq: cutoff tail}
    \end{equation}
    by \eqref{eq: V growth} and $\pi(x) \Le C e^{-\frac{m}{8}|x|^2} \Le C e^{-\frac{m \delta^2}{16 h^2}}\, e^{-\frac{m}{16}|x|^2}$ on $\{|x| \Ge \delta/h\}$.
    Inserting \eqref{eq: exponential taylor} into the denominator of \eqref{eq: modified target ratio}, and using $\pi(V) = 0$, \eqref{eq: r bound} and \eqref{eq: cutoff tail}, we obtain
    \begin{equation}
        \pi\bigl( e^{-(\tilde U_h - U)} \bigr) = 1 - \frac h4\, \pi(V \theta_h) + \pi(\varrho_2) = 1 + \frac h4\, \pi\bigl( V (1 - \theta_h) \bigr) + \pi(\varrho_2) = 1 + \mathcal O(h^2) ,
        \label{eq: ratio expansion}
    \end{equation}
    where $\mathcal O(h^2)$ is bounded by $C h^2$. After shrinking $h_0$, $\pi\bigl( e^{-(\tilde U_h - U)} \bigr) \Ge 1 - C h^2 \Ge \frac12$, hence
    \begin{equation}
        \biggl| \frac{1}{\pi\bigl( e^{-(\tilde U_h - U)} \bigr)} - 1 \biggr| \Le C h^2 .
        \label{eq: ratio bound}
    \end{equation}

    Let $\chi \sim \mathcal N(0,I_d)$, so that
    \begin{equation*}
        \biggl( \tilde\pi_h * \mathcal N\Bigl( 0, \frac h2 I_d \Bigr) \biggr)(x) = \mathbb E\biggl[ \tilde\pi_h\Bigl( x - \sqrt{\frac h2}\, \chi \Bigr) \biggr] .
    \end{equation*}
    Since $\tilde U_h \in C^5$ by Lemma~\ref{lem: modified class}, $\tilde\pi_h \in C^5$. Expanding $\tilde\pi_h$ about $x$ to the order $4$ along $-\sqrt{h/2}\, \chi$,
    \begin{equation*}
        \tilde\pi_h\Bigl( x - \sqrt{\frac h2}\, \chi \Bigr) = \tilde\pi_h(x) - \sqrt{\frac h2}\, \nabla \tilde\pi_h(x) \cdot \chi + \frac h4\, \nabla^2 \tilde\pi_h(x)[\chi,\chi] - \frac16 \Bigl( \frac h2 \Bigr)^{\frac32} \nabla^3\tilde\pi_h(x)\bigl[ \chi^{\otimes3} \bigr] + \varrho_4^{\chi}(x) ,
    \end{equation*}
    with
    \begin{equation}
        |\varrho_4^{\chi}(x)| \Le \frac{h^2}{24}\, |\chi|^4 \int_0^1 (1-s)^3\, \biggl\| \nabla^4 \tilde\pi_h\Bigl( x - s\sqrt{\frac h2}\, \chi \Bigr) \biggr\| \,\D s .
        \label{eq: fourth remainder}
    \end{equation}
    Under the expectation the terms of odd order in $\chi$ vanish and the term of order $2$ has mean $\frac h4\, \Delta \tilde\pi_h(x)$, hence
    \begin{equation}
        \biggl( \tilde\pi_h * \mathcal N\Bigl( 0, \frac h2 I_d \Bigr) \biggr)(x) = \tilde\pi_h(x) + \frac h4\, \Delta \tilde\pi_h(x) + \mathbb E\bigl[ \varrho_4^{\chi}(x) \bigr] .
        \label{eq: shift of density}
    \end{equation}
    Differentiating $\tilde\pi_h \propto e^{-\tilde U_h}$ to the order $2$,
    \begin{equation}
        \Delta \tilde\pi_h = \tilde\pi_h\, \bigl( |\nabla \tilde U_h|^2 - \Delta \tilde U_h \bigr) .
        \label{eq: laplacian of density}
    \end{equation}
    By \eqref{eq: modified target ratio}, \eqref{eq: shift of density} and \eqref{eq: laplacian of density},
    \begin{equation}
        \biggl( \tilde\pi_h * \mathcal N\Bigl( 0, \frac h2 I_d \Bigr) \biggr)(x) = \frac{\pi(x)\, e^{-(\tilde U_h - U)(x)}}{\pi\bigl( e^{-(\tilde U_h - U)} \bigr)} \Bigl( 1 + \frac h4 \bigl( |\nabla \tilde U_h(x)|^2 - \Delta \tilde U_h(x) \bigr) \Bigr) + \mathbb E\bigl[ \varrho_4^{\chi}(x) \bigr] .
        \label{eq: product form}
    \end{equation}
    By \eqref{eq: modified class bounds}, \eqref{eq: modified class} and \eqref{eq: V}, $\bigl| |\nabla \tilde U_h|^2 - \Delta \tilde U_h \bigr| \Le C (1 + |x|)^2$ and
    \begin{equation}
        \bigl| |\nabla \tilde U_h|^2 - \Delta \tilde U_h - V \bigr| \Le |\nabla \tilde U_h - \nabla U|\, |\nabla \tilde U_h + \nabla U| + |\Delta \tilde U_h - \Delta U| \Le C h \bigl( 1 + |x| \bigr)^2 .
        \label{eq: modified V gap}
    \end{equation}
    Taylor's formula for the exponential, with $\tilde U_h - U = \frac h4\, V\theta_h$, the bound $|V| \Le C(1+|x|)^2$ of \eqref{eq: V growth}, the bound on $|\nabla \tilde U_h|^2 - \Delta \tilde U_h$ and the bound on $e^{|\tilde U_h - U|}$ in \eqref{eq: weighted moments} give
    \begin{equation*}
        \biggl| e^{-(\tilde U_h - U)} \Bigl( 1 + \frac h4 \bigl( |\nabla \tilde U_h|^2 - \Delta \tilde U_h \bigr) \Bigr) - 1 - \frac h4 \bigl( |\nabla \tilde U_h|^2 - \Delta \tilde U_h - V\theta_h \bigr) \biggr| \Le C h^2 \bigl( 1 + |x| \bigr)^6 e^{\frac{m}{16}|x|^2} .
    \end{equation*}
    Since $V\theta_h = V - V(1 - \theta_h)$, \eqref{eq: modified V gap} gives
    \begin{equation}
        \biggl| e^{-(\tilde U_h - U)} \Bigl( 1 + \frac h4 \bigl( |\nabla \tilde U_h|^2 - \Delta \tilde U_h \bigr) \Bigr) - 1 \biggr|
        \Le C h^2 \bigl( 1 + |x| \bigr)^6 e^{\frac{m}{16}|x|^2} + \frac h4\, |V| \bigl( 1 - \theta_h \bigr) .
        \label{eq: factor bound}
    \end{equation}
    Multiplying \eqref{eq: factor bound} by $(1+|x|)\pi(x)$ and integrating over $\mathbb R^d$, the term $C h^2 (1+|x|)^6 e^{\frac{m}{16}|x|^2}$ contributes at most $C h^2 \int_{\mathbb R^d} (1 + |x|)^{7} e^{-\frac{m}{16}|x|^2} \,\D x \Le C h^2$ by the density bound $\pi(x) \Le C e^{-\frac{m}{8}|x|^2}$ of Lemma~\ref{lem: dissipative}, and the term $\frac h4\, |V| (1 - \theta_h)$ contributes at most $C h^3$ by \eqref{eq: cutoff tail}, so
    \begin{equation}
        \int_{\mathbb R^d} \bigl( 1 + |x| \bigr) \biggl| e^{-(\tilde U_h - U)} \Bigl( 1 + \frac h4 \bigl( |\nabla \tilde U_h|^2 - \Delta \tilde U_h \bigr) \Bigr) - 1 \biggr|\, \pi(x) \,\D x \Le C h^2 .
        \label{eq: factor integral}
    \end{equation}

    Differentiating $\tilde\pi_h \propto e^{-\tilde U_h}$ to the order $4$ and using \eqref{eq: modified class bounds} with $|\nabla \tilde U_h(y)| \Le 2M(1+|y|)$,
    \begin{equation*}
        \norm{\nabla^4 \tilde\pi_h(y)} \Le C \bigl( 1 + |y| \bigr)^4 \tilde\pi_h(y) , \qquad y \in \mathbb R^d .
    \end{equation*}
    By \eqref{eq: modified target ratio}, \eqref{eq: weighted moments}, the lower bound $\pi( e^{-(\tilde U_h - U)} ) \Ge \frac12$ and the density bound of Lemma~\ref{lem: dissipative}, $\tilde\pi_h\bigl( (1 + |x|)^q \bigr) \Le 2C\, \pi\bigl( (1 + |x|)^q e^{\frac{m}{16}|x|^2} \bigr) \Le C_q$ for $q \Ge 1$. For each fixed $s \in [0,1]$, the substitution $y = x - s\sqrt{h/2}\,\chi$ and the bound on $\tilde\pi_h\bigl( (1 + |x|)^q \bigr)$ with $q = 5$ give
    \begin{equation*}
        \int_{\mathbb R^d} \bigl( 1 + |x| \bigr) \biggl( 1 + \Bigl| x - s\sqrt{\frac h2}\, \chi \Bigr| \biggr)^4 \tilde\pi_h\Bigl( x - s\sqrt{\frac h2}\, \chi \Bigr) \,\D x
        \Le \bigl( 1 + |\chi| \bigr) \int_{\mathbb R^d} \bigl( 1 + |y| \bigr)^5 \tilde\pi_h(y) \,\D y \Le C \bigl( 1 + |\chi| \bigr) .
    \end{equation*}
    Hence, by \eqref{eq: fourth remainder}, Tonelli's theorem and $\mathbb E\bigl[ |\chi|^4 (1+|\chi|) \bigr] \Le C$,
    \begin{equation}
    \begin{aligned}
        & \quad \int_{\mathbb R^d} \bigl( 1 + |x| \bigr) \mathbb E|\varrho_4^{\chi}(x)| \,\D x \\
        & \Le C h^2\, \mathbb E\Biggl[ |\chi|^4 \int_0^1 (1-s)^3 \int_{\mathbb R^d} \bigl( 1 + |x| \bigr) \biggl( 1 + \Bigl| x - s\sqrt{\frac h2}\, \chi \Bigr| \biggr)^4 \tilde\pi_h\Bigl( x - s\sqrt{\frac h2}\, \chi \Bigr) \,\D x \,\D s \Biggr] \Le C h^2 .
    \end{aligned}
    \label{eq: remainder integral}
    \end{equation}

    By \eqref{eq: product form}, the difference between $\tilde\pi_h * \mathcal N(0, \frac h2 I_d)$ and $\pi$ splits into the error $\mathcal T_1$ of the normalizing constant, the error $\mathcal T_2$ of the factor and the remainder $\mathcal T_3$ of the expansion along the Gaussian shift,
    \begin{align}
        & \biggl( \tilde\pi_h * \mathcal N\Bigl( 0, \frac h2 I_d \Bigr) \biggr)(x) - \pi(x) \notag \\
        & = \underbrace{\biggl( \frac{1}{\pi\bigl( e^{-(\tilde U_h - U)} \bigr)} - 1 \biggr) \pi(x)}_{=: \mathcal T_1} + \underbrace{\frac{\pi(x)}{\pi\bigl( e^{-(\tilde U_h - U)} \bigr)} \biggl( e^{-(\tilde U_h - U)(x)} \Bigl( 1 + \frac h4 \bigl( |\nabla \tilde U_h(x)|^2 - \Delta \tilde U_h(x) \bigr) \Bigr) - 1 \biggr)}_{=: \mathcal T_2} + \underbrace{\mathbb E\bigl[ \varrho_4^{\chi}(x) \bigr]}_{=: \mathcal T_3} .
        \label{eq: density split}
    \end{align}
    Integrated against $1 + |x|$, each term of \eqref{eq: density split} is at most $C h^2$: $\mathcal T_1$ by \eqref{eq: ratio bound} and $\pi(1+|x|) \Le C$, $\mathcal T_2$ by the lower bound $\pi\bigl( e^{-(\tilde U_h - U)} \bigr) \Ge \frac12$ and \eqref{eq: factor integral}, and $\mathcal T_3$ by \eqref{eq: remainder integral}, which gives
    \begin{equation}
        \int_{\mathbb R^d} \bigl( 1 + |x| \bigr)\, \Biggl| \biggl( \tilde\pi_h * \mathcal N\Bigl( 0, \frac h2 I_d \Bigr) \biggr)(x) - \pi(x) \Biggr| \,\D x \Le C h^2 .
        \label{eq: density comparison}
    \end{equation}
    Finally, the symmetry of $\mathcal N(0, \frac h2 I_d)$ gives
    \begin{equation*}
        \tilde\pi_h(\tilde f_h) = \int_{\mathbb R^d} f(x)\, \biggl( \tilde\pi_h * \mathcal N\Bigl( 0, \frac h2 I_d \Bigr) \biggr)(x) \,\D x ,
    \end{equation*}
    and since both $\tilde\pi_h * \mathcal N(0, \frac h2 I_d)$ and $\pi$ are probability densities, $f$ may be replaced by $f - f(0)$, with $|f(x) - f(0)| \Le \norm{\nabla f}_\infty |x|$; \eqref{eq: density comparison} then gives \eqref{eq: shifted bias}.
\end{proof}

\section{Proofs for Section~\ref{sec: local}}
\label{app: local}

Statements about a single step are made at the index $0$, with the step $D_0 = Y_1 - Y_0$ of \eqref{eq: chain step} and $\chi_0 \sim \mathcal N(0, I_d)$ independent of $\mathcal F_0$, so that under $\mathbb E_0$ the state $Y_0$ is fixed and
\begin{equation*}
    \mathbb E_0[\chi_0] = 0 , \qquad \mathbb E_0\bigl[ \chi_0^{\otimes 2} \bigr] = I_d , \qquad \mathbb E_0\bigl[ \chi_0^{\otimes 3} \bigr] = 0 , \qquad \mathbb E_0 |\chi_0|^p \Le C , \quad p \Ge 1 .
\end{equation*}
The displacement of the modified dynamics started at $Y_0$ is written $\tilde W_s := \tilde Y_0(s) - Y_0$, $0 \Le s \Le h$. Throughout this appendix, the constant $C$ depends only on $m$, $M$, $R$ and $d$, and also on $p$ where a moment order $p$ appears. The largest moment order used is $p = 12$, for the bounds \eqref{eq: moment bounds} in the proof of Lemma~\ref{lem: moments}; all other proofs in this appendix use $p \Le 6$.

\subsection{Proof of Lemma~\ref{lem: moments}}
\label{sec: moments proof}

The proof uses the bounds \eqref{eq: exact sizes} on one step of the modified dynamics. Lemma~\ref{lem: exact moments} states them first, together with the moments of that step, which are used in Appendix~\ref{sec: matching proof}.

\begin{lemma}
\label{lem: exact moments}
    Let Assumption~\ref{ass: U} hold and let $\tilde U_h$ be the modified potential \eqref{eq: modified potential}. Then, for $0 \Le s \Le h$ and $p \Ge 1$, with $C$ depending only on $M$, $d$ and $p$ and $\pnorm{\cdot}_p$ taken under $\mathbb E_0$,
    \begin{equation}
        \pnorm{\tilde Y_0(s) - Y_0}_p \Le C \sqrt h \bigl( 1 + |Y_0| \bigr) , \qquad
        \pnorm{\tilde Y_0(h) - Y_0 - \sqrt{2h}\, \chi_0}_p \Le C h \bigl( 1 + |Y_0| \bigr) ,
        \label{eq: exact sizes}
    \end{equation}
    and, with $C$ depending only on $M$ and $d$, all derivatives evaluated at $Y_0$ and $1 \Le i,j,k \Le d$,
    \begin{align}
        \Bigl| \mathbb E_0\bigl[ \tilde Y_0(h) - Y_0 \bigr] + h\, \nabla \tilde U_h - \frac{h^2}{2} \bigl( \nabla^2 \tilde U_h\, \nabla \tilde U_h - \nabla \Delta \tilde U_h \bigr) \Bigr| &\Le C h^3 \bigl( 1+|Y_0| \bigr)^3 , \label{eq: exact m1}\\
        \Bigl\| \mathbb E_0\bigl[ (\tilde Y_0(h) - Y_0)^{\otimes 2} \bigr] - 2h\, I_d - h^2 \bigl( \nabla \tilde U_h^{\otimes 2} - 2\, \nabla^2 \tilde U_h \bigr) \Bigr\| &\Le C h^3 \bigl( 1+|Y_0| \bigr)^4 , \label{eq: exact m2}\\
        \Bigl| \bigl( \mathbb E_0\bigl[ (\tilde Y_0(h) - Y_0)^{\otimes 3} \bigr] \bigr)_{ijk} + 2h^2 \bigl( \delta_{ij}\, \partial_k \tilde U_h + \delta_{ik}\, \partial_j \tilde U_h + \delta_{jk}\, \partial_i \tilde U_h \bigr) \Bigr| &\Le C h^3 \bigl( 1+|Y_0| \bigr)^{6} . \label{eq: exact m3}
    \end{align}
    The same statements hold with $U$ in place of $\tilde U_h$ and $X_s$ in place of $\tilde Y_0(s)$, for the exact dynamics \eqref{eq: langevin} started at $X_0 = Y_0$, with the same constants.
\end{lemma}

\begin{proof}
    By Lemma~\ref{lem: modified class}, $|\nabla \tilde U_h(0)| \Le 2M$ and $\norm{\nabla^2 \tilde U_h} \Le 2M$, so that $|\nabla \tilde U_h(x)| \Le 2M(1+|x|)$ for every $x$. The displacement of the modified dynamics \eqref{eq: modified langevin} started at $Y_0$ is
    \begin{equation*}
        \tilde W_s = \tilde Y_0(s) - Y_0 = -\int_0^s \nabla \tilde U_h\bigl( Y_0 + \tilde W_r \bigr) \D r + \sqrt2\, B_s , \qquad 0 \Le s \Le h ,
    \end{equation*}
    where $(B_t)_{t \Ge 0}$ is the Brownian motion of the step, so that $\sqrt2\, B_h = \sqrt{2h}\, \chi_0$. Hence
    \begin{equation*}
        |\tilde W_s| \Le 2M \int_0^s \bigl( 1 + |Y_0| + |\tilde W_r| \bigr) \D r + \sqrt2\, \sup_{r \Le h} |B_r| ,
    \end{equation*}
    and Gr\"onwall's inequality gives
    \begin{equation*}
        \sup_{s \Le h} |\tilde W_s| \Le e^{2Mh} \Bigl( 2Mh \bigl( 1 + |Y_0| \bigr) + \sqrt2\, \sup_{r \Le h} |B_r| \Bigr) .
    \end{equation*}
    Since $\mathbb E \big[\sup_{r \Le h} |B_r|^p \big] \Le C h^{\frac p2}$ by the Burkholder--Davis--Gundy inequality and $h \Le 1$, the bound on $\tilde Y_0(s) - Y_0$ in \eqref{eq: exact sizes} follows. For the bound on $\tilde Y_0(h) - Y_0 - \sqrt{2h}\, \chi_0$,
    \begin{equation*}
        \bigl| \tilde W_h - \sqrt{2h}\, \chi_0 \bigr| = \biggl| \int_0^h \nabla \tilde U_h\bigl( Y_0 + \tilde W_r \bigr) \D r \biggr| \Le 2Mh \Bigl( 1 + |Y_0| + \sup_{r \Le h} |\tilde W_r| \Bigr) ,
    \end{equation*}
    and the Gr\"onwall bound on $\sup_{s \Le h} |\tilde W_s|$ above applies. This proves the bound on $\tilde Y_0(h) - Y_0 - \sqrt{2h}\, \chi_0$ in \eqref{eq: exact sizes}. In particular, $\mathbb E_0 |\tilde Y_0(s)|^q \Le C ( 1 + |Y_0| )^q$ for $s \Le h$ and every $q \Ge 1$.

    Next we consider the moments of $\tilde W_h = \tilde Y_0(h) - Y_0$ shown in \eqref{eq: exact m1}--\eqref{eq: exact m3}. Since $\tilde Y_0(h)$ is driven by the modified dynamics \eqref{eq: modified langevin} with the generator $\tilde{\mathcal L}_h \psi = \Delta \psi - \nabla \tilde U_h \cdot \nabla \psi$ of \eqref{eq: modified generator}, the estimate of these moments relies on a Taylor expansion with respect to $\tilde{\mathcal L}_h$.
    
    Let $\ell(z) := z - Y_0$, a linear function of $z \in \mathbb R^d$ with components $\ell_i$, so that $\tilde W_h = \ell(\tilde Y_0(h))$ and the entries of $\mathbb E_0\bigl[ \tilde W_h^{\otimes n} \bigr]$, $n = 1, 2, 3$, are $\mathbb E_0\bigl[ \psi(\tilde Y_0(h)) \bigr]$ for the polynomials $\psi = \ell_i$, $\ell_i \ell_j$, $\ell_i \ell_j \ell_k$. Here and in the rest of the proof, $z$ denotes the variable of $\ell$ and of these polynomials.
    
    For each such $\psi$, since $\tilde U_h \in C^5(\mathbb R^d)$ by Lemma~\ref{lem: modified class}, the functions $\tilde{\mathcal L}_h \psi$ and $\tilde{\mathcal L}_h^2 \psi$ are of class $C^2$ and $\tilde{\mathcal L}_h^3 \psi$ is continuous. By \eqref{eq: modified class bounds}, $\tilde{\mathcal L}_h \psi$, $\tilde{\mathcal L}_h^2 \psi$, $\tilde{\mathcal L}_h^3 \psi$ and the gradients of $\psi$, $\tilde{\mathcal L}_h \psi$, $\tilde{\mathcal L}_h^2 \psi$ have polynomial growth, so that, by $\mathbb E_0 |\tilde Y_0(s)|^q \Le C ( 1 + |Y_0| )^q$, the stochastic integral in It\^o's formula is a martingale. Dynkin's formula, applied to $\psi$, then to $\tilde{\mathcal L}_h \psi$ and $\tilde{\mathcal L}_h^2 \psi$ inside the integrals, gives with Fubini's theorem the Taylor expansion to the order $2$ with integral remainder, where $\psi(Y_0) = 0$ is omitted since $\ell(Y_0) = 0$,
    \begin{equation}
        \mathbb E_0\bigl[ \psi(\tilde Y_0(h)) \bigr] = h\, (\tilde{\mathcal L}_h\psi)(Y_0) + \frac{h^2}{2}\, (\tilde{\mathcal L}_h^2\psi)(Y_0)
        + \frac12 \int_0^h (h-s)^2\, \mathbb E_0\bigl[ (\tilde{\mathcal L}_h^3 \psi)(\tilde Y_0(s)) \bigr] \D s .
        \label{eq: ito taylor}
    \end{equation}
    Here $\tilde{\mathcal L}_h \psi(Y_0)$ and $\tilde{\mathcal L}_h^2 \psi(Y_0)$ are computed with the product rule
    \begin{equation}
        \tilde{\mathcal L}_h(g_1 g_2) = g_1 \tilde{\mathcal L}_h g_2 + g_2 \tilde{\mathcal L}_h g_1 + 2\, \nabla g_1 \cdot \nabla g_2 , \qquad g_1, g_2 \in C^2(\mathbb R^d) .
        \label{eq: generator product}
    \end{equation}

    Next we evaluate the explicit $\tilde{\mathcal L}_h$ and $\tilde{\mathcal L}_h^2$ operations on $\psi = \ell_i$, $\ell_i \ell_j$, $\ell_i \ell_j \ell_k$, to give the explicit $h$ and $h^2$ coefficients in the moment estimates \eqref{eq: exact m1}--\eqref{eq: exact m3}.
    \begin{enumerate}
    \item For $\psi = \ell_i$, the linearity of $\ell_i$ gives $\tilde{\mathcal L}_h \ell_i = -\partial_i \tilde U_h$, and
    \begin{equation*}
        \tilde{\mathcal L}_h^2 \ell_i = -\tilde{\mathcal L}_h \partial_i \tilde U_h = -\partial_i \Delta \tilde U_h + \sum_{j} \partial_j \tilde U_h\, \partial_i \partial_j \tilde U_h = \bigl( \nabla^2 \tilde U_h\, \nabla \tilde U_h - \nabla \Delta \tilde U_h \bigr)_i .
    \end{equation*}
    At $z = Y_0$, $h\, \tilde{\mathcal L}_h \ell_i + \frac{h^2}{2}\, \tilde{\mathcal L}_h^2 \ell_i$ is the component $i$ of $-h\, \nabla \tilde U_h + \frac{h^2}{2} \bigl( \nabla^2 \tilde U_h\, \nabla \tilde U_h - \nabla \Delta \tilde U_h \bigr)$, the approximation of $\mathbb E_0[\tilde W_h]$ in \eqref{eq: exact m1}.

    \item For $\psi = \ell_i \ell_j$, \eqref{eq: modified generator} gives
    \begin{equation*}
        \tilde{\mathcal L}_h(\ell_i \ell_j) = 2\delta_{ij} - \partial_i \tilde U_h\, \ell_j - \partial_j \tilde U_h\, \ell_i ,
    \end{equation*}
    which equals $2\delta_{ij}$ at $z = Y_0$. Since $\tilde{\mathcal L}_h$ annihilates the constant $2\delta_{ij}$, $\tilde{\mathcal L}_h^2(\ell_i \ell_j) = -\tilde{\mathcal L}_h(\partial_i \tilde U_h\, \ell_j) - \tilde{\mathcal L}_h(\partial_j \tilde U_h\, \ell_i)$, and \eqref{eq: generator product} with $g_1 = \partial_i \tilde U_h$ and $g_2 = \ell_j$ gives, at $z = Y_0$, where $\ell_j = 0$,
    \begin{equation*}
        \tilde{\mathcal L}_h\bigl( \partial_i \tilde U_h\, \ell_j \bigr)(Y_0) = \partial_i \tilde U_h\, \tilde{\mathcal L}_h \ell_j + 2\, \nabla \partial_i \tilde U_h \cdot \nabla \ell_j = -\partial_i \tilde U_h\, \partial_j \tilde U_h + 2\, \partial_i \partial_j \tilde U_h .
    \end{equation*}
    This identity and its copy with $i$ and $j$ exchanged give
    \begin{equation*}
        \tilde{\mathcal L}_h^2(\ell_i \ell_j)(Y_0) = 2\, \partial_i \tilde U_h\, \partial_j \tilde U_h - 4\, \partial_i \partial_j \tilde U_h ,
    \end{equation*}
    so that, at $z = Y_0$, $h\, \tilde{\mathcal L}_h(\ell_i \ell_j) + \frac{h^2}{2}\, \tilde{\mathcal L}_h^2(\ell_i \ell_j)$ is the entry $(i,j)$ of $2h\, I_d + h^2 \bigl( \nabla \tilde U_h^{\otimes 2} - 2\, \nabla^2 \tilde U_h \bigr)$, the approximation of $\mathbb E_0[\tilde W_h^{\otimes 2}]$ in \eqref{eq: exact m2}.

    \item For $\psi = \ell_i \ell_j \ell_k$,
    \begin{equation*}
        \tilde{\mathcal L}_h\bigl( \ell_i \ell_j \ell_k \bigr) = 2 \bigl( \delta_{ij}\, \ell_k + \delta_{ik}\, \ell_j + \delta_{jk}\, \ell_i \bigr) - \bigl( \partial_i \tilde U_h\, \ell_j \ell_k + \partial_j \tilde U_h\, \ell_i \ell_k + \partial_k \tilde U_h\, \ell_i \ell_j \bigr) ,
    \end{equation*}
    which vanishes at $z = Y_0$. In $\tilde{\mathcal L}_h^2(\ell_i \ell_j \ell_k)(Y_0)$, the terms $2\delta_{ij}\, \ell_k$, $2\delta_{ik}\, \ell_j$, $2\delta_{jk}\, \ell_i$ contribute $-2\delta_{ij}\, \partial_k \tilde U_h$, $-2\delta_{ik}\, \partial_j \tilde U_h$, $-2\delta_{jk}\, \partial_i \tilde U_h$, since $\tilde{\mathcal L}_h \ell_k = -\partial_k \tilde U_h$. By \eqref{eq: generator product} with $g_1 = \partial_i \tilde U_h$ and $g_2 = \ell_j \ell_k$, whose value and gradient vanish at $Y_0$, the term $-\partial_i \tilde U_h\, \ell_j \ell_k$ contributes $-\partial_i \tilde U_h\, \tilde{\mathcal L}_h(\ell_j \ell_k)(Y_0) = -2\delta_{jk}\, \partial_i \tilde U_h$, and in the same way the terms $-\partial_j \tilde U_h\, \ell_i \ell_k$ and $-\partial_k \tilde U_h\, \ell_i \ell_j$ contribute $-2\delta_{ik}\, \partial_j \tilde U_h$ and $-2\delta_{ij}\, \partial_k \tilde U_h$. Hence
    \begin{equation*}
        \tilde{\mathcal L}_h^2 \bigl( \ell_i \ell_j \ell_k \bigr)(Y_0) = -4 \bigl( \delta_{ij}\, \partial_k \tilde U_h + \delta_{ik}\, \partial_j \tilde U_h + \delta_{jk}\, \partial_i \tilde U_h \bigr) ,
    \end{equation*}
    so that $\frac{h^2}{2}\, \tilde{\mathcal L}_h^2(\ell_i \ell_j \ell_k)(Y_0)$ is the entry $(i,j,k)$ of the approximation of $\mathbb E_0[\tilde W_h^{\otimes 3}]$ in \eqref{eq: exact m3}.
    \end{enumerate}

    It remains to bound the integral term of \eqref{eq: ito taylor}. Each of $\psi$, $\tilde{\mathcal L}_h \psi$, $\tilde{\mathcal L}_h^2 \psi$ and $\tilde{\mathcal L}_h^3 \psi$ is a linear combination of products of components of $\ell$, components of $\nabla \tilde U_h$ and derivatives of $\tilde U_h$ of orders $2$ to $5$; call the degree of such a product the number of its factors that are components of $\ell$ or of $\nabla \tilde U_h$. By the linear growth of $\nabla \tilde U_h$ and \eqref{eq: modified class bounds}, a product of degree $q$ is bounded by $C (1 + |z| + |Y_0|)^q$. In $\tilde{\mathcal L}_h$, the Laplacian does not raise the degree, and $-\nabla \tilde U_h \cdot \nabla$ raises it by at most $1$, and only where it differentiates a derivative of $\tilde U_h$ of order at least $2$. No such derivative occurs in $\psi$ or $\tilde{\mathcal L}_h \psi$, so for $\psi$ a product of $n$ components of $\ell$ the products in $\tilde{\mathcal L}_h \psi$ and $\tilde{\mathcal L}_h^2 \psi$ have degree at most $n$, those in $\tilde{\mathcal L}_h^3 \psi$ have degree at most $n + 1$, and
    \begin{equation*}
        \bigl| (\tilde{\mathcal L}_h^3 \psi)(z) \bigr| \Le C \bigl( 1 + |z| + |Y_0| \bigr)^{n+1} .
    \end{equation*}
    With $\mathbb E_0 (1 + |\tilde Y_0(s)| + |Y_0|)^{n+1} \Le C(1+|Y_0|)^{n+1}$ for $s \Le h$, the integral term is bounded by $C h^3 (1+|Y_0|)^{n+1}$, which is at most the remainders stated in \eqref{eq: exact m1}--\eqref{eq: exact m3}.

    For the exact Langevin dynamics \eqref{eq: langevin}, the same argument applies with $U$ and $X_s$ in place of $\tilde U_h$ and $\tilde Y_0(s)$, since by Assumption~\ref{ass: U} the derivatives of $U$ satisfy \eqref{eq: modified class bounds} with $M$ in place of $2M$.
\end{proof}

\begin{proof}[Proof of Lemma~\ref{lem: moments}]
    We prove the bounds with the exponents $2q$ and $q$ in place of $12$ and $6$ for an integer $q \Ge 1$, the case $q = 6$ being \eqref{eq: moment bounds}; in this proof $C$ may depend on $q$ as well. The proof is a Lyapunov argument: a bound on the displacement over one step and a drift condition give a contraction of $\mathbb E_k |Y_{k+1}|^{2q}$ against $|Y_k|^{2q}$, which is iterated over $k$. By \eqref{eq: lm markov}, the step of the chain at the index $k$ is
    \begin{equation*}
        D_k := Y_{k+1} - Y_k = \sqrt{2h}\, \chi_k - h\, \nabla U(Z_k) , \qquad Z_k = Y_k + \sqrt{\frac h2}\, \chi_k ,
    \end{equation*}
    and Assumption~\ref{ass: U} gives $|\nabla U(Z_k)| \Le M ( 1 + |Y_k| + \sqrt{h/2}\, |\chi_k| )$.

    We first bound the displacement over one step. Since $\mathbb E_k |\sqrt{2h}\, \chi_k|^{p} \Le C h^{\frac p2}$, the bound on $|\nabla U(Z_k)|$ gives, for every $p \Ge 1$,
    \begin{equation}
        \mathbb E_k |D_k|^{p} \Le C \bigl( \mathbb E_k |\sqrt{2h}\, \chi_k|^{p} + h^{p}\, \mathbb E_k |\nabla U(Z_k)|^{p} \bigr)
        \Le C h^{\frac p2} + C h^{p} \bigl( 1 + |Y_k| \bigr)^{p} ,
        \label{eq: step displacement}
    \end{equation}
    and, since \eqref{eq: exact sizes} at the index $k$ gives $\pnorm{\tilde Y_k(h) - Y_k - \sqrt{2h}\,\chi_k}_p \Le C h (1 + |Y_k|)$ under $\mathbb E_k$,
    \begin{equation}
        \mathbb E_k |\tilde Y_k(h) - Y_k|^{p} \Le C h^{\frac p2} + C h^{p} \bigl( 1 + |Y_k| \bigr)^{p} .
        \label{eq: exact displacement}
    \end{equation}

    Next we derive the drift condition. Splitting $\langle Y_k, \nabla U(Z_k)\rangle = \langle Z_k, \nabla U(Z_k)\rangle - \sqrt{h/2}\, \langle \chi_k, \nabla U(Z_k)\rangle$ and applying the drift bound in \eqref{eq: dissipative bounds} at $Z_k$ and the bound on $|\nabla U(Z_k)|$,
    \begin{equation*}
        \langle Y_k, \nabla U(Z_k) \rangle \Ge \frac m4 |Z_k|^2 - C - M \sqrt{\frac h2}\, |\chi_k| \Bigl( 1 + |Y_k| + \sqrt{\frac h2}\, |\chi_k| \Bigr) .
    \end{equation*}
    Taking $\mathbb E_k$, with $\mathbb E_k |Z_k|^2 = |Y_k|^2 + \frac{dh}{2}$ and $\mathbb E_k |\chi_k|^2 = d$,
    \begin{equation}
        \langle Y_k, \mathbb E_k [\nabla U(Z_k)] \rangle \Ge \frac m4 |Y_k|^2 - C - C \sqrt h \bigl( 1 + |Y_k| \bigr) \Ge \frac{m}{8} |Y_k|^2 - C ,
        \label{eq: drift condition}
    \end{equation}
    since $C \sqrt h\, |Y_k| \Le \frac m8 |Y_k|^2 + C$ by Young's inequality and $h \Le 1$.

    Next we combine \eqref{eq: step displacement} and \eqref{eq: drift condition} into the Lyapunov estimate. Let $S_k := |Y_{k+1}|^2 - |Y_k|^2 = 2\langle Y_k, D_k\rangle + |D_k|^2$. Since $\chi_k$ is centered and independent of $\mathcal F_k$, \eqref{eq: drift condition} and \eqref{eq: step displacement} with $p = 2$ give
    \begin{equation*}
        \mathbb E_k S_k = -2h \langle Y_k, \mathbb E_k [\nabla U(Z_k)] \rangle + \mathbb E_k |D_k|^2 \Le -\frac m4 h |Y_k|^2 + C h + C h^2 \bigl( 1 + |Y_k| \bigr)^2 \Le -\frac m8 h |Y_k|^2 + C h
    \end{equation*}
    for $h \Le h_0$ with $C h_0 \Le \frac m8$, while $|S_k| \Le 2|Y_k||D_k| + |D_k|^2$ and \eqref{eq: step displacement} give, for $1 \Le j \Le q$,
    \begin{equation*}
        \mathbb E_k |S_k|^{j} \Le C \bigl( |Y_k|^j\, \mathbb E_k |D_k|^j + \mathbb E_k |D_k|^{2j} \bigr) \Le C h^{\frac j2} |Y_k|^{j} + C h^{j} \bigl( 1 + |Y_k| \bigr)^{2j} .
    \end{equation*}
    Expanding
    \begin{equation*}
        |Y_{k+1}|^{2q} = \bigl( |Y_k|^2 + S_k \bigr)^q = \sum_{j=0}^{q} \binom{q}{j} |Y_k|^{2(q-j)} S_k^{j}
    \end{equation*}
    and inserting the bounds on $\mathbb E_k S_k$ and $\mathbb E_k |S_k|^{j}$, with $h \Le 1$,
    \begin{equation*}
        q\, |Y_k|^{2q-2}\, \mathbb E_k S_k \Le -\frac{qm}{8} h |Y_k|^{2q} + C h |Y_k|^{2q-2} , \qquad
        |Y_k|^{2(q-j)}\, \mathbb E_k |S_k|^{j} \Le C h |Y_k|^{2q - j} + C h^2 (1+|Y_k|)^{2q} , \quad j \Ge 2 .
    \end{equation*}
    By Young's inequality, the terms of degree below $2q$ in $|Y_k|$ satisfy
    \begin{equation*}
        C h |Y_k|^{2q-2} + C h \sum_{j=2}^{q} |Y_k|^{2q-j} \Le \frac{qm}{32}\, h |Y_k|^{2q} + C h ,
    \end{equation*}
    and, for $h \Le h_0$ with $C h_0$ small enough, the terms of order $h^2$ satisfy
    \begin{equation*}
        C h^2 \bigl( 1+|Y_k| \bigr)^{2q} \Le \frac{qm}{32}\, h |Y_k|^{2q} + C h .
    \end{equation*}
    With these bounds on the terms of degree below $2q$ and of order $h^2$, the expansion of $\mathbb E_k |Y_{k+1}|^{2q}$ gives
    \begin{equation*}
        \mathbb E_k |Y_{k+1}|^{2q} \Le |Y_k|^{2q} - \frac{qm}{8}\, h |Y_k|^{2q} + 2 \cdot \frac{qm}{32}\, h |Y_k|^{2q} + C h ,
    \end{equation*}
    that is,
    \begin{equation}
        \mathbb E_k |Y_{k+1}|^{2q} \Le \Bigl( 1 - \frac{qm}{16} h \Bigr) |Y_k|^{2q} + C h ,
        \label{eq: lyapunov}
    \end{equation}
    with $h_0 \Le 16/(qm)$ making the factor lie in $[0,1)$. Taking expectations in \eqref{eq: lyapunov} and iterating over $k$ gives the bound on $\mathbb E|Y_k|^{2q}$ in \eqref{eq: moment bounds}, and averaging \eqref{eq: step displacement} and \eqref{eq: exact displacement} with $p = 2q$ over the law of $Y_k$ gives the displacement bounds in \eqref{eq: moment bounds}:
    \begin{equation*}
        \mathbb E |Y_k|^{2q} \Le |Y_0|^{2q} + \frac{16 C}{qm} , \qquad
        \max\bigl\{ \mathbb E |D_k|^{2q} ,\; \mathbb E |\tilde Y_k(h) - Y_k|^{2q} \bigr\} \Le C h^{q} + C h^{2q}\, \mathbb E \bigl( 1 + |Y_k| \bigr)^{2q} \Le C h^{q} \bigl( 1 + |Y_0| \bigr)^{2q} .
    \end{equation*}

    The argument carries over to the stochastic gradient chain \eqref{eq: sglm markov}, with $\hat{\mathbb E}_k$ of \eqref{eq: sg filtration} in place of $\mathbb E_k$, once the displacement bounds and the drift condition are verified for it. The step of this chain is
    \begin{equation*}
        D_k = Y_{k+1} - Y_k = \sqrt{2h}\, \chi_k - h\, \nabla \hat U(Z_k; \xi_k) .
    \end{equation*}
    By the moment bound of Assumption~\ref{ass: sg} and Jensen's inequality, for $1 \Le p \Le 12$,
    \begin{equation*}
        \hat{\mathbb E}_k |\nabla \hat U(Z_k; \xi_k)|^p \Le M^p\, \hat{\mathbb E}_k \bigl( 1 + |Z_k| \bigr)^p \Le C \bigl( 1 + |Y_k| \bigr)^p ,
    \end{equation*}
    which is the bound used for $\nabla U(Z_k)$ in \eqref{eq: step displacement}, with the same constant, so that
    \begin{equation*}
        \hat{\mathbb E}_k |D_k|^{p} \Le C h^{\frac p2} + C h^{p} \bigl( 1 + |Y_k| \bigr)^{p} , \qquad 1 \Le p \Le 12 .
    \end{equation*}
    The bound \eqref{eq: exact displacement} holds with $\hat{\mathbb E}_k$ in place of $\mathbb E_k$, since $\tilde Y_k(h) - Y_k$ is driven by $(B_{kh+t} - B_{kh})_{0 \Le t \Le h}$, which is independent of $\hat{\mathcal F}_k$. Since the estimator is unbiased and $\xi_k$ is independent of $\hat{\mathcal F}_{k}$ and $\chi_k$, the derivation of \eqref{eq: drift condition} with $\hat{\mathbb E}_k$ in place of $\mathbb E_k$ gives
    \begin{equation*}
        \bigl\langle Y_k, \hat{\mathbb E}_k [\nabla \hat U(Z_k; \xi_k)] \bigr\rangle = \bigl\langle Y_k, \hat{\mathbb E}_k [\nabla U(Z_k)] \bigr\rangle \Ge \frac{m}{8} |Y_k|^2 - C .
    \end{equation*}
    The Lyapunov estimate uses only $\hat{\mathbb E}_k[\chi_k] = 0$, the drift condition and the bound on $\hat{\mathbb E}_k |D_k|^{p}$ with $p \Le 2q$, so for $q \Le 6$ it gives \eqref{eq: lyapunov} with $\hat{\mathbb E}_k$ in place of $\mathbb E_k$ and the same constants. Averaging the bound on $\hat{\mathbb E}_k |D_k|^{p}$ and \eqref{eq: exact displacement} under $\hat{\mathbb E}_k$, with $p = 2q \Le 12$, over the law of $Y_k$ then gives \eqref{eq: moment bounds} for \eqref{eq: sglm markov}.
\end{proof}

\subsection{Proof of Lemma~\ref{lem: sizes}}
\label{sec: sizes proof}

\begin{proof}[Proof of Lemma~\ref{lem: sizes}]
    We prove the bounds in $\pnorm{\cdot}_p$ for every $p \Ge 1$, with $C$ depending also on $p$; the case $p = 4$ is \eqref{eq: sizes}. The bounds on $Y_1 - Y_0$ and on $\tilde Y_0(s) - Y_0$ in \eqref{eq: sizes} are direct. For the step $D_0 = Y_1 - Y_0 = \sqrt{2h}\, \chi_0 - h\, \nabla U(Z_0)$ of \eqref{eq: chain step}, the moments of $\chi_0$ and Assumption~\ref{ass: U} give
    \begin{equation*}
        \pnorm{\sqrt{2h}\, \chi_0}_p \Le C\sqrt h , \qquad \pnorm{h\, \nabla U(Z_0)}_p \Le M h \Bigl( 1 + |Y_0| + \sqrt{\frac h2}\, \pnorm{\chi_0}_p \Bigr) \Le C h (1+|Y_0|) ,
    \end{equation*}
    and the bound on $\tilde Y_0(s) - Y_0$ is that in \eqref{eq: exact sizes}.

    Next we bound the local error $Y_1 - \tilde Y_0(h) = D_0 - \tilde W_h$. Subtracting the displacement $\tilde W_h = \sqrt{2h}\, \chi_0 - \int_0^h \nabla \tilde U_h(Y_0 + \tilde W_r)\,\D r$ of \eqref{eq: modified langevin} from $D_0$ gives
    \begin{equation}
        Y_1 - \tilde Y_0(h) = \underbrace{\int_0^h \Bigl( \nabla U\bigl( Y_0 + \tilde W_r \bigr) - \nabla U( Z_0 ) \Bigr) \D r}_{=: \mathcal T_1} + \underbrace{\int_0^h \Bigl( \nabla \tilde U_h\bigl( Y_0 + \tilde W_r \bigr) - \nabla U\bigl( Y_0 + \tilde W_r \bigr) \Bigr) \D r}_{=: \mathcal T_2} .
        \label{eq: local error explicit}
    \end{equation}
    The integrands of $\mathcal T_1$ and $\mathcal T_2$ satisfy
    \begin{equation*}
        \bigl| \nabla U( Y_0 + \tilde W_r ) - \nabla U( Z_0 ) \bigr| \Le M \Bigl| \tilde W_r - \sqrt{\frac h2}\, \chi_0 \Bigr| , \qquad
        \bigl| ( \nabla \tilde U_h - \nabla U )( Y_0 + \tilde W_r ) \bigr| \Le C h \bigl( 1 + |Y_0| + |\tilde W_r| \bigr) ,
    \end{equation*}
    by $\norm{\nabla^2 U} \Le M$ of Assumption~\ref{ass: U} with $Z_0 = Y_0 + \sqrt{h/2}\, \chi_0$, and by \eqref{eq: modified class}, respectively. Minkowski's inequality for the integrals over $r \in [0,h]$ and the bound $\pnorm{\tilde W_r}_p \Le C \sqrt h\, (1 + |Y_0|)$ of \eqref{eq: exact sizes} then give
    \begin{equation*}
    \begin{aligned}
        \pnorm{\mathcal T_1}_p & \Le M \int_0^h \Bigl( \pnorm{\tilde W_r}_p + \sqrt{\frac h2}\, \pnorm{\chi_0}_p \Bigr) \D r \Le C h^{\frac32} (1+|Y_0|) , \\
        \pnorm{\mathcal T_2}_p & \Le C h \int_0^h \bigl( 1 + |Y_0| + \pnorm{\tilde W_r}_p \bigr) \D r \Le C h^2 (1+|Y_0|) .
    \end{aligned}
    \end{equation*}
    Adding the bounds on $\pnorm{\mathcal T_1}_p$ and $\pnorm{\mathcal T_2}_p$ gives $\pnorm{Y_1 - \tilde Y_0(h)}_p \Le C h^{\frac32} (1+|Y_0|)$, as stated in \eqref{eq: sizes} for $p = 4$.
\end{proof}

\subsection{Proof of Lemma~\ref{lem: moment matching}}
\label{sec: matching proof}

The proof compares the moments of $\tilde W_h$ in Lemma~\ref{lem: exact moments} with the following expansion of the moments of the step $Y_1 - Y_0$ of the chain.

\begin{lemma}
\label{lem: step moments}
    Let Assumption~\ref{ass: U} hold and let $h \in (0,1]$. Then, with all derivatives of $U$ evaluated at $Y_0$ and $C$ depending only on $M$ and $d$, the moments of the step \eqref{eq: chain step} of the chain satisfy, for $1 \Le i,j,k \Le d$,
    \begin{align}
        \Bigl| \mathbb E_0[Y_1 - Y_0] + h\, \nabla U + \frac{h^2}{4}\, \nabla \Delta U \Bigr| &\Le C h^3 , \label{eq: step m1}\\
        \Bigl\| \mathbb E_0\bigl[ (Y_1 - Y_0)^{\otimes 2} \bigr] - 2h\, I_d - h^2 \bigl( \nabla U^{\otimes 2} - 2\, \nabla^2 U \bigr) \Bigr\| &\Le C h^3 \bigl( 1+|Y_0| \bigr) , \label{eq: step m2}\\
        \Bigl| \bigl( \mathbb E_0\bigl[ (Y_1 - Y_0)^{\otimes 3} \bigr] \bigr)_{ijk} + 2h^2 \bigl( \delta_{ij}\, \partial_k U + \delta_{ik}\, \partial_j U + \delta_{jk}\, \partial_i U \bigr) \Bigr| &\Le C h^3 \bigl( 1+|Y_0| \bigr)^{3} . \label{eq: step m3}
    \end{align}
\end{lemma}

\begin{proof}
    The step \eqref{eq: chain step} reads $Y_1 - Y_0 = \sqrt{2h}\, \chi_0 - h\, \nabla U(Z_0)$ with $Z_0 - Y_0 = \sqrt{h/2}\, \chi_0$, so its moments follow from expansions of $\nabla U(Z_0)$ about $Y_0$ in powers of $\chi_0$. In this proof, $\mathcal O(h^q)$ denotes a tensor of norm at most $C h^q$. Taylor's formula to the order $3$, with the remainder controlled by $\norm{\nabla^5 U} \Le M$, gives
    \begin{equation}
        \nabla U(Z_0) = \nabla U + \sqrt{\frac h2}\, \nabla^2 U\, \chi_0 + \frac h4\, \nabla^3 U\bigl[ \chi_0^{\otimes 2} \bigr] + \frac16 \Bigl( \frac h2 \Bigr)^{\frac32} \nabla^4 U\bigl[ \chi_0^{\otimes 3} \bigr] + \varrho_4 , \qquad |\varrho_4| \Le \frac{M}{24} \Bigl( \frac h2 \Bigr)^2 |\chi_0|^4 ,
        \label{eq: gradient taylor}
    \end{equation}
    and Taylor's formula to the order $1$, with the remainder controlled by $\norm{\nabla^3 U} \Le M$, gives
    \begin{equation}
        \nabla U(Z_0) = \nabla U + \sqrt{\frac h2}\, \nabla^2 U\, \chi_0 + \varrho_2 , \qquad |\varrho_2| \Le \frac M2\, \frac h2\, |\chi_0|^2 .
        \label{eq: gradient taylor crude}
    \end{equation}
    The expansion \eqref{eq: gradient taylor} is used where $\nabla U(Z_0)$ carries a factor $h$ or $h^{\frac32}$, namely in the mean and in the cross term of the moment of order $2$, and \eqref{eq: gradient taylor crude} where it carries a factor $h^2$ or $h^{\frac52}$. In addition, $|\nabla U(x)| \Le M(1+|x|)$ gives $|\nabla U(Z_0)| \Le M ( 1 + |Y_0| + \sqrt{h/2}\, |\chi_0| )$, and under $\mathbb E_0$ every term of odd degree in $\chi_0$ in \eqref{eq: gradient taylor} and \eqref{eq: gradient taylor crude}, and in their products with powers of $\chi_0$, has mean zero.

    Next we compute the moments of orders $1$, $2$, $3$ of $Y_1 - Y_0$ from \eqref{eq: gradient taylor} and \eqref{eq: gradient taylor crude}, to give the explicit $h$ and $h^2$ coefficients in \eqref{eq: step m1}--\eqref{eq: step m3}.
    \begin{enumerate}
    \item For the mean, the terms of degrees $1$ and $3$ in $\chi_0$ of \eqref{eq: gradient taylor} have mean zero, and $\nabla^3 U[\mathbb E_0 \chi_0^{\otimes 2}] = \nabla^3 U[I_d] = \nabla \Delta U$ by the contraction convention of Section~\ref{subsec: notation}, so that
    \begin{equation*}
        \mathbb E_0\bigl[ \nabla U(Z_0) \bigr] = \nabla U + \frac h4\, \nabla \Delta U + \mathbb E_0[\varrho_4] , \qquad |\mathbb E_0[\varrho_4]| \Le C h^2 .
    \end{equation*}
    Multiplied by $-h$, this is \eqref{eq: step m1}, since $\mathbb E_0[Y_1 - Y_0] = -h\, \mathbb E_0[\nabla U(Z_0)]$.

    \item For the moment of order $2$,
    \begin{equation*}
        (Y_1 - Y_0)^{\otimes 2} = 2h\, \chi_0^{\otimes 2} - 2h\sqrt{2h}\, \mathrm{Sym}\bigl( \chi_0 \otimes \nabla U(Z_0) \bigr) + h^2\, \nabla U(Z_0)^{\otimes 2} .
    \end{equation*}
    Here $\mathrm{Sym}$ is essential to ensure that both sides are symmetric matrices. The term $2h\, \chi_0^{\otimes 2}$ has conditional mean $2h I_d$. For the cross term $-2h\sqrt{2h}\, \mathrm{Sym}( \chi_0 \otimes \nabla U(Z_0) )$, insert \eqref{eq: gradient taylor} into $\chi_0 \otimes \nabla U(Z_0)$. The terms $\chi_0 \otimes \nabla U$ and $\chi_0 \otimes \nabla^3 U[\chi_0^{\otimes 2}]$ have odd degree and mean zero, $\mathbb E_0[\chi_0 \otimes \nabla^2 U\, \chi_0] = \nabla^2 U$ since $\mathbb E_0[\chi_0^{\otimes 2}] = I_d$ and $\nabla^2 U$ is symmetric, and the terms $\frac16 (\frac h2)^{\frac32}\, \chi_0 \otimes \nabla^4 U[\chi_0^{\otimes 3}]$ and $\chi_0 \otimes \varrho_4$ have conditional means of norm at most $C h^{\frac32}$. Hence
    \begin{equation*}
        \mathbb E_0\bigl[ \chi_0 \otimes \nabla U(Z_0) \bigr] = \sqrt{\frac h2}\, \nabla^2 U + \mathcal O\bigl( h^{\frac32} \bigr) , \qquad
        -2h\sqrt{2h}\, \mathrm{Sym}\, \mathbb E_0\bigl[ \chi_0 \otimes \nabla U(Z_0) \bigr] = -2h^2\, \nabla^2 U + \mathcal O( h^3 ) .
    \end{equation*}
    For the term $h^2\, \nabla U(Z_0)^{\otimes 2}$, \eqref{eq: gradient taylor crude} gives
    \begin{equation}
    \begin{aligned}
        \nabla U(Z_0)^{\otimes 2} & = \nabla U^{\otimes 2} + 2\, \mathrm{Sym}\Bigl( \nabla U \otimes \sqrt{\frac h2}\, \nabla^2 U\, \chi_0 \Bigr) + 2\, \mathrm{Sym}\bigl( \nabla U \otimes \varrho_2 \bigr) \\
        & \quad + \frac h2 \bigl( \nabla^2 U\, \chi_0 \bigr)^{\otimes 2} + 2\, \mathrm{Sym}\Bigl( \sqrt{\frac h2}\, \nabla^2 U\, \chi_0 \otimes \varrho_2 \Bigr) + \varrho_2^{\otimes 2} .
    \end{aligned}
    \label{eq: gradient square}
    \end{equation}
    The term $2\, \mathrm{Sym}( \nabla U \otimes \sqrt{h/2}\, \nabla^2 U\, \chi_0 )$ has mean zero. By $|\nabla U| \Le M(1+|Y_0|)$ and the bound on $\varrho_2$ in \eqref{eq: gradient taylor crude}, the conditional means of the remaining $4$ terms of \eqref{eq: gradient square} satisfy
    \begin{equation*}
    \begin{aligned}
        & \Bigl\| \mathbb E_0\bigl[ 2\, \mathrm{Sym}(\nabla U \otimes \varrho_2) \bigr] \Bigr\| \Le C h (1+|Y_0|) , \qquad
        \biggl\| \mathbb E_0\Bigl[ \frac h2 (\nabla^2 U\, \chi_0)^{\otimes 2} \Bigr] \biggr\| \Le C h , \\
        & \Bigl\| \mathbb E_0\bigl[ 2\, \mathrm{Sym}\bigl(\sqrt{h/2}\, \nabla^2 U\, \chi_0 \otimes \varrho_2\bigr) \bigr] \Bigr\| \Le C h^{\frac32} , \qquad
        \Bigl\| \mathbb E_0\bigl[ \varrho_2^{\otimes 2} \bigr] \Bigr\| \Le C h^2 ,
    \end{aligned}
    \end{equation*}
    so that
    \begin{equation*}
        \Bigl\| h^2\, \mathbb E_0\bigl[ \nabla U(Z_0)^{\otimes 2} \bigr] - h^2\, \nabla U^{\otimes 2} \Bigr\| \Le C h^3 (1+|Y_0|) .
    \end{equation*}
    Adding the conditional means of the $3$ terms of $(Y_1 - Y_0)^{\otimes 2}$, namely $2h I_d$, $-2h^2\, \nabla^2 U + \mathcal O(h^3)$ and $h^2\, \nabla U^{\otimes 2}$ up to $C h^3 (1+|Y_0|)$, gives \eqref{eq: step m2}.

    \item For the moment of order $3$,
    \begin{equation*}
        (Y_1 - Y_0)^{\otimes 3} = (2h)^{\frac32}\, \chi_0^{\otimes 3} - 6h^2\, \mathrm{Sym}\bigl( \chi_0^{\otimes 2} \otimes \nabla U(Z_0) \bigr) + 3h^2 \sqrt{2h}\, \mathrm{Sym}\bigl( \chi_0 \otimes \nabla U(Z_0)^{\otimes 2} \bigr) - h^3\, \nabla U(Z_0)^{\otimes 3} .
    \end{equation*}
    The term $(2h)^{\frac32}\, \chi_0^{\otimes 3}$ has mean zero. For the term $-6h^2\, \mathrm{Sym}( \chi_0^{\otimes 2} \otimes \nabla U(Z_0) )$, insert \eqref{eq: gradient taylor crude} into the entry $\mathbb E_0[\chi_{0,i}\, \chi_{0,j}\, \nabla U(Z_0)_k]$: the term $\nabla U$ gives $\delta_{ij}\, \partial_k U$, the term linear in $\chi_0$ has mean zero, and the remainder $\varrho_2$ contributes at most $C h$, so that
    \begin{equation*}
        \mathbb E_0\bigl[ \chi_{0,i}\, \chi_{0,j}\, \nabla U(Z_0)_k \bigr] = \delta_{ij}\, \partial_k U + \mathcal O(h) ,
    \end{equation*}
    and, since $\delta_{ij}\, \partial_k U$ is symmetric in $i$ and $j$, its symmetrization is $\frac13 ( \delta_{ij}\, \partial_k U + \delta_{ik}\, \partial_j U + \delta_{jk}\, \partial_i U )$, so that
    \begin{equation*}
        -6h^2\, \mathrm{Sym}\, \mathbb E_0\bigl[ \chi_0^{\otimes 2} \otimes \nabla U(Z_0) \bigr]_{ijk} = -2h^2 \bigl( \delta_{ij}\, \partial_k U + \delta_{ik}\, \partial_j U + \delta_{jk}\, \partial_i U \bigr) + \mathcal O(h^3) .
    \end{equation*}
    For the term $3h^2 \sqrt{2h}\, \mathrm{Sym}( \chi_0 \otimes \nabla U(Z_0)^{\otimes 2} )$, insert \eqref{eq: gradient square} into $\chi_0 \otimes \nabla U(Z_0)^{\otimes 2}$: the term $\chi_0 \otimes \nabla U^{\otimes 2}$ has mean zero, and every other term carries a factor $\sqrt{h/2}\, \nabla^2 U\, \chi_0$ or $\varrho_2$, so its conditional mean has norm at most $C \sqrt h (1+|Y_0|)$. For the term $h^3\, \nabla U(Z_0)^{\otimes 3}$, the bound on $|\nabla U(Z_0)|$ suffices. Hence
    \begin{equation*}
    \begin{aligned}
        & \Bigl\| 3h^2 \sqrt{2h}\, \mathrm{Sym}\, \mathbb E_0\bigl[ \chi_0 \otimes \nabla U(Z_0)^{\otimes 2} \bigr] \Bigr\| \Le C h^3 (1+|Y_0|) , \\
        & \Bigl\| h^3\, \mathbb E_0\bigl[ \nabla U(Z_0)^{\otimes 3} \bigr] \Bigr\| \Le h^3\, \mathbb E_0 |\nabla U(Z_0)|^3 \Le C h^3 (1+|Y_0|)^3 .
    \end{aligned}
    \end{equation*}
    Adding the conditional means of the $4$ terms of $(Y_1 - Y_0)^{\otimes 3}$ gives \eqref{eq: step m3}.
    \end{enumerate}
\end{proof}

\begin{proof}[Proof of Lemma~\ref{lem: moment matching}]
    The moments of $\tilde W_h = \tilde Y_0(h) - Y_0$ are expanded in \eqref{eq: exact m1}--\eqref{eq: exact m3} with coefficients in $\tilde U_h$, and those of $D_0 = Y_1 - Y_0$ in \eqref{eq: step m1}--\eqref{eq: step m3} with coefficients in $U$; Lemma~\ref{lem: step moments} applies since $h \Le h_0 \Le 1$, and all derivatives are evaluated at $Y_0$. Next we compare the two expansions. Their coefficients differ through $\tilde U_h - U$, and by \eqref{eq: modified class} and the linear growth of $\nabla U$, and of $\nabla \tilde U_h$ from Lemma~\ref{lem: modified class},
    \begin{equation}
        |\nabla \tilde U_h - \nabla U| \Le C h \bigl( 1+|Y_0| \bigr) , \quad \norm{\nabla \tilde U_h^{\otimes 2} - \nabla U^{\otimes 2}} \Le C h \bigl( 1+|Y_0| \bigr)^2 , \quad \max_{2 \Le i \Le 3} \norm{\nabla^i \tilde U_h - \nabla^i U} \Le C h \bigl( 1+|Y_0| \bigr) .
        \label{eq: coefficient replacement}
    \end{equation}
    
    Next we prove \eqref{eq: moment matching} for $j = 1, 2, 3$ with \eqref{eq: coefficient replacement}, first for $j = 2, 3$ and then for $j = 1$.
    \begin{enumerate}
    \item For $j = 2$ and $j = 3$, replacing $\tilde U_h$ by $U$ in the terms of order $h^2$ of \eqref{eq: exact m2} and \eqref{eq: exact m3} changes them in norm by at most $C h^3 (1+|Y_0|)^2$, by \eqref{eq: coefficient replacement}. After this replacement, the approximations of $\mathbb E_0[\tilde W_h^{\otimes j}]$ in \eqref{eq: exact m2} and \eqref{eq: exact m3} are those of $\mathbb E_0[D_0^{\otimes j}]$ in \eqref{eq: step m2} and \eqref{eq: step m3}, which proves \eqref{eq: moment matching} for $j = 2$ and $j = 3$.
    \item For $j = 1$, the term $-h\, \nabla \tilde U_h$ of the approximation of $\mathbb E_0[\tilde W_h]$ in \eqref{eq: exact m1} is kept, since \eqref{eq: coefficient replacement} bounds $h(\nabla \tilde U_h - \nabla U)$ only by $C h^2 (1+|Y_0|)$. In the term of order $h^2$, the replacement of $\tilde U_h$ by $U$ introduces an additional error of at most $C h^3 (1+|Y_0|)^2$ in norm, by \eqref{eq: coefficient replacement}, the linear growth of $\nabla \tilde U_h$ and $\norm{\nabla^2 U} \Le M$, the bound on $\nabla^3 \tilde U_h - \nabla^3 U$ in \eqref{eq: coefficient replacement} controlling $\nabla \Delta \tilde U_h - \nabla \Delta U$. Subtracting \eqref{eq: step m1} then gives
    \begin{equation}
        \Bigl| \mathbb E_0[\tilde W_h] - \mathbb E_0[D_0] + h \bigl( \nabla \tilde U_h - \nabla U \bigr) - \frac{h^2}{2} \bigl( \nabla^2 U\, \nabla U - \nabla \Delta U \bigr) - \frac{h^2}{4}\, \nabla \Delta U \Bigr| \Le C h^3 \bigl( 1+|Y_0| \bigr)^3 .
        \label{eq: m1 difference}
    \end{equation}
    It remains to bound the terms of \eqref{eq: m1 difference} other than $\mathbb E_0[\tilde W_h] - \mathbb E_0[D_0]$ by $C h^3 (1+|Y_0|)^2$, and here the cutoff $\theta_h$ of \eqref{eq: modified potential} enters. On the ball $\{|Y_0| \Le \delta/h\}$, where $\theta_h = 1$ and $\nabla \theta_h = 0$, $\nabla \tilde U_h - \nabla U = \frac h4 \nabla V = \frac h4 ( 2\, \nabla^2 U\, \nabla U - \nabla \Delta U )$ by \eqref{eq: V}, so that
    \begin{equation}
    \begin{aligned}
        & \quad h \bigl( \nabla \tilde U_h - \nabla U \bigr) - \frac{h^2}{2} \bigl( \nabla^2 U\, \nabla U - \nabla \Delta U \bigr) - \frac{h^2}{4}\, \nabla \Delta U \\
        & = \frac{h^2}{4} \bigl( 2\, \nabla^2 U\, \nabla U - \nabla \Delta U \bigr) - \frac{h^2}{2}\, \nabla^2 U\, \nabla U + \frac{h^2}{4}\, \nabla \Delta U = 0 .
    \end{aligned}
    \label{eq: m1 cancellation}
    \end{equation}
    Outside that ball, \eqref{eq: coefficient replacement} and Assumption~\ref{ass: U} give
    \begin{equation*}
        h\, \bigl| \nabla \tilde U_h(Y_0) - \nabla U(Y_0) \bigr| \Le C h^2 (1+|Y_0|) , \qquad
        \Bigl| \frac{h^2}{2} \bigl( \nabla^2 U\, \nabla U - \nabla \Delta U \bigr) + \frac{h^2}{4}\, \nabla \Delta U \Bigr| \Le C h^2 (1+|Y_0|) ,
    \end{equation*}
    and $1 \Le h|Y_0|/\delta$ there, so the left side of \eqref{eq: m1 cancellation} is at most
    \begin{equation*}
        C h^2 (1+|Y_0|) \cdot \frac{h |Y_0|}{\delta} \Le C h^3 (1+|Y_0|)^2
    \end{equation*}
    in norm. On the ball and outside it, \eqref{eq: m1 difference} therefore gives
    \begin{equation*}
        \bigl| \mathbb E_0[\tilde W_h] - \mathbb E_0[D_0] \bigr| \Le C h^3 (1+|Y_0|)^3 \Le C h^3 (1+|Y_0|)^6 ,
    \end{equation*}
    which proves \eqref{eq: moment matching} for $j = 1$.
    \end{enumerate}
\end{proof}

\subsection{Proof of Lemma~\ref{lem: weak}}
\label{sec: weak proof}

\begin{proof}[Proof of Lemma~\ref{lem: weak}]
    The proof expands $G(Y_1) - G(\tilde Y_0(h))$ about $Y_0$ to the order $3$, so that the terms of orders $1$ to $3$ are differences of the moments of $D_0$ and $\tilde W_h$, bounded by Lemma~\ref{lem: moment matching}, and the remainder is bounded by Lemma~\ref{lem: sizes}.
    
    Work at the index $0$, with $D_0$ and $\tilde W_h$ as in the opening of Appendix~\ref{app: local}. Write $Q_\varsigma := \varsigma Y_1 + (1 - \varsigma) \tilde Y_0(h)$ for $\varsigma \in [0,1]$, the interpolated position of \citep{ye2026mse}, so that $\frac{\D}{\D\varsigma} Q_\varsigma = Y_1 - \tilde Y_0(h)$, the local error, and
    \begin{equation}
        G(Y_1) - G(\tilde Y_0(h)) = \int_0^1 \bigl\langle \nabla G(Q_\varsigma),\, Y_1 - \tilde Y_0(h) \bigr\rangle \,\D\varsigma .
        \label{eq: weak start}
    \end{equation}
    Expand $\nabla G$ about $Y_0$ to the order $2$ with the integral remainder,
    \begin{equation*}
        \nabla G(Y_0 + v) = \nabla G(Y_0) + \nabla^2 G(Y_0)\, v + \frac12\, \nabla^3 G(Y_0)\bigl[ v^{\otimes2} \bigr] + \varrho_3(v) ,
        \qquad |\varrho_3(v)| \Le \frac16\, \norm{\nabla^4 G}_\infty\, |v|^3 ,
    \end{equation*}
    and insert it into \eqref{eq: weak start} with $v = Q_\varsigma - Y_0$. Since $\nabla^2 G(Y_0)$ and $\nabla^3 G(Y_0)$ are symmetric, $\nabla^j G(Y_0) : T = \nabla^j G(Y_0) : \mathrm{Sym}\, T$ for every tensor $T$ of order $j = 2, 3$. Moreover,
    \begin{equation*}
        \frac{\D}{\D\varsigma} \bigl( Q_\varsigma - Y_0 \bigr)^{\otimes j} = j\, \mathrm{Sym}\bigl( (Q_\varsigma - Y_0)^{\otimes (j-1)} \otimes ( Y_1 - \tilde Y_0(h) ) \bigr) , \qquad j = 2, 3 ,
    \end{equation*}
    and integrating over $\varsigma \in [0,1]$, with $Q_0 - Y_0 = \tilde W_h$ and $Q_1 - Y_0 = D_0$,
    \begin{equation*}
    \begin{aligned}
        \mathrm{Sym} \biggl( \int_0^1 ( Q_\varsigma - Y_0 ) \,\D\varsigma \otimes \bigl( Y_1 - \tilde Y_0(h) \bigr) \biggr) & = \frac12 \Bigl( D_0^{\otimes2} - \tilde W_h^{\otimes2} \Bigr) , \\
        \mathrm{Sym} \biggl( \int_0^1 ( Q_\varsigma - Y_0 )^{\otimes2} \,\D\varsigma \otimes \bigl( Y_1 - \tilde Y_0(h) \bigr) \biggr) & = \frac13 \Bigl( D_0^{\otimes3} - \tilde W_h^{\otimes3} \Bigr) .
    \end{aligned}
    \end{equation*}
    These identities turn \eqref{eq: weak start} into
    \begin{equation}
        G(Y_1) - G(\tilde Y_0(h)) = \sum_{j=1}^{3} \frac{1}{j!}\, \nabla^j G(Y_0) : \Bigl( D_0^{\otimes j} - \tilde W_h^{\otimes j} \Bigr) + \varrho_4 ,
        \label{eq: G expansion full}
    \end{equation}
    where the remainder $\varrho_4 := \int_0^1 \bigl\langle \varrho_3(Q_\varsigma - Y_0),\, Y_1 - \tilde Y_0(h) \bigr\rangle \,\D\varsigma$ satisfies
    \begin{equation*}
        |\varrho_4| \Le \frac16\, \norm{\nabla^4 G}_\infty\, \bigl| Y_1 - \tilde Y_0(h) \bigr|\, \max\bigl\{ |D_0| ,\, |\tilde W_h| \bigr\}^3 .
    \end{equation*}
    In \eqref{eq: G expansion full}, $G(Y_1)$ and $G(\tilde Y_0(h))$ are both expanded about the initial value $Y_0$, which is why Lemma~\ref{lem: moment matching} estimates the moments of $D_0$ and $\tilde W_h$ up to the order $3$. Taking $\mathbb E_0$, \eqref{eq: moment matching} bounds the terms of orders $1$ to $3$ in \eqref{eq: G expansion full}:
    \begin{equation*}
        \biggl| \sum_{j=1}^{3} \frac{1}{j!}\, \nabla^j G(Y_0) : \mathbb E_0\bigl[ D_0^{\otimes j} - \tilde W_h^{\otimes j} \bigr] \biggr|
        \Le C h^3 \sum_{j=1}^{3} \norm{\nabla^j G}_\infty \bigl( 1 + |Y_0| \bigr)^{6} .
    \end{equation*}
    For the remainder $\varrho_4$, \eqref{eq: sizes} gives
    \begin{equation*}
        \pnorm{Y_1 - \tilde Y_0(h)}_4 \Le C h^{\frac32}(1+|Y_0|) , \qquad \pnorm{D_0}_4 + \pnorm{\tilde W_h}_4 \Le C \sqrt h (1+|Y_0|) ,
    \end{equation*}
    and H\"older's inequality gives
    \begin{equation*}
        \mathbb E_0 |\varrho_4|
        \Le \frac16 \norm{\nabla^4 G}_\infty\, \pnorm{Y_1 - \tilde Y_0(h)}_4 \bigl( \pnorm{D_0}_{4} + \pnorm{\tilde W_h}_{4} \bigr)^3
        \Le C \norm{\nabla^4 G}_\infty\, h^3 \bigl( 1+|Y_0| \bigr)^{4} .
    \end{equation*}
    Adding the bounds on the terms of orders $1$ to $3$ and on $\varrho_4$ proves \eqref{eq: weak} at the index $0$, and by time homogeneity at every $k$.
\end{proof}

\subsection{Proof of Lemma~\ref{lem: sg weak}}
\label{sec: sg weak proof}

The proof uses two bounds on the covariance of the gradient noise, which also serve Appendices~\ref{app: bias} to \ref{app: w1}. We write the gradient noise and its covariance \eqref{eq: covariance} as
\begin{equation*}
    \eta(x;\xi) := \nabla \hat U(x;\xi) - \nabla U(x) , \qquad \mathcal C(x) = \mathbb E_\xi\bigl[ \eta(x;\xi) \otimes \eta(x;\xi) \bigr] .
\end{equation*}
Let
\begin{equation}
    S(x;\xi) := \sup_{|x'-x| \Le 1} \norm{\nabla^2 \hat U(x';\xi)} , \qquad \mathbb E_\xi \big[S(x;\xi)^2 \big] \Le M^2 ,
    \label{eq: hessian sup}
\end{equation}
the bound being the Hessian condition of Assumption~\ref{ass: sg}. Let $|x - x'| \Le 1$. Then
\begin{equation*}
    \eta(x;\xi) - \eta(x';\xi) = \int_0^1 \bigl( \nabla^2 \hat U(x_s;\xi) - \nabla^2 U(x_s) \bigr) (x - x') \D s , \qquad x_s := x' + s(x - x') ,
\end{equation*}
so the definition of $S$ and $\norm{\nabla^2 U} \Le M$ give
\begin{equation}
    |\eta(x;\xi) - \eta(x';\xi)| \Le \bigl( S(x;\xi) + M \bigr) |x - x'| .
    \label{eq: noise lipschitz}
\end{equation}
The identity
\begin{equation*}
    \eta(x;\xi)^{\otimes 2} - \eta(x';\xi)^{\otimes 2} = \eta(x;\xi) \otimes \bigl( \eta(x;\xi) - \eta(x';\xi) \bigr) + \bigl( \eta(x;\xi) - \eta(x';\xi) \bigr) \otimes \eta(x';\xi)
\end{equation*}
gives
\begin{equation*}
    \bigl\| \eta(x;\xi)^{\otimes 2} - \eta(x';\xi)^{\otimes 2} \bigr\| \Le \bigl( |\eta(x;\xi)| + |\eta(x';\xi)| \bigr) |\eta(x;\xi) - \eta(x';\xi)| .
\end{equation*}
Minkowski's inequality, with Assumption~\ref{ass: sg} for $\eta(x;\xi)$ and $\eta(x';\xi)$ and with \eqref{eq: hessian sup} for $S$, gives
\begin{equation}
    \mathbb E_\xi \Big[\bigl( |\eta(x;\xi)| + |\eta(x';\xi)| \bigr)^2 \Big]  \Le 4\Sigma , \qquad \mathbb E_\xi \big[( S(x;\xi) + M )^2\big] \Le 4M^2 ,
    \label{eq: noise moments}
\end{equation}
so \eqref{eq: covariance}, the bound on $\eta(x;\xi)^{\otimes 2} - \eta(x';\xi)^{\otimes 2}$, \eqref{eq: noise lipschitz} and the Cauchy--Schwarz inequality with \eqref{eq: noise moments} give
\begin{equation*}
\begin{aligned}
    \norm{\mathcal C(x) - \mathcal C(x')} & = \Bigl\| \mathbb E_\xi\bigl[ \eta(x;\xi)^{\otimes 2} - \eta(x';\xi)^{\otimes 2} \bigr] \Bigr\| \Le \mathbb E_\xi\Bigl[ \bigl( |\eta(x;\xi)| + |\eta(x';\xi)| \bigr) |\eta(x;\xi) - \eta(x';\xi)| \Bigr] \\
    & \Le \mathbb E_\xi\Bigl[ \bigl( |\eta(x;\xi)| + |\eta(x';\xi)| \bigr) \bigl( S(x;\xi) + M \bigr) \Bigr]\, |x - x'| \Le 2\Sigma^{\frac12} \cdot 2M\, |x - x'| .
\end{aligned}
\end{equation*}
For $|x - x'| > 1$, since $\mathcal C(x)$ and $\mathcal C(x')$ are positive semidefinite, \eqref{eq: covariance} gives
\begin{equation*}
    \norm{\mathcal C(x) - \mathcal C(x')} \Le \mathrm{Tr}\bigl[ \mathcal C(x) \bigr] + \mathrm{Tr}\bigl[ \mathcal C(x') \bigr] \Le 2\Sigma \Le 2\Sigma\, |x - x'| .
\end{equation*}
Since $\Sigma \Le M^2$,
\begin{equation*}
    4 \Sigma^{\frac12} M \Le 4 M^2 , \qquad 2\Sigma \Le 4 M^2 ,
\end{equation*}
so in both cases
\begin{equation}
    \norm{\mathcal C(x) - \mathcal C(x')} \Le 4 M^2 |x - x'| , \qquad x, x' \in \mathbb R^d .
    \label{eq: covariance lipschitz}
\end{equation}
Writing $\mathcal C(x) = \sum_{i=1}^d \lambda_i\, u_i \otimes u_i$ with $\lambda_i \Ge 0$ and $(u_i)$ orthonormal, \eqref{eq: covariance} gives
\begin{equation}
    |A : \mathcal C(x)| = \Bigl| \sum_{i=1}^d \lambda_i \langle u_i, A u_i \rangle \Bigr| \Le \norm{A}\, \mathrm{Tr}\bigl[ \mathcal C(x) \bigr] \Le \Sigma \norm{A} , \qquad A \in \mathbb R^{d \times d} , \ x \in \mathbb R^d .
    \label{eq: trace bound}
\end{equation}

\begin{proof}[Proof of Lemma~\ref{lem: sg weak}]
    Work at the index $0$, where $\hat{\mathbb E}_0$ and $\mathbb E_0$ agree on functions of $Y_0$ and $(B_s)_{s \Le h}$, both fixing $Y_0$ and averaging over $(B_s)_{s \Le h}$. With $D_0 = \sqrt{2h}\, \chi_0 - h\, \nabla U(Z_0)$, the right side of \eqref{eq: chain step}, the step of \eqref{eq: sglm markov} adds the gradient noise $\eta_0$ at $Z_0$, multiplied by $-h$,
    \begin{equation*}
        Y_1 = Y_0 + D_0 - h \eta_0 , \qquad \eta_0 := \eta(Z_0;\xi_0) = \nabla \hat U(Z_0;\xi_0) - \nabla U(Z_0) .
    \end{equation*}
    The proof expands $G(Y_1)$ about $Y_0 + D_0$ to the order $2$ in $h\eta_0$, compares $G(Y_0 + D_0)$ with $G(\tilde Y_0(h))$ by Lemma~\ref{lem: weak}, and finally replaces $\nabla^2 G(Y_0 + D_0)$ by $\nabla^2 G(Y_0)$ in the term of order $h^2$.
    
    Let $\hat{\mathcal H}_0 := \sigma\bigl( \hat{\mathcal F}_0, (B_s)_{s \Le h} \bigr)$. Then $Y_0$, $D_0$, $Z_0$ and $\tilde Y_0(h)$ are $\hat{\mathcal H}_0$-measurable, and since $\xi_0$ is independent of $\hat{\mathcal H}_0$, Assumption~\ref{ass: sg} and \eqref{eq: covariance} give
    \begin{equation*}
        \mathbb E\bigl[ \eta_0 \,\big|\, \hat{\mathcal H}_0 \bigr] = 0 , \qquad \mathbb E\bigl[ \eta_0 \otimes \eta_0 \,\big|\, \hat{\mathcal H}_0 \bigr] = \mathcal C(Z_0) .
    \end{equation*}
    Taylor's formula for $G$ at $Y_0 + D_0$, to the order $2$ in $h\eta_0$, gives
    \begin{equation*}
        G(Y_1) = G(Y_0 + D_0) - h\, \nabla G(Y_0+D_0) \cdot \eta_0 + \frac{h^2}{2}\, \nabla^2 G(Y_0+D_0)[\eta_0, \eta_0] + \varrho_3^{\eta_0} , \qquad |\varrho_3^{\eta_0}| \Le \frac{h^3}{6}\, \norm{\nabla^3 G}_\infty\, |\eta_0|^3 .
    \end{equation*}
    Under $\mathbb E[\,\cdot \,|\, \hat{\mathcal H}_0]$, the term linear in $\eta_0$ vanishes and the term quadratic in $\eta_0$ becomes $\frac{h^2}{2}\, \nabla^2 G(Y_0+D_0) : \mathcal C(Z_0)$. For the remainder, $|\eta_0| \Le |\nabla \hat U(Z_0;\xi_0)| + |\nabla U(Z_0)|$, and the moment bound of Assumption~\ref{ass: sg}, applied given $\hat{\mathcal H}_0$, with $|\nabla U(x)| \Le M(1+|x|)$ and $h \Le 1$, gives
    \begin{equation*}
        \hat{\mathbb E}_0 |\eta_0|^3 \Le C\, \hat{\mathbb E}_0 \bigl( 1 + |Z_0| \bigr)^3 \Le C \bigl( 1 + |Y_0| \bigr)^3 .
    \end{equation*}
    Hence
    \begin{equation}
        \Bigl| \hat{\mathbb E}_0\bigl[ G(Y_1) \bigr] - \mathbb E_0\bigl[ G(Y_0 + D_0) \bigr] - \frac{h^2}{2}\, \hat{\mathbb E}_0\bigl[ \nabla^2 G(Y_0 + D_0) : \mathcal C(Z_0) \bigr] \Bigr| \Le C h^3 \norm{\nabla^3 G}_\infty \bigl( 1+|Y_0| \bigr)^3 .
        \label{eq: sg noise expansion}
    \end{equation}
    Next, $Y_0 + D_0$ is the state after one step of \eqref{eq: lm markov} from $Y_0$, so $\mathbb E_0[G(Y_0+D_0)] - \mathbb E_0[G(\tilde Y_0(h))]$ is the local weak error of the chain with the full gradient, and \eqref{eq: weak} at $k = 0$ gives
    \begin{equation}
        \bigl| \mathbb E_0\bigl[ G(Y_0 + D_0) - G(\tilde Y_0(h)) \bigr] \bigr| \Le C h^3 \sum_{j=1}^{4} \norm{\nabla^j G}_\infty \bigl( 1+|Y_0| \bigr)^6 .
        \label{eq: sg exact part}
    \end{equation}

    It remains to replace $\nabla^2 G(Y_0 + D_0)$ by $\nabla^2 G(Y_0)$ in the term carrying $\mathcal C(Z_0)$, at the cost $C h^3$. By Taylor's formula,
    \begin{equation*}
        \nabla^2 G(Y_0 + D_0) = \nabla^2 G(Y_0) + \nabla^3 G(Y_0)[D_0] + \varrho_2^{D_0} , \qquad \norm{\varrho_2^{D_0}} \Le \frac12 \norm{\nabla^4 G}_\infty |D_0|^2 ,
    \end{equation*}
    and, with $D_0 = \sqrt{2h}\, \chi_0 - h\, \nabla U(Z_0)$,
    \begin{equation}
    \begin{aligned}
        & \bigl( \nabla^2 G(Y_0 + D_0) - \nabla^2 G(Y_0) \bigr) : \mathcal C(Z_0) \\
        & \quad = \varrho_2^{D_0} : \mathcal C(Z_0) - h\, \nabla^3 G(Y_0)[\nabla U(Z_0)] : \mathcal C(Z_0) + \sqrt{2h}\, \nabla^3 G(Y_0)[\chi_0] : \mathcal C(Z_0) .
    \end{aligned}
    \label{eq: sg hessian split}
    \end{equation}
    Assumption~\ref{ass: U} gives
    \begin{equation*}
        |\nabla U(Z_0)| \Le M \bigl( 1 + |Z_0| \bigr) \Le M \Bigl( 1 + |Y_0| + \sqrt{\frac h2}\, |\chi_0| \Bigr) ,
    \end{equation*}
    hence
    \begin{equation}
        \hat{\mathbb E}_0 |D_0|^2 \Le 4h\, \hat{\mathbb E}_0 |\chi_0|^2 + 2h^2\, \hat{\mathbb E}_0 |\nabla U(Z_0)|^2 \Le C h (1+|Y_0|)^2 .
        \label{eq: sg step second moment}
    \end{equation}
    By \eqref{eq: trace bound}, $\Sigma \Le M^2$ and \eqref{eq: sg step second moment}, the terms of \eqref{eq: sg hessian split} carrying $\varrho_2^{D_0}$ and $\nabla U(Z_0)$ contribute
    \begin{equation}
    \begin{aligned}
        & \bigl| \hat{\mathbb E}_0[\varrho_2^{D_0} : \mathcal C(Z_0)] \bigr| \Le \frac12 \norm{\nabla^4 G}_\infty\, \Sigma\, \hat{\mathbb E}_0 |D_0|^2 \Le C h \norm{\nabla^4 G}_\infty (1+|Y_0|)^2 , \\
        & h\, \bigl| \hat{\mathbb E}_0\bigl[ \nabla^3 G(Y_0)[\nabla U(Z_0)] : \mathcal C(Z_0) \bigr] \bigr| \Le \Sigma h \norm{\nabla^3 G}_\infty\, \hat{\mathbb E}_0 |\nabla U(Z_0)| \Le C h \norm{\nabla^3 G}_\infty (1+|Y_0|) .
    \end{aligned}
    \label{eq: sg remainder terms}
    \end{equation}
    For the term of \eqref{eq: sg hessian split} carrying $\chi_0$, split
    \begin{equation}
        \nabla^3 G(Y_0)[\chi_0] : \mathcal C(Z_0) = \nabla^3 G(Y_0)[\chi_0] : \mathcal C(Y_0) + \nabla^3 G(Y_0)[\chi_0] : \bigl( \mathcal C(Z_0) - \mathcal C(Y_0) \bigr) .
        \label{eq: sg chi split}
    \end{equation}
    Since $\chi_0$ has mean $0$ and is independent of $\hat{\mathcal F}_0$, while $\nabla^3 G(Y_0)$ and $\mathcal C(Y_0)$ are $\hat{\mathcal F}_0$-measurable,
    \begin{equation}
        \hat{\mathbb E}_0 \bigl[ \nabla^3 G(Y_0)[\chi_0] : \mathcal C(Y_0) \bigr] = 0 ,
        \label{eq: sg centered term}
    \end{equation}
    and, since $Z_0 - Y_0 = \sqrt{h/2}\,\chi_0$, \eqref{eq: covariance lipschitz} gives
    \begin{equation*}
        \norm{\mathcal C(Z_0) - \mathcal C(Y_0)} \Le 4 M^2 \sqrt{\frac h2}\, |\chi_0| ,
    \end{equation*}
    and the facts
    \begin{equation*}
        |A : B| \Le d \norm{A} \norm{B} , \qquad \norm{\nabla^3 G(Y_0)[\chi_0]} \Le \norm{\nabla^3 G}_\infty |\chi_0| , \qquad \hat{\mathbb E}_0 |\chi_0|^2 = d
    \end{equation*}
    give
    \begin{equation}
    \begin{aligned}
        & \sqrt{2h}\, \Bigl| \hat{\mathbb E}_0 \bigl[ \nabla^3 G(Y_0)[\chi_0] : \bigl( \mathcal C(Z_0) - \mathcal C(Y_0) \bigr) \bigr] \Bigr| \Le \sqrt{2h}\, d\, \norm{\nabla^3 G}_\infty\, \hat{\mathbb E}_0\bigl[ |\chi_0|\, \norm{\mathcal C(Z_0) - \mathcal C(Y_0)} \bigr] \\
        & \quad \Le \sqrt{2h}\, d\, \norm{\nabla^3 G}_\infty \cdot 4 M^2 \sqrt{\frac h2}\, \hat{\mathbb E}_0 |\chi_0|^2 = 4 d^2 M^2 h \norm{\nabla^3 G}_\infty .
    \end{aligned}
    \label{eq: sg covariance shift}
    \end{equation}
    Inserting \eqref{eq: sg remainder terms}, \eqref{eq: sg centered term} and \eqref{eq: sg covariance shift} into \eqref{eq: sg hessian split} and \eqref{eq: sg chi split}, and multiplying by $\frac{h^2}{2}$, gives
    \begin{equation}
        \frac{h^2}{2} \Bigl| \hat{\mathbb E}_0\bigl[ \bigl( \nabla^2 G(Y_0 + D_0) - \nabla^2 G(Y_0) \bigr) : \mathcal C(Z_0) \bigr] \Bigr| \Le C h^3 \bigl( \norm{\nabla^3 G}_\infty + \norm{\nabla^4 G}_\infty \bigr) (1+|Y_0|)^2 .
        \label{eq: sg hessian shift}
    \end{equation}
    Finally, the left side of \eqref{eq: sg weak} at the index $0$ splits as
    \begin{align}
        & \quad \hat{\mathbb E}_0\bigl[ G(Y_1) - G(\tilde Y_0(h)) \bigr] - \frac{h^2}{2}\, \nabla^2 G(Y_0) : \hat{\mathbb E}_0\bigl[ \mathcal C(Z_0) \bigr] \notag \\
        & = \Bigl( \hat{\mathbb E}_0\bigl[ G(Y_1) \bigr] - \mathbb E_0\bigl[ G(Y_0 + D_0) \bigr] - \frac{h^2}{2}\, \hat{\mathbb E}_0\bigl[ \nabla^2 G(Y_0 + D_0) : \mathcal C(Z_0) \bigr] \Bigr) + \mathbb E_0\bigl[ G(Y_0 + D_0) - G(\tilde Y_0(h)) \bigr] \notag \\
        & \quad + \frac{h^2}{2}\, \hat{\mathbb E}_0\bigl[ \bigl( \nabla^2 G(Y_0 + D_0) - \nabla^2 G(Y_0) \bigr) : \mathcal C(Z_0) \bigr] .
        \label{eq: sg weak split}
    \end{align}
	On the right side of \eqref{eq: sg weak split}, \eqref{eq: sg noise expansion} bounds the first bracket, \eqref{eq: sg exact part} the second term, and \eqref{eq: sg hessian shift} the third term. Adding these bounds gives
    \begin{equation*}
    \begin{aligned}
        & \quad \Bigl| \hat{\mathbb E}_0\bigl[ G(Y_1) - G(\tilde Y_0(h)) \bigr] - \frac{h^2}{2}\, \nabla^2 G(Y_0) : \hat{\mathbb E}_0\bigl[ \mathcal C(Z_0) \bigr] \Bigr| \\
        & \Le C h^3 \norm{\nabla^3 G}_\infty (1+|Y_0|)^3 + C h^3 \sum_{j=1}^{4} \norm{\nabla^j G}_\infty (1+|Y_0|)^6 + C h^3 \bigl( \norm{\nabla^3 G}_\infty + \norm{\nabla^4 G}_\infty \bigr) (1+|Y_0|)^2 \\
        & \Le C h^3 \sum_{j=1}^{4} \norm{\nabla^j G}_\infty (1+|Y_0|)^6 ,
    \end{aligned}
    \end{equation*}
    which is \eqref{eq: sg weak} at the index $0$; by time homogeneity, \eqref{eq: sg weak} holds at every $k$.
\end{proof}

\section{Proofs of the MSE analysis}
\label{app: mse}

Throughout this appendix, $C$ depends only on $m$, $M$, $R$ and $d$, and $h \Le h_0$ with $h_0 \Le \min\{\lambda^{-1}, \frac12\}$, $\lambda$ the rate of Theorem~\ref{thm: u decay} for $\tilde U_h$, so that Corollary~\ref{cor: modified poisson} applies. By \eqref{eq: modified poisson bound}, $\tilde\phi_h$ is $CL$-Lipschitz.

\subsection{Proof of Theorem~\ref{thm: lm mse}: MSE of LM}
\label{sec: lm mse proof}

\begin{proof}[Proof of Theorem~\ref{thm: lm mse}]
    Let first $Nh \Ge 1$, so that $N \Ge 2$ since $h \Le \frac12$. The first step of \eqref{eq: lm} gives
    \begin{equation*}
        Z_k = Y_k + \sqrt{\frac h2}\,\chi_k , \quad k \Ge 1 , \qquad Y_1 := Z_0 - h \nabla U(Z_0) + \sqrt{\frac h2}\,\chi_0 ,
    \end{equation*}
    with $(Y_k)_{k \Ge 1}$ the chain \eqref{eq: lm markov} started from $Y_1$. Here $Y_1$ is independent of $(B_t - B_h)_{t \Ge h}$, and $\mathbb E |Y_1|^{12} \Le C (1 + |Z_0|)^{12}$. Splitting off the term $k = 0$ in \eqref{eq: mse},
    \begin{equation*}
        \MSE(N,h) \Le \frac{2}{N^2} \bigl| f(Z_0) - \pi(f) \bigr|^2 + 2 \Bigl( \frac{N-1}{N} \Bigr)^2\, \mathbb E \Biggl[ \biggl| \frac{1}{N-1} \sum_{k=1}^{N-1} f(Z_k) - \pi(f) \biggr|^2 \Biggr] ,
    \end{equation*}
    where
    \begin{equation*}
        \frac{2}{N^2} \bigl| f(Z_0) - \pi(f) \bigr|^2 \Le \frac{C L^2}{N^2} \bigl( 1+|Z_0| \bigr)^2 \Le \frac{C L^2}{Nh} \bigl( 1+|Z_0| \bigr)^2 ,
    \end{equation*}
    and the expectation of the time average over $Z_1, \dots, Z_{N-1}$ is the MSE of $N-1$ samples of a chain of the same form started from $Y_1$, with $\frac{1}{(N-1)h} \Le \frac{2}{Nh}$ for $N \Ge 2$. Conditioning on $Y_1$, it therefore suffices to prove \eqref{eq: lm mse} for the samples $Z_k = Y_k + \sqrt{h/2}\,\chi_k$, $k \Ge 0$, of the chain \eqref{eq: lm markov} started from a given $Y_0$, with $Y_0$ in place of $Z_0$, and then to apply it with $Y_1$ in place of $Y_0$ and take the expectation, with $\mathbb E(1 + |Y_1|)^2 \Le C (1 + |Z_0|)^2$ and $\mathbb E(1 + |Y_1|)^{12} \Le C (1 + |Z_0|)^{12}$. By \eqref{eq: mse split modified} and $\frac{h}{N} \Le \frac{1}{Nh}$ for $h \Le 1$, it suffices to prove
    \begin{equation}
        \mathbb E \Biggl[ \biggl| \frac1N \sum_{k=0}^{N-1} \bigl( \tilde f_h(Y_k) - \tilde\pi_h(\tilde f_h) \bigr) \biggr|^2 \Biggr] \Le C L^2 \Bigl( \bigl( 1 + |Y_0| \bigr)^2 \frac{1}{Nh} + \bigl( 1 + |Y_0| \bigr)^{12} h^4 \Bigr) .
        \label{eq: markov mse}
    \end{equation}
    By \eqref{eq: time average formula} and Cauchy's inequality, the left side of \eqref{eq: markov mse} is at most
    \begin{equation}
        \frac{3}{N^2h^2}\, \mathbb E\bigl| \tilde\phi_h(Y_0) - \tilde\phi_h(Y_N) \bigr|^2
        + \frac{3}{N^2}\, \mathbb E\biggl| \sum_{k=0}^{N-1} \Upsilon_k \biggr|^2
        + \frac{3}{N^2}\, \mathbb E\biggl| \sum_{k=0}^{N-1} \Psi_k \biggr|^2 .
        \label{eq: three terms}
    \end{equation}
    The terms $\Upsilon_k$ and $\Psi_k$ are $\mathcal F_{k+1}$-measurable, since $\tilde Y_k(h)$ is driven by the Brownian motion on $[kh, (k+1)h]$, and $\mathbb E_k \Psi_k = 0$ by \eqref{eq: modified poisson}. Hence the $\Psi_k$, and the centered terms $\Upsilon_k - \mathbb E_k \Upsilon_k$, are martingale differences with respect to $(\mathcal F_k)_{k \Ge 0}$, and their sums have additive second moments. Lemmas~\ref{lem: moments} and \ref{lem: sizes} also give
    \begin{equation}
        \mathbb E|Y_k|^2 \Le C \bigl( 1+|Y_0| \bigr)^2 , \qquad \mathbb E|Y_{k+1} - Y_k|^2 + \mathbb E|\tilde Y_k(h) - Y_k|^2 \Le C h \bigl( 1+|Y_0| \bigr)^2 .
        \label{eq: markov second moments}
    \end{equation}

    For the boundary term of \eqref{eq: three terms}, the Lipschitz bound on $\tilde\phi_h$ and \eqref{eq: markov second moments} give
    \begin{equation}
        \frac{3}{N^2h^2}\, \mathbb E\bigl| \tilde\phi_h(Y_0) - \tilde\phi_h(Y_N) \bigr|^2 \Le \frac{C L^2}{N^2h^2}\, \mathbb E|Y_N - Y_0|^2 \Le \frac{C L^2}{N^2 h^2} \bigl( 1+|Y_0| \bigr)^2 \Le \frac{C L^2}{Nh} \bigl( 1+|Y_0| \bigr)^2 ,
    \label{eq: boundary term bound}
    \end{equation}
    using $Nh \Ge \frac12$. This holds for the $N-1$ samples after the first step, since $(N-1)h \Ge 1 - h \Ge \frac12$.

    For the discretization term of \eqref{eq: three terms}, carrying $\Upsilon_k$ with $h \Upsilon_k = \tilde\phi_h(Y_{k+1}) - \tilde\phi_h(\tilde Y_k(h))$ by \eqref{eq: Upsilon Psi}, Lemma~\ref{lem: weak} with $G = \tilde\phi_h$ and \eqref{eq: modified poisson bound} give
    \begin{equation*}
        |\mathbb E_k \Upsilon_k| \Le C L h^2 \bigl( 1+|Y_k| \bigr)^6 ,
    \end{equation*}
    so by Cauchy's inequality and Lemma~\ref{lem: moments},
    \begin{equation*}
        \mathbb E\biggl| \sum_{k=0}^{N-1} \mathbb E_k \Upsilon_k \biggr|^2 \Le N \sum_{k=0}^{N-1} \mathbb E\bigl| \mathbb E_k \Upsilon_k \bigr|^2 \Le C L^2 N^2 h^4 \bigl( 1 + |Y_0| \bigr)^{12} .
    \end{equation*}
    For its fluctuation, the Lipschitz bound on $\tilde\phi_h$ and Lemma~\ref{lem: sizes} give
    \begin{equation*}
        h |\Upsilon_k| \Le C L\, \bigl| Y_{k+1} - \tilde Y_k(h) \bigr| , \qquad \mathbb E_k |\Upsilon_k|^2 \Le C L^2 h \bigl( 1+|Y_k| \bigr)^2 ,
    \end{equation*}
    so by the orthogonality of martingale differences,
    \begin{equation*}
        \mathbb E\biggl| \sum_{k=0}^{N-1} \bigl( \Upsilon_k - \mathbb E_k \Upsilon_k \bigr) \biggr|^2 = \sum_{k=0}^{N-1} \mathbb E\bigl| \Upsilon_k - \mathbb E_k \Upsilon_k \bigr|^2 \Le C L^2 N h \bigl( 1 + |Y_0| \bigr)^2 .
    \end{equation*}
    The discretization term of \eqref{eq: three terms} is therefore at most
    \begin{equation}
        \frac{3}{N^2}\, \mathbb E\biggl| \sum_{k=0}^{N-1} \Upsilon_k \biggr|^2 \Le C L^2 \Bigl( \bigl( 1+|Y_0| \bigr)^{12} h^4 + \bigl( 1+|Y_0| \bigr)^2 \frac hN \Bigr) .
        \label{eq: discretization term bound}
    \end{equation}

    For the martingale term of \eqref{eq: three terms}, carrying $\Psi_k$ with $h \Psi_k = \tilde\phi_h(\tilde Y_k(h)) - \mathbb E_k[\tilde\phi_h(\tilde Y_k(h))]$ by \eqref{eq: Upsilon Psi} and \eqref{eq: modified poisson}, the Lipschitz bound and \eqref{eq: markov second moments} give
    \begin{equation*}
        h^2\, \mathbb E|\Psi_k|^2 \Le C L^2\, \mathbb E\bigl| \tilde Y_k(h) - Y_k \bigr|^2 \Le C L^2 h \bigl( 1+|Y_0| \bigr)^2 ,
    \end{equation*}
    and by orthogonality,
    \begin{equation}
        \frac{3}{N^2}\, \mathbb E\biggl| \sum_{k=0}^{N-1} \Psi_k \biggr|^2 = \frac{3}{N^2} \sum_{k=0}^{N-1} \mathbb E|\Psi_k|^2 \Le \frac{C L^2}{Nh} \bigl( 1+|Y_0| \bigr)^2 .
        \label{eq: martingale term bound}
    \end{equation}

    Adding \eqref{eq: boundary term bound}, \eqref{eq: discretization term bound} and \eqref{eq: martingale term bound}, and using $\frac{h}{N} \Le \frac{1}{Nh}$ for $h \Le 1$, proves \eqref{eq: markov mse}.

    Let finally $Nh < 1$. By Assumption~\ref{ass: f} and the moments of $\pi$ in Lemma~\ref{lem: dissipative}, $|f(x) - \pi(f)| \Le C L (1 + |x|)$, and $\mathbb E|Z_k|^2 \Le C (1 + |Z_0|)^2$ for every $k \Ge 0$, by the moment bound on $Y_1$ and Lemma~\ref{lem: moments}. Hence \eqref{eq: mse} and Cauchy's inequality give
    \begin{equation*}
        \MSE(N,h) \Le \frac1N \sum_{k=0}^{N-1} \mathbb E\bigl| f(Z_k) - \pi(f) \bigr|^2 \Le C L^2 \bigl( 1 + |Z_0| \bigr)^2 \Le \frac{C L^2}{Nh} \bigl( 1 + |Z_0| \bigr)^2 ,
    \end{equation*}
    which is at most the right side of \eqref{eq: lm mse}.
\end{proof}

\subsection{Proof of Theorem~\ref{thm: sglm mse}: MSE of SGLM}
\label{sec: sglm mse proof}

\begin{proof}[Proof of Theorem~\ref{thm: sglm mse}]
    As in the proof of Theorem~\ref{thm: lm mse}, with $Y_1 := Z_0 - h \nabla \hat U(Z_0;\xi_0) + \sqrt{h/2}\,\chi_0$, it suffices to bound the MSE of the samples $Z_k = Y_k + \sqrt{h/2}\,\chi_k$, $k \Ge 0$, of the chain \eqref{eq: sglm markov} started from a given $Y_0$ by the right side of \eqref{eq: sglm mse} with $Y_0$ in place of $Z_0$, and then to apply it with $Y_1$ in place of $Y_0$ and take the expectation, with $\mathbb E(1 + |Y_1|)^2 \Le C (1 + |Z_0|)^2$ and $\mathbb E(1 + |Y_1|)^{12} \Le C (1 + |Z_0|)^{12}$ by Assumption~\ref{ass: sg}. The bound \eqref{eq: mse split modified} and the formula \eqref{eq: time average formula} involve the modified test function, the modified dynamics and $\tilde\phi_h$ only, and hold for this chain with $\hat{\mathcal F}_k$ and $\hat{\mathbb E}_k$ in place of $\mathcal F_k$ and $\mathbb E_k$. The proof of Theorem~\ref{thm: lm mse} is repeated with two changes.

    First, the moment bounds of Lemma~\ref{lem: moments} hold for \eqref{eq: sglm markov} with the same constants, so the boundary term and the martingale term are bounded by $C L^2 (1+|Y_0|)^2 / (Nh)$ as before.

    Second, in the discretization term, Lemma~\ref{lem: sg weak} with $G = \tilde\phi_h$, \eqref{eq: modified poisson bound} and \eqref{eq: trace bound} give
    \begin{equation*}
        \bigl| \hat{\mathbb E}_k \Upsilon_k \bigr| \Le C L \Bigl( \Sigma h + h^2 \bigl( 1 + |Y_k| \bigr)^6 \Bigr) ,
    \end{equation*}
    since $\frac{h^2}{2} |\nabla^2 \tilde\phi_h(Y_k) : \hat{\mathbb E}_k[\mathcal C(Z_k)]| \Le C L \Sigma h^2$. Cauchy's inequality and Lemma~\ref{lem: moments} for the term $h^2 (1 + |Y_k|)^6$ give
    \begin{equation*}
        \mathbb E\biggl| \sum_{k=0}^{N-1} \hat{\mathbb E}_k \Upsilon_k \biggr|^2 \Le C L^2 N^2 \Bigl( \Sigma^2 h^2 + h^4 \bigl( 1 + |Y_0| \bigr)^{12} \Bigr) .
    \end{equation*}
    For its fluctuation, the step of \eqref{eq: sglm markov} differs from that of \eqref{eq: lm markov} by $-h \bigl( \nabla \hat U(Z_k;\xi_k) - \nabla U(Z_k) \bigr)$, and
    \begin{equation*}
        \hat{\mathbb E}_k \bigl| h \bigl( \nabla \hat U(Z_k;\xi_k) - \nabla U(Z_k) \bigr) \bigr|^2 \Le \Sigma h^2 \Le M^2 h^2 .
    \end{equation*}
    With the Lipschitz bound on $\tilde\phi_h$ and Lemma~\ref{lem: sizes}, this gives
    \begin{equation*}
        \hat{\mathbb E}_k |\Upsilon_k|^2 \Le C L^2 \Bigl( h \bigl( 1+|Y_k| \bigr)^2 + \Sigma \Bigr) \Le C L^2 \bigl( 1+|Y_k| \bigr)^2 ,
    \end{equation*}
    and, by orthogonality and $h \Le 1$,
    \begin{equation*}
        \frac{3}{N^2}\, \mathbb E\biggl| \sum_{k=0}^{N-1} \bigl( \Upsilon_k - \hat{\mathbb E}_k \Upsilon_k \bigr) \biggr|^2 \Le \frac{C L^2}{N} \bigl( 1+|Y_0| \bigr)^2 \Le \frac{C L^2}{Nh} \bigl( 1+|Y_0| \bigr)^2 .
    \end{equation*}
    Adding the bounds proves \eqref{eq: sglm mse} for $Nh \Ge 1$. For $Nh < 1$, it follows as at the end of the proof of Theorem~\ref{thm: lm mse}, since the moment bounds of Lemma~\ref{lem: moments} hold for \eqref{eq: sglm markov} and $\mathbb E(1 + |Y_1|)^2 \Le C (1 + |Z_0|)^2$.

    For the average of $P$ independent estimators, the average is unbiased, its moment of order $12$ satisfies the bound of Assumption~\ref{ass: sg} with the same $M$ by Minkowski's inequality, and its variance is at most $\Sigma / P \Le M^2$. Minkowski's inequality in $L^2$ also gives the Hessian bound of Assumption~\ref{ass: sg} with the same $M$. Assumption~\ref{ass: sg} therefore holds with $\Sigma / P$ in place of $\Sigma$, and \eqref{eq: sglm mse} gives the term $\Sigma^2 h^2 / P^2$.
\end{proof}

\section{Proofs of the numerical bias estimate}
\label{app: bias}

Throughout this appendix, $(Y_k)_{k \Ge 0}$ is the chain \eqref{eq: sglm markov}, $\mathcal P_h$ its transition kernel, and $h \Le h_0$ with $M h_0 \Le 1$. The local Hessian bound \eqref{eq: hessian sup} and the bounds \eqref{eq: covariance lipschitz} and \eqref{eq: trace bound} on the covariance $\mathcal C(x)$ of the gradient noise are those of Appendix~\ref{sec: sg weak proof}.

\subsection{Proof of Lemma~\ref{lem: invariant}: invariant distribution}
\label{sec: invariant proof}

\begin{proof}[Proof of Lemma~\ref{lem: invariant}]
    The proof obtains the existence and uniqueness of $\mu_h$ from Harris's ergodic theorem, deduces the moment bound of $\mu_h$ from Lemma~\ref{lem: moments}, and proves the limit of time averages first for the chain \eqref{eq: sglm markov} and then, in \eqref{eq: ergodic limit}, for the samples of \eqref{eq: sglm}. Harris's ergodic theorem in the form of \citep[Theorem~1.2]{hairer2008yet} assumes that a measurable $W \colon \mathbb R^d \to [0,\infty)$, constants $\gamma, \alpha_h \in (0,1)$, $K \Ge 0$, a level $R_W > 0$ and a probability measure $\nu$ satisfy the drift condition, the minorization on the level set $\{W \Le R_W\}$, and the condition on the level,
    \begin{equation}
        (\mathcal P_h W)(y) \Le \gamma W(y) + K , \quad y \in \mathbb R^d , \qquad
        \mathcal P_h(y, \cdot) \Ge \alpha_h\, \nu , \quad W(y) \Le R_W , \qquad
        R_W > \frac{2K}{1-\gamma} .
        \label{eq: harris conditions}
    \end{equation}
    That theorem then gives a unique invariant distribution $\mu_h$ and constants $C_h > 0$ and $\bar\gamma \in (0,1)$, which may depend on $h$, with
    \begin{equation}
        \sup_{y \in \mathbb R^d} \frac{\bigl|(\mathcal P_h^n \varphi)(y) - \mu_h(\varphi)\bigr|}{1 + W(y)} \Le C_h\, \bar\gamma^{\,n} \sup_{y \in \mathbb R^d} \frac{|\varphi(y) - \mu_h(\varphi)|}{1 + W(y)} , \qquad n \Ge 0 ,
        \label{eq: harris}
    \end{equation}
    for every measurable $\varphi$ with $\sup_{y \in \mathbb R^d} |\varphi(y)|/(1 + W(y)) < \infty$.

    For the drift condition in \eqref{eq: harris conditions}, take $W(y) := |y|^2$. The Lyapunov estimate \eqref{eq: lyapunov} with $q = 1$, which holds for \eqref{eq: sglm markov} as shown in the proof of Lemma~\ref{lem: moments}, gives
    \begin{equation*}
        (\mathcal P_h W)(y) \Le \Bigl( 1 - \frac{m}{16} h \Bigr) W(y) + C h , \qquad y \in \mathbb R^d .
    \end{equation*}
    This is the drift condition in \eqref{eq: harris conditions} with
    \begin{equation*}
        \gamma := 1 - \frac{m}{16} h , \qquad K := C h ,
    \end{equation*}
    where $\gamma \in (0,1)$ since $m \Le M$ by Assumption~\ref{ass: U} and $Mh \Le 1$. With the constant $C$ in $K = C h$, set
    \begin{equation*}
        R_W := \frac{64 C}{m} > \frac{32 C}{m} = \frac{2K}{1 - \gamma} ,
    \end{equation*}
    which depends only on $m$, $M$, $R$ and $d$.

    Next we prove the minorization on the level set $\{W \Le R_W\}$, with a constant $\alpha_h$ that may depend on $h$. Let $h_0$ also satisfy $4 M h_0 (2 + R_W^{\frac12}) \Le 1$, let $\mathbb B$ be the closed unit ball and $\nu$ the uniform distribution on $\mathbb B$, and fix $y$ with $|y|^2 \Le R_W$. From $Y_0 = y$, the step \eqref{eq: sglm markov} reads
    \begin{equation}
        Y_1 = T_{\xi_0}\Bigl( \frac y2 + \sqrt{\frac h2}\,\chi_0 \Bigr) , \qquad T_\xi(v) := 2v - h \nabla \hat U\Bigl( \frac y2 + v; \xi \Bigr) , \qquad v \in \mathbb R^d .
        \label{eq: integrator map}
    \end{equation}
    By Markov's inequality, \eqref{eq: hessian sup} and the moment bound of Assumption~\ref{ass: sg}, the event $E_0 := \{ \xi : S(\frac y2;\xi) \Le 4M ,\ |\nabla \hat U(\frac y2;\xi)| \Le 4M (1 + \frac{|y|}{2}) \}$ satisfies $\mathbb P_\xi(E_0^c) \Le \frac1{16} + 4^{-12} \Le \frac12$. For $\xi \in E_0$ and $v, v' \in \mathbb B$,
    \begin{equation*}
        h \Bigl| \nabla \hat U\Bigl( \frac y2 + v;\xi \Bigr) - \nabla \hat U\Bigl( \frac y2 + v';\xi \Bigr) \Bigr| \Le 4Mh\, |v - v'| , \qquad
        h \Bigl| \nabla \hat U\Bigl( \frac y2 + v;\xi \Bigr) \Bigr| \Le 4Mh \Bigl( 2 + \frac{|y|}{2} \Bigr) \Le 1 .
    \end{equation*}
    Hence $T_\xi$ is Lipschitz with constant $3$ on $\mathbb B$, and for every $w \in \mathbb B$ the map
    \begin{equation*}
        v \mapsto \frac12 \Bigl( w + h \nabla \hat U\Bigl( \frac y2 + v;\xi \Bigr) \Bigr)
    \end{equation*}
    maps $\mathbb B$ into itself with Lipschitz constant $\frac12$, so, by the contraction mapping theorem, it has a fixed point, which solves $T_\xi(v) = w$; hence $\mathbb B \subset T_\xi(\mathbb B)$. For a Borel set $A$, the set $A \cap \mathbb B$ is therefore contained in $T_\xi(T_\xi^{-1}(A) \cap \mathbb B)$, whose outer Lebesgue measure is at most $3^d\, \mathrm{Leb}(T_\xi^{-1}(A) \cap \mathbb B)$ by the Lipschitz bound for the $d$-dimensional Hausdorff measure, which equals the Lebesgue measure $\mathrm{Leb}$ on $\mathbb R^d$; on $\mathbb B$, whose points lie within $1 + \frac12 R_W^{\frac12}$ of $\frac y2$, the density of $\frac y2 + \sqrt{h/2}\,\chi_0 \sim \mathcal N(\frac y2, \frac h2 I_d)$ is at least $\beta_h := (\pi h)^{-\frac d2} \exp\bigl( -(1 + \frac12 R_W^{\frac12})^2 / h \bigr)$. Conditioning \eqref{eq: integrator map} on $\xi_0$, which is independent of $\chi_0$, gives
    \begin{equation*}
        \mathcal P_h(y, A) \Ge \mathbb E_\xi\biggl[ \mathbf 1_{\{\xi \in E_0\}}\, \mathbb P\Bigl( \frac y2 + \sqrt{\frac h2}\,\chi_0 \in T_\xi^{-1}(A) \cap \mathbb B \Bigr) \biggr] \Ge \frac{\beta_h}{2 \cdot 3^d}\, \mathrm{Leb}(A \cap \mathbb B) = \frac{\beta_h\, \mathrm{Leb}(\mathbb B)}{2 \cdot 3^d}\, \nu(A) ,
    \end{equation*}
    which is the minorization with $\alpha_h := \min\{ \frac12 , \beta_h\, \mathrm{Leb}(\mathbb B) / (2 \cdot 3^d) \}$, and \eqref{eq: harris} holds.

    For the moment bound, let $\varphi_\ell(y) := \min\{|y|^{12}, \ell\}$, $\ell \Ge 1$. Since $0 \Le \varphi_\ell \Le \ell$, \eqref{eq: harris} with $\varphi = \varphi_\ell$ and $n = k$ gives, for the chain from $Y_0 = 0$,
    \begin{equation*}
        \bigl| \mathbb E\big[ \varphi_\ell(Y_k) \big] - \mu_h(\varphi_\ell) \bigr| = \bigl| (\mathcal P_h^k \varphi_\ell)(0) - \mu_h(\varphi_\ell) \bigr| \Le \ell\, C_h\, \bar\gamma^{\,k} \to 0 , \qquad k \to \infty .
    \end{equation*}
    With Lemma~\ref{lem: moments} at $Y_0 = 0$, this gives
    \begin{equation*}
        \mu_h(\varphi_\ell) = \lim_{k \to \infty} \mathbb E\big[ \varphi_\ell(Y_k)\big] \Le \sup_{k \Ge 0} \mathbb E |Y_k|^{12} \Le C .
    \end{equation*}
    Monotone convergence as $\ell \to \infty$ gives $\mu_h(|y|^{12}) \Le C$.

    Next we prove the limit of time averages for the chain \eqref{eq: sglm markov}. By the minorization, the level set $\{W \Le R_W\}$ is petite \citep[Section~3.1]{meyn1992stability}. Outside this set, $W > R_W > \frac{2K}{1-\gamma}$, so the drift condition in \eqref{eq: harris conditions} gives $\mathcal P_h W \Le \frac{1+\gamma}{2}\, W$. For $Y_0 = y$ with $W(y) > R_W$ and the entrance time $\tau$ of the chain into $\{W \Le R_W\}$, the process $(\frac{1+\gamma}{2})^{-(k \wedge \tau)}\, W(Y_{k \wedge \tau})$ is therefore a nonnegative supermartingale; it exceeds $(\frac{1+\gamma}{2})^{-k} R_W$ on $\{\tau > k\}$ and has mean at most $W(y)$, so
    \begin{equation*}
        \mathbb P(\tau > k) \Le \Bigl( \frac{1+\gamma}{2} \Bigr)^{k} \frac{W(y)}{R_W} \to 0 , \qquad k \to \infty .
    \end{equation*}
    By the Markov property at the time $1$, the chain started from any initial value returns to $\{W \Le R_W\}$ almost surely, so it is Harris recurrent by \citep[Theorem~4.1(i)]{meyn1992stability}, and, since it has the invariant probability $\mu_h$, positive Harris recurrent \citep[Section~1.3]{meyn1992stability}. Let $g$ be measurable and of linear growth, so that $\mu_h(|g|) < \infty$ by the moment bound. The strong law of large numbers for positive Harris recurrent chains \citep[Theorem~17.0.1(i)]{meyn1993markov} gives, for every initial value $Y_0$,
    \begin{equation}
        \frac1N \sum_{k=0}^{N-1} g(Y_k) \to \mu_h(g) \qquad \text{almost surely} .
        \label{eq: chain limit}
    \end{equation}

    For the samples of \eqref{eq: sglm} from the given $Z_0$ of \eqref{eq: ergodic limit}, start \eqref{eq: sglm markov}, as in Section~\ref{subsec: mse analysis}, at
    \begin{equation*}
        Y_1 = Z_0 - h\nabla \hat U(Z_0;\xi_0) + \sqrt{\frac h2}\,\chi_0 , \qquad \text{so that} \qquad Z_k = Y_k + \sqrt{\frac h2}\,\chi_k , \quad k \Ge 1 .
    \end{equation*}
    Since $Y_1$ is independent of $(\chi_k, \xi_k)_{k \Ge 1}$, the sequence $(Y_{j+1})_{j \Ge 0}$ is the chain \eqref{eq: sglm markov} started at $Y_1$, and \eqref{eq: chain limit}, which holds for every initial value, gives by Fubini's theorem, almost surely,
    \begin{equation}
        \frac1N \sum_{j=0}^{N-1} g(Y_{j+1}) \to \mu_h(g) , \qquad N \to \infty .
        \label{eq: shifted chain limit}
    \end{equation}
    For $k \Ge 1$, as in Section~\ref{subsec: modified test function}, $\mathfrak M_k = f(Z_k) - \tilde f_h(Y_k)$, where $\hat{\mathbb E}_k[\mathfrak M_k] = 0$ by \eqref{eq: heat shift} and $\mathbb E[\mathfrak M_k^2] \Le \frac{dh}{2} L^2$ by \eqref{eq: martingale variance}, both with $\hat{\mathbb E}_k$ of \eqref{eq: sg filtration} in place of $\mathbb E_k$, since $Y_k$ is $\hat{\mathcal F}_k$-measurable and $\chi_k$ is independent of $\hat{\mathcal F}_k$; the $\mathfrak M_k$ are therefore $\hat{\mathcal F}_{k+1}$-measurable martingale differences. By \eqref{eq: shifted chain limit} with $g = \tilde f_h$, which is Lipschitz with constant $L$ by Section~\ref{subsec: modified test function}, hence of linear growth,
    \begin{equation*}
        \frac1N \sum_{k=1}^{N-1} \tilde f_h(Y_k) \to \mu_h(\tilde f_h) \qquad \text{almost surely} .
    \end{equation*}
    By the bound on $\mathbb E [ \mathfrak M_k^2 ]$, the martingale $\sum_{k=1}^{N} \mathfrak M_k / k$ is bounded in $L^2$ and converges almost surely, and Kronecker's lemma gives
    \begin{equation*}
        \frac1N \sum_{k=1}^{N-1} \mathfrak M_k \to 0 \qquad \text{almost surely} .
    \end{equation*}
    Let $Y \sim \mu_h$ and $\chi \sim \mathcal N(0, I_d)$ be independent. Fubini's theorem, \eqref{eq: heat shift} and the definition of $\pi_h$ in Lemma~\ref{lem: invariant} give
    \begin{equation}
        \mu_h(\tilde f_h) = \mathbb E \biggl[ f\Big( Y + \sqrt{\frac h2}\,\chi \Big) \biggr] = \pi_h(f) .
        \label{eq: heat shift invariant}
    \end{equation}
    Combining the limits of the averages of $\tilde f_h(Y_k)$ and of $\mathfrak M_k$ with \eqref{eq: heat shift invariant} and with $f(Z_0)/N \to 0$, since $Z_0$ is given,
    \begin{equation*}
        \frac1N \sum_{k=0}^{N-1} f(Z_k) = \frac{f(Z_0)}{N} + \frac1N \sum_{k=1}^{N-1} \tilde f_h(Y_k) + \frac1N \sum_{k=1}^{N-1} \mathfrak M_k \to \pi_h(f) \qquad \text{almost surely} ,
    \end{equation*}
    which is \eqref{eq: ergodic limit}.
\end{proof}

\subsection{Proof of Theorem~\ref{thm: numerical bias}: numerical bias}
\label{sec: numerical bias proof}

\begin{proof}[Proof of Theorem~\ref{thm: numerical bias}]
    The proof writes $\mu_h(\tilde f_h) - \tilde\pi_h(\tilde f_h)$ through the modified discrete Poisson equation \eqref{eq: modified poisson} as the local weak error of Lemma~\ref{lem: sg weak} averaged over $\mu_h$, and then adds the bound of Lemma~\ref{lem: shifted bias} on $\tilde\pi_h(\tilde f_h) - \pi(f)$.

    Let $Y_0 \sim \mu_h$, the invariant distribution of \eqref{eq: sglm markov} given by Lemma~\ref{lem: invariant}, independent of $(B_t)_{t \Ge 0}$ and $(\xi_k)_{k \Ge 0}$. Then $Y_1 \sim \mu_h$. Add $Y_0$ to the generators of $\hat{\mathcal F}_k$ in \eqref{eq: sg filtration}, so that $\hat{\mathcal F}_0 = \sigma(Y_0)$, while $\chi_k$ and $\xi_k$ remain independent of $\hat{\mathcal F}_k$. The process $\tilde Y_0(s)$ evolves by the modified dynamics \eqref{eq: modified langevin}, with the generator $\tilde{\mathcal L}_h$ of \eqref{eq: modified generator}, from $\tilde Y_0(0) = Y_0$, so that
    \begin{equation*}
        (e^{h\tilde{\mathcal L}_h} \tilde\phi_h)(Y_0) = \hat{\mathbb E}_0\bigl[ \tilde\phi_h(\tilde Y_0(h)) \bigr] .
    \end{equation*}
    The modified discrete Poisson equation \eqref{eq: modified poisson} at $Y_0$ therefore reads
    \begin{equation*}
        \tilde f_h(Y_0) - \tilde\pi_h(\tilde f_h) = \frac1h\, \hat{\mathbb E}_0\bigl[ \tilde\phi_h(Y_0) - \tilde\phi_h(\tilde Y_0(h)) \bigr] .
    \end{equation*}
    Taking expectations, with $\mathbb E\, \tilde f_h(Y_0) = \mu_h(\tilde f_h)$, gives
    \begin{equation}
        \mu_h(\tilde f_h) - \tilde\pi_h(\tilde f_h) = \frac1h\, \mathbb E\bigl[ \tilde\phi_h(Y_0) - \tilde\phi_h(\tilde Y_0(h)) \bigr] = \frac1h\, \mathbb E\bigl[ \tilde\phi_h(Y_1) - \tilde\phi_h(\tilde Y_0(h)) \bigr] ,
        \label{eq: bias identity}
    \end{equation}
    where $\mathbb E\, \tilde\phi_h(Y_0) = \mathbb E\, \tilde\phi_h(Y_1)$ since $Y_1 \sim \mu_h$. All expectations in \eqref{eq: bias identity} are finite, since $\tilde\phi_h$ is Lipschitz with constant $C L$ by \eqref{eq: modified poisson bound}, $\mu_h(|y|^{12}) \Le C$ by Lemma~\ref{lem: invariant}, and $\mathbb E|\tilde Y_0(h) - Y_0| \Le C \sqrt h\, ( 1 + \mathbb E|Y_0| )$ by the bound on $\tilde Y_0(s) - Y_0$ in \eqref{eq: exact sizes}.

    Next we expand $\mathbb E\bigl[ \tilde\phi_h(Y_1) - \tilde\phi_h(\tilde Y_0(h)) \bigr]$ in \eqref{eq: bias identity} by Lemma~\ref{lem: sg weak}. By \eqref{eq: modified poisson bound}, the derivatives of $\tilde\phi_h$ of orders $1$ to $4$ are bounded by $C L$, so \eqref{eq: sg weak} with $G = \tilde\phi_h$ at $k = 0$ gives
    \begin{equation*}
        \Bigl| \hat{\mathbb E}_0\bigl[ \tilde\phi_h(Y_1) - \tilde\phi_h(\tilde Y_0(h)) \bigr] - \frac{h^2}{2}\, \nabla^2 \tilde\phi_h(Y_0) : \hat{\mathbb E}_0\bigl[ \mathcal C(Z_0) \bigr] \Bigr| \Le C L h^3 \bigl( 1 + |Y_0| \bigr)^6 .
    \end{equation*}
    Take expectations in the bound on $\hat{\mathbb E}_0\bigl[ \tilde\phi_h(Y_1) - \tilde\phi_h(\tilde Y_0(h)) \bigr]$. Its term of order $h^2$ carries the $\hat{\mathcal F}_0$-measurable factor $\nabla^2 \tilde\phi_h(Y_0)$, so the tower property gives
    \begin{equation*}
        \mathbb E\bigl[ \nabla^2 \tilde\phi_h(Y_0) : \hat{\mathbb E}_0[\mathcal C(Z_0)] \bigr] = \mathbb E\bigl[ \nabla^2 \tilde\phi_h(Y_0) : \mathcal C(Z_0) \bigr] ,
    \end{equation*}
    and hence
    \begin{equation*}
        \Bigl| \mathbb E\bigl[ \tilde\phi_h(Y_1) - \tilde\phi_h(\tilde Y_0(h)) \bigr] - \frac{h^2}{2}\, \mathbb E\bigl[ \nabla^2 \tilde\phi_h(Y_0) : \mathcal C(Z_0) \bigr] \Bigr| \Le C L h^3\, \mathbb E \bigl( 1 + |Y_0| \bigr)^6 .
    \end{equation*}
    Dividing by $h$ and inserting \eqref{eq: bias identity},
    \begin{equation}
        \Bigl| \mu_h(\tilde f_h) - \tilde\pi_h(\tilde f_h) - \frac h2\, \mathbb E\bigl[ \nabla^2 \tilde\phi_h(Y_0) : \mathcal C(Z_0) \bigr] \Bigr| \Le C L h^2\, \mathbb E \bigl( 1 + |Y_0| \bigr)^6 \Le C L h^2 ,
        \label{eq: shifted numerical bias}
    \end{equation}
    by the moment bound $\mu_h(|y|^{12}) \Le C$ of Lemma~\ref{lem: invariant}. By \eqref{eq: heat shift invariant}, $\pi_h(f) = \mu_h(\tilde f_h)$, so
    \begin{equation}
    \begin{aligned}
        & \quad \pi_h(f) - \pi(f) - \frac h2\, \mathbb E\bigl[ \nabla^2 \tilde\phi_h(Y_0) : \mathcal C(Z_0) \bigr] \\
        & = \Bigl( \mu_h(\tilde f_h) - \tilde\pi_h(\tilde f_h) - \frac h2\, \mathbb E\bigl[ \nabla^2 \tilde\phi_h(Y_0) : \mathcal C(Z_0) \bigr] \Bigr) + \bigl( \tilde\pi_h(\tilde f_h) - \pi(f) \bigr) .
    \end{aligned}
        \label{eq: bias split}
    \end{equation}
    The first bracket of \eqref{eq: bias split} is bounded by \eqref{eq: shifted numerical bias}, and the second by \eqref{eq: shifted bias} of Lemma~\ref{lem: shifted bias}, which proves \eqref{eq: numerical bias}. Finally, by \eqref{eq: trace bound} and \eqref{eq: modified poisson bound},
    \begin{equation*}
        \frac h2 \Bigl| \mathbb E\bigl[ \nabla^2 \tilde\phi_h(Y_0) : \mathcal C(Z_0) \bigr] \Bigr| \Le \frac{\Sigma h}{2}\, \norm{\nabla^2 \tilde\phi_h}_\infty \Le C L \Sigma h ,
    \end{equation*}
    which with \eqref{eq: numerical bias} gives
    \begin{equation*}
        |\pi_h(f) - \pi(f)| \Le C L (\Sigma h + h^2) .
    \end{equation*}
\end{proof}

\subsection{Proof of Corollary~\ref{cor: bias coefficient}: coefficient of the numerical bias}
\label{sec: bias coefficient proof}

\begin{proof}[Proof of Corollary~\ref{cor: bias coefficient}]
    With $Y_0 \sim \mu_h$ and $Z_0$ as in Theorem~\ref{thm: numerical bias}, let
    \begin{equation}
        a_h := \mathbb E\bigl[ \nabla^2 \tilde\phi_h(Y_0) : \mathcal C(Z_0) \bigr] .
        \label{eq: coefficient mean}
    \end{equation}
    Dividing \eqref{eq: numerical bias} by $h$ gives
    \begin{equation*}
        \biggl| \frac{\pi_h(f) - \pi(f)}{h} - \frac12\, a_h \biggr| \Le C L h ,
    \end{equation*}
    so it suffices to prove
    \begin{equation*}
        a_h \to \pi\bigl( \nabla^2 \phi_0(x) : \mathcal C(x) \bigr) , \qquad h \to 0 .
    \end{equation*}
    Since $Z_0 = Y_0 + \sqrt{h/2}\,\chi_0$ with $\chi_0 \sim \mathcal N(0, I_d)$ independent of $Y_0 \sim \mu_h$, conditioning on $Y_0$ in \eqref{eq: coefficient mean} gives
    \begin{equation}
        a_h = \mathbb E\Bigl[ \nabla^2 \tilde\phi_h(Y_0) : \mathbb E\bigl[ \mathcal C(Z_0) \,\big|\, Y_0 \bigr] \Bigr] = \mu_h(A_h) ,
        \label{eq: coefficient conditioning}
    \end{equation}
    with
    \begin{equation}
        A_h(y) := \nabla^2 \tilde\phi_h(y) : \tilde{\mathcal C}_h(y) , \qquad \tilde{\mathcal C}_h(y) := \mathbb E\bigg[ \mathcal C\Big( y + \sqrt{\frac h2}\, \chi \Big) \bigg] , \quad \chi \sim \mathcal N(0, I_d) .
        \label{eq: coefficient integrand}
    \end{equation}
    
    First, $\mu_h \to \pi$ weakly as $h\to0$. The family $(\mu_h)$ is tight: since $\mu_h(|y|^{12}) \Le C$ uniformly in $h$ by Lemma~\ref{lem: invariant}, Markov's inequality gives
    \begin{equation}
        \mu_h\bigl( |y| > \rho \bigr) \Le \rho^{-12}\, \mu_h\bigl( |y|^{12} \bigr) \Le C \rho^{-12} , \qquad \forall \rho > 0 .
        \label{eq: mu tail}
    \end{equation}
    Let $\varphi \in C_c^\infty(\mathbb R^d)$; it satisfies Assumption~\ref{ass: f} with some constant $L_\varphi$. By \eqref{eq: heat shift invariant} and \eqref{eq: shift expansion},
    \begin{equation*}
        \pi_h(\varphi) = \mu_h\bigl( e^{\frac h4\Delta} \varphi \bigr) , \qquad \bigl| (e^{\frac h4\Delta} \varphi)(y) - \varphi(y) \bigr| \Le C L_\varphi h , \quad y \in \mathbb R^d ,
    \end{equation*}
    and with Theorem~\ref{thm: numerical bias} applied to $\varphi$ these give
    \begin{equation*}
        |\mu_h(\varphi) - \pi(\varphi)| \Le \Bigl| \mu_h\bigl( \varphi - e^{\frac h4\Delta} \varphi \bigr) \Bigr| + |\pi_h(\varphi) - \pi(\varphi)| \Le C L_\varphi h + C L_\varphi (\Sigma h + h^2) .
    \end{equation*}
    Hence $\mu_h(\varphi) \to \pi(\varphi)$ for every $\varphi \in C_c^\infty(\mathbb R^d)$. Let $b$ be a bounded continuous function on $\mathbb R^d$, and let $\rho > 0$ and $\varepsilon > 0$. With the cutoff $\theta$ of \eqref{eq: modified potential}, which satisfies $\mathbf 1_{\{|y| \Le \rho\}} \Le \theta(y/\rho) \Le 1$, choose a mollification $\varphi \in C_c^\infty(\mathbb R^d)$ of $b\,\theta(\cdot/\rho)$ with $\norm{b\,\theta(\cdot/\rho) - \varphi}_\infty \Le \varepsilon$. The split
    \begin{equation*}
        b = b\,\bigl( 1 - \theta(\cdot/\rho) \bigr) + \bigl( b\,\theta(\cdot/\rho) - \varphi \bigr) + \varphi , \qquad \Bigl| b\,\bigl( 1 - \theta(\cdot/\rho) \bigr) \Bigr| \Le \norm{b}_\infty\, \mathbf 1_{\{|y| > \rho\}} ,
    \end{equation*}
    taken under $\mu_h$ and under $\pi$, gives
    \begin{equation*}
        |\mu_h(b) - \pi(b)| \Le 2\varepsilon + |\mu_h(\varphi) - \pi(\varphi)| + \norm{b}_\infty \bigl( \mu_h(|y| > \rho) + \pi(|y| > \rho) \bigr) ,
    \end{equation*}
    and \eqref{eq: mu tail}, with $\mu_h(\varphi) \to \pi(\varphi)$, gives
    \begin{equation*}
        \limsup_{h \to 0}\, |\mu_h(b) - \pi(b)| \Le 2\varepsilon + \norm{b}_\infty \bigl( C \rho^{-12} + \pi(|y| > \rho) \bigr) ,
    \end{equation*}
    which tends to $0$ as $\varepsilon \to 0$ and $\rho \to \infty$. So $\mu_h \to \pi$ weakly.

    Second, $\nabla^2 \tilde\phi_h \to \nabla^2 \phi_0$ uniformly on compact sets as $h \to 0$. Let
    \begin{equation}
        \psi_h(x) := \tilde\phi_h(x) - \tilde\phi_h(0) , \qquad x \in \mathbb R^d ,
        \label{eq: psi h}
    \end{equation}
    so that $\psi_h(0) = 0$. By \eqref{eq: modified poisson bound}, the derivatives of $\psi_h$ of orders $1$ to $4$ are bounded by $C L$ uniformly in $h$; hence, for $x, x' \in \mathbb R^d$ and $j = 0, 1, 2, 3$,
    \begin{equation*}
        \norm{\nabla^j \psi_h(x) - \nabla^j \psi_h(x')} \Le C L |x - x'| , \qquad |\psi_h(x)| \Le C L |x| ,
    \end{equation*}
    so, on every ball $\{|x| \Le \rho\}$ and uniformly in $h$,
    \begin{equation*}
        \norm{\nabla^j \psi_h(x)} \Le C L (1 + \rho) , \qquad j = 0, 1, 2, 3 .
    \end{equation*}
    The Arzel\`a--Ascoli theorem gives, for every sequence $h_n \to 0$, a subsequence, still written $h_n$, along which each $\nabla^j \psi_{h_n}$ converges uniformly on compact sets. Since the uniform limit of derivatives is the derivative of the limit, the limits of $\nabla^j \psi_{h_n}$ are $\nabla^j \psi$ for a function $\psi \in C^3(\mathbb R^d)$ with $\nabla \psi$, $\nabla^2 \psi$, $\nabla^3 \psi$ bounded by $C L$, and the uniform convergence on compact sets reads
    \begin{equation}
        \lim_{n \to \infty} \sup_{|x| \Le \rho}\Big\{\max_{0 \Le j \Le 3} \norm{\nabla^j \psi_{h_n}(x) - \nabla^j \psi(x)} \Big\} = 0 , \qquad \forall \rho > 0 .
        \label{eq: psi convergence}
    \end{equation}
    
    To identify the limit $\psi$, fix a deterministic initial point $Y_0 \in \mathbb R^d$. By the modified discrete Poisson equation \eqref{eq: modified poisson} at $Y_0$, with $(e^{h\tilde{\mathcal L}_h} \tilde\phi_h)(Y_0) = \mathbb E[\tilde\phi_h(\tilde Y_0(h))]$,
    \begin{equation}
        \mathbb E\bigl[ \tilde\phi_h(\tilde Y_0(h)) \bigr] - \tilde\phi_h(Y_0) = -h \bigl( \tilde f_h(Y_0) - \tilde\pi_h(\tilde f_h) \bigr) .
        \label{eq: modified poisson expectation}
    \end{equation}
    Since $\tilde Y_0(h)$ is driven by the modified dynamics \eqref{eq: modified langevin} from $Y_0$, with the generator $\tilde{\mathcal L}_h$ of \eqref{eq: modified generator}, It\^o's formula for $\tilde\phi_h$, whose stochastic integral is a martingale since $\nabla\tilde\phi_h$ is bounded, expands the difference on the left side of \eqref{eq: modified poisson expectation} as
    \begin{equation}
        \mathbb E\bigl[ \tilde\phi_h(\tilde Y_0(h)) \bigr] - \tilde\phi_h(Y_0) = \int_0^h \mathbb E\bigl[ \tilde g_h(\tilde Y_0(s)) \bigr] \D s ,
        \label{eq: modified dynkin}
    \end{equation}
    where, with $\psi_h$ of \eqref{eq: psi h},
    \begin{equation}
        \tilde g_h := \tilde{\mathcal L}_h \tilde\phi_h = \Delta \psi_h - \nabla \tilde U_h \cdot \nabla \psi_h .
        \label{eq: modified g}
    \end{equation}
    Compared with the modified discrete Poisson equation \eqref{eq: modified poisson}, $-\tilde g_h$ is the counterpart of $\tilde f_h - \tilde\pi_h(\tilde f_h)$, with the difference quotient $(\tilde\phi_h - e^{h\tilde{\mathcal L}_h} \tilde\phi_h)/h$ replaced by $-\tilde{\mathcal L}_h \tilde\phi_h$. Next we bound $\tilde g_h(Y_0) + \tilde f_h(Y_0) - \tilde\pi_h(\tilde f_h)$. Differentiating $\tilde g_h$ in \eqref{eq: modified g},
    \begin{equation*}
        \nabla \tilde g_h = \nabla \Delta \psi_h - \nabla^2 \tilde U_h\, \nabla \psi_h - \nabla^2 \psi_h\, \nabla \tilde U_h ,
    \end{equation*}
    and \eqref{eq: modified class bounds} gives
    \begin{equation*}
        \norm{\nabla^2 \tilde U_h(z)} \Le C , \qquad |\nabla \tilde U_h(z)| \Le C (1 + |z|) , \qquad z \in \mathbb R^d .
    \end{equation*}
    Hence the bounds of \eqref{eq: modified poisson bound} on the derivatives of $\psi_h$ of orders $1$ to $3$ give, for $z \in \mathbb R^d$ and $0 \Le s \Le h$,
    \begin{equation*}
        |\nabla \tilde g_h(z)| \Le C L (1 + |z|) , \qquad |\tilde g_h(\tilde Y_0(s)) - \tilde g_h(Y_0)| \Le C L \bigl( 1 + |Y_0| + |\tilde Y_0(s) - Y_0| \bigr) |\tilde Y_0(s) - Y_0| .
    \end{equation*}
    The bound on $\tilde Y_0(s) - Y_0$ in \eqref{eq: exact sizes} gives, for $0 \Le s \Le h \Le 1$,
    \begin{equation*}
        \mathbb E |\tilde Y_0(s) - Y_0| \Le C \sqrt h (1+|Y_0|) , \qquad \mathbb E |\tilde Y_0(s) - Y_0|^2 \Le C h (1+|Y_0|)^2 \Le C \sqrt h (1+|Y_0|)^2 .
    \end{equation*}
    Taking the expectation of the bound on $|\tilde g_h(\tilde Y_0(s)) - \tilde g_h(Y_0)|$ and inserting the bounds on $\mathbb E |\tilde Y_0(s) - Y_0|$ and $\mathbb E |\tilde Y_0(s) - Y_0|^2$ gives
    \begin{equation*}
        \Bigl| \mathbb E\bigl[ \tilde g_h(\tilde Y_0(s)) \bigr] - \tilde g_h(Y_0) \Bigr| \Le C L (1+|Y_0|)\, \mathbb E |\tilde Y_0(s) - Y_0| + C L\, \mathbb E |\tilde Y_0(s) - Y_0|^2 \Le C L \sqrt h\, (1+|Y_0|)^2 .
    \end{equation*}
    By \eqref{eq: modified poisson expectation} and \eqref{eq: modified dynkin}, averaging the bound on $\mathbb E\bigl[ \tilde g_h(\tilde Y_0(s)) \bigr] - \tilde g_h(Y_0)$ over $s \in [0,h]$ gives
    \begin{equation}
        \bigl| \tilde g_h(Y_0) + \tilde f_h(Y_0) - \tilde\pi_h(\tilde f_h) \bigr| = \biggl| \frac1h \int_0^h \Bigl( \mathbb E\bigl[ \tilde g_h(\tilde Y_0(s)) \bigr] - \tilde g_h(Y_0) \Bigr) \D s \biggr| \Le C L \sqrt h\, (1+|Y_0|)^2 , \quad Y_0 \in \mathbb R^d .
        \label{eq: g consistency}
    \end{equation}
    To pass to the limit in \eqref{eq: g consistency} along $h_n$, we write, with $\mathcal L$ as in \eqref{eq: generator} and $\tilde g_h$ as in \eqref{eq: modified g},
    \begin{equation}
        \tilde g_{h_n} - \mathcal L \psi = \bigl( \Delta \psi_{h_n} - \Delta \psi \bigr) - \bigl( \nabla \tilde U_{h_n} - \nabla U \bigr) \cdot \nabla \psi_{h_n} - \nabla U \cdot \bigl( \nabla \psi_{h_n} - \nabla \psi \bigr) ,
        \label{eq: g limit split}
    \end{equation}
    and, for every $x \in \mathbb R^d$, with $|\nabla \psi_{h_n}| \Le C L$ by \eqref{eq: modified poisson bound}, \eqref{eq: g limit split} gives
    \begin{equation}
        \bigl| \tilde g_{h_n}(x) - \mathcal L \psi(x) \bigr| \Le \bigl| \Delta \psi_{h_n}(x) - \Delta \psi(x) \bigr| + C L\, \bigl| \nabla \tilde U_{h_n}(x) - \nabla U(x) \bigr| + |\nabla U(x)|\, \bigl| \nabla \psi_{h_n}(x) - \nabla \psi(x) \bigr| .
        \label{eq: g limit bound}
    \end{equation}
    On the right side of \eqref{eq: g limit bound}, the terms in $\psi_{h_n} - \psi$ vanish by \eqref{eq: psi convergence} at $j = 1$ and $j = 2$, and the term in $\nabla \tilde U_{h_n} - \nabla U$ vanishes by \eqref{eq: modified class}:
    \begin{equation*}
    \begin{gathered}
        \lim_{n \to \infty} \bigl| \Delta \psi_{h_n}(x) - \Delta \psi(x) \bigr| \Le d \lim_{n \to \infty} \norm{\nabla^2 \psi_{h_n}(x) - \nabla^2 \psi(x)} = 0 , \\
        \lim_{n \to \infty} \bigl| \nabla \psi_{h_n}(x) - \nabla \psi(x) \bigr| = 0 , \qquad \lim_{n \to \infty} \bigl| \nabla \tilde U_{h_n}(x) - \nabla U(x) \bigr| \Le \lim_{n \to \infty} C h_n (1+|x|) = 0 .
    \end{gathered}
    \end{equation*}
    Hence $\lim_{n \to \infty} \tilde g_{h_n}(x) = \mathcal L \psi(x)$, and \eqref{eq: shift expansion} and \eqref{eq: shifted bias} give
    \begin{equation*}
        \lim_{n \to \infty} \bigl( \tilde f_{h_n}(x) - \tilde\pi_{h_n}(\tilde f_{h_n}) \bigr) = f(x) - \pi(f) .
    \end{equation*}
    Letting $n \to \infty$ in \eqref{eq: g consistency} at $h = h_n$ shows that $\psi$ solves the continuous Poisson equation \eqref{eq: continuous poisson}.
    
    It remains to show that $\nabla^2 \psi = \nabla^2 \phi_0$. The continuous Poisson equation \eqref{eq: continuous poisson} has the solution
    \begin{equation*}
        \phi_0 = \int_0^\infty u(\cdot, t)\, \D t .
    \end{equation*}
    The integral is well defined: by \eqref{eq: u bound}, $u(\cdot,t)$ decays like $e^{-\lambda t}$ uniformly on compact sets, and its derivatives of orders $1$ to $4$ decay like $e^{-\lambda t}$ uniformly on $\mathbb R^d$. Hence
    \begin{equation}
        \norm{\nabla^j \phi_0}_\infty \Le C L , \qquad j = 1, \dots, 4 ,
        \label{eq: continuous poisson bound}
    \end{equation}
    and, since $\partial_t u = \mathcal L u$ and $u(\cdot,0) = f - \pi(f)$, the integral solves \eqref{eq: continuous poisson}. The difference
    \begin{equation*}
        w := \psi - \phi_0 \in C^3(\mathbb R^d) \qquad \text{satisfies} \qquad \mathcal L w = 0 , \qquad \norm{\nabla w}_\infty < \infty .
    \end{equation*}
    Since $\mathcal L w = 0$ and $\nabla w$ is bounded, It\^o's formula gives $(e^{t\mathcal L} w)(x) = \mathbb E^x[w(X_t)] = w(x)$ for $t \Ge 0$.
    The contraction of \citep[Lemma~5]{ye2026mse}, which holds under Assumption~\ref{ass: U} and requires of its initial function only a bounded gradient, gives
    \begin{equation*}
        \norm{\nabla w}_\infty = \norm{\nabla (e^{t\mathcal L} w)}_\infty \Le C e^{-\lambda t} \norm{\nabla w}_\infty \to 0 , \qquad t \to \infty .
    \end{equation*}
    So $w$ is constant and $\nabla^2 \psi = \nabla^2 \phi_0$. The argument for $w$ uses only $\mathcal L w = 0$ and the bounded gradient of $w$, so \eqref{eq: continuous poisson} determines $\phi_0$ up to an additive constant among the functions with a bounded gradient, as stated in Section~\ref{subsec: poisson}. Since $\nabla^2 \psi_h = \nabla^2 \tilde\phi_h$ and $\nabla^2 \psi = \nabla^2 \phi_0$, \eqref{eq: psi convergence} at $j = 2$ shows that every sequence $h_n \to 0$ has a subsequence along which
    \begin{equation*}
        \lim_{n \to \infty} \sup_{|x| \Le \rho} \norm{\nabla^2 \tilde\phi_{h_n}(x) - \nabla^2 \phi_0(x)} = 0 , \qquad \forall \rho > 0 ,
    \end{equation*}
    and hence $\nabla^2 \tilde\phi_h \to \nabla^2 \phi_0$ uniformly on compact sets as $h \to 0$.

    Third, $\mu_h(A_h) \to \pi(A_0)$, with $A_h$ of \eqref{eq: coefficient integrand} and
    \begin{equation}
        A_0(y) := \nabla^2 \phi_0(y) : \mathcal C(y) , \qquad y \in \mathbb R^d .
        \label{eq: coefficient integrand limit}
    \end{equation}
    By \eqref{eq: covariance lipschitz} and the definition of $\tilde{\mathcal C}_h$ in \eqref{eq: coefficient integrand}, for every $y \in \mathbb R^d$,
    \begin{equation}
        \norm{\tilde{\mathcal C}_h(y) - \mathcal C(y)} \Le \mathbb E \biggl\| \mathcal C\Bigl( y + \sqrt{\frac h2}\, \chi \Bigr) - \mathcal C(y) \biggr\| \Le 4 M^2 \sqrt{\frac h2}\, \mathbb E|\chi| .
        \label{eq: covariance smoothing}
    \end{equation}
    By \eqref{eq: coefficient integrand} and \eqref{eq: coefficient integrand limit},
    \begin{equation}
        A_h(y) - A_0(y) = \nabla^2 \tilde\phi_h(y) : \bigl( \tilde{\mathcal C}_h(y) - \mathcal C(y) \bigr) + \bigl( \nabla^2 \tilde\phi_h(y) - \nabla^2 \phi_0(y) \bigr) : \mathcal C(y) .
        \label{eq: coefficient integrand split}
    \end{equation}
    For the term carrying $\tilde{\mathcal C}_h(y) - \mathcal C(y)$ in \eqref{eq: coefficient integrand split}, $|A : B| \Le d \norm{A} \norm{B}$, \eqref{eq: modified poisson bound}, \eqref{eq: covariance smoothing} and $\mathbb E|\chi| \Le \sqrt d$ give
    \begin{equation*}
        \Bigl| \nabla^2 \tilde\phi_h(y) : \bigl( \tilde{\mathcal C}_h(y) - \mathcal C(y) \bigr) \Bigr| \Le d\, \norm{\nabla^2 \tilde\phi_h}_\infty \norm{\tilde{\mathcal C}_h(y) - \mathcal C(y)} \Le d \cdot C L \cdot 4 M^2 \sqrt{\frac h2}\, \sqrt d \Le C L \sqrt h ,
    \end{equation*}
    and for the term carrying $\nabla^2 \tilde\phi_h(y) - \nabla^2 \phi_0(y)$, \eqref{eq: trace bound} gives
    \begin{equation*}
        \Bigl| \bigl( \nabla^2 \tilde\phi_h(y) - \nabla^2 \phi_0(y) \bigr) : \mathcal C(y) \Bigr| \Le \Sigma\, \norm{\nabla^2 \tilde\phi_h(y) - \nabla^2 \phi_0(y)} .
    \end{equation*}
    Hence
    \begin{equation}
        |A_h(y) - A_0(y)| \Le C L \sqrt h + \Sigma\, \norm{\nabla^2 \tilde\phi_h(y) - \nabla^2 \phi_0(y)} , \qquad y \in \mathbb R^d .
        \label{eq: coefficient integrand difference}
    \end{equation}
    With the uniform convergence of $\nabla^2 \tilde\phi_h$ to $\nabla^2 \phi_0$ on compact sets, \eqref{eq: coefficient integrand difference} gives $A_h \to A_0$ uniformly on compact sets. Next we bound $A_h$ and $A_0$ uniformly in $y$ and $h$. Since $A : \tilde{\mathcal C}_h(y) = \mathbb E\bigl[ A : \mathcal C(y + \sqrt{h/2}\,\chi) \bigr]$ by \eqref{eq: coefficient integrand}, \eqref{eq: trace bound} gives
    \begin{equation*}
        |A : \tilde{\mathcal C}_h(y)| \Le \Sigma \norm{A} , \qquad A \in \mathbb R^{d \times d} , \ y \in \mathbb R^d .
    \end{equation*}
    With \eqref{eq: modified poisson bound} and \eqref{eq: continuous poisson bound}, the bound on $A : \tilde{\mathcal C}_h(y)$ and \eqref{eq: trace bound} give
    \begin{equation*}
        |A_h(y)| + |A_0(y)| \Le \Sigma \bigl( \norm{\nabla^2 \tilde\phi_h}_\infty + \norm{\nabla^2 \phi_0}_\infty \bigr) \Le C L \Sigma , \qquad y \in \mathbb R^d ,
    \end{equation*}
    and $A_0$ is continuous, since $\nabla^2 \phi_0$ is continuous and $x \mapsto \mathcal C(x)$ is Lipschitz by \eqref{eq: covariance lipschitz}. For every $\rho > 0$,
    \begin{equation*}
        |\mu_h(A_h) - \pi(A_0)| \Le \Bigl| \mu_h\bigl( (A_h - A_0)\, \mathbf 1_{\{|y| \Le \rho\}} \bigr) \Bigr| + \Bigl| \mu_h\bigl( (A_h - A_0)\, \mathbf 1_{\{|y| > \rho\}} \bigr) \Bigr| + |\mu_h(A_0) - \pi(A_0)| ,
    \end{equation*}
    and the bound $|A_h| + |A_0| \Le C L \Sigma$ on the term carrying $\mathbf 1_{\{|y| > \rho\}}$, followed by \eqref{eq: mu tail}, gives
    \begin{equation*}
    \begin{aligned}
        |\mu_h(A_h) - \pi(A_0)| & \Le \sup_{|y| \Le \rho} |A_h(y) - A_0(y)| + C L \Sigma\, \mu_h\bigl( |y| > \rho \bigr) + |\mu_h(A_0) - \pi(A_0)| \\
        & \Le \sup_{|y| \Le \rho} |A_h(y) - A_0(y)| + C L \Sigma\, \rho^{-12} + |\mu_h(A_0) - \pi(A_0)| .
    \end{aligned}
    \end{equation*}
    As $h \to 0$, the uniform convergence of $A_h$ to $A_0$ on compact sets and the weak convergence $\mu_h \to \pi$, applied to the bounded continuous $A_0$, give
    \begin{equation*}
        \sup_{|y| \Le \rho} |A_h(y) - A_0(y)| \to 0 , \qquad |\mu_h(A_0) - \pi(A_0)| \to 0 ,
    \end{equation*}
    so, with $a_h = \mu_h(A_h)$ by \eqref{eq: coefficient conditioning} and $\pi(A_0) = \pi\bigl( \nabla^2 \phi_0(x) : \mathcal C(x) \bigr)$ by \eqref{eq: coefficient integrand limit},
    \begin{equation*}
        \limsup_{h \to 0}\, \Bigl| a_h - \pi\bigl( \nabla^2 \phi_0(x) : \mathcal C(x) \bigr) \Bigr| \Le C L \Sigma\, \rho^{-12} , \qquad \rho > 0 .
    \end{equation*}
    Letting $\rho \to \infty$ gives $a_h \to \pi\bigl( \nabla^2 \phi_0(x) : \mathcal C(x) \bigr)$, which proves \eqref{eq: bias coefficient}.
\end{proof}


\section{Proofs for the numerical bias comparison with SGLD}
\label{app: sgld}

Throughout this appendix, $(Z_k)_{k \Ge 0}$ is the chain \eqref{eq: sgld} and $\mathcal P_h$ its transition kernel. Since $(\chi_k)_{k \Ge 1}$ is independent of $(\xi_k)_{k \Ge 0}$, the law of the chain is unchanged if its Gaussian increments are taken as $\sqrt{2h}\, \chi_{k+1} = \sqrt2\, (B_{(k+1)h} - B_{kh})$, with $(B_t)_{t \Ge 0}$ a Brownian motion independent of $(\xi_k)_{k \Ge 0}$; then, for a given $Z_0$, $Z_k$ is $\hat{\mathcal F}_k$-measurable, with $\hat{\mathcal F}_k$ and $\hat{\mathbb E}_k$ of \eqref{eq: sg filtration}. Statements about a single step are made at the index $0$, with the solution $X_s$ of \eqref{eq: langevin} from $X_0 = Z_0$, driven by the same Brownian motion. Then
\begin{equation}
    Z_1 = Z_0 + D_0 - h\, \eta_0 , \qquad D_0 := \sqrt2\, B_h - h\, \nabla U(Z_0) , \qquad \eta_0 := \eta(Z_0;\xi_0) = \nabla \hat U(Z_0;\xi_0) - \nabla U(Z_0) ,
    \label{eq: sgld step}
\end{equation}
where $Z_0 + D_0$ is the one-step update of EM from $Z_0$ and $\eta_0$ is the gradient noise $\eta(x;\xi)$ of Appendix~\ref{sec: sg weak proof} evaluated at $Z_0$ and $\xi_0$.

\subsection{Invariant distribution of SGLD}
\label{sec: sgld invariant proof}

\begin{lemma}
\label{lem: sgld invariant}
    Let Assumptions~\ref{ass: U} and \ref{ass: sg} hold. For every $h \in (0,h_0]$, \eqref{eq: sgld} has a unique invariant distribution $\bar\pi_h$, with $\bar\pi_h(|x|^{12}) \Le C$. For every $f$ satisfying Assumption~\ref{ass: f} and every given $Z_0 \in \mathbb R^d$, 
    \begin{equation*}
        \lim_{N \to \infty} \frac1N \sum_{k=0}^{N-1} f(Z_k) = \bar\pi_h(f) \qquad \text{almost surely} .
    \end{equation*}
\end{lemma}

\begin{proof}
    The proof follows that of Lemma~\ref{lem: invariant} in Appendix~\ref{sec: invariant proof}, and we indicate the changes only. The step of \eqref{eq: sgld} from $Z_k$ is $\sqrt{2h}\, \chi_{k+1} - h\, \nabla \hat U(Z_k;\xi_k)$, with the stochastic gradient evaluated at $Z_k$ itself. The proof of Lemma~\ref{lem: moments} for \eqref{eq: sglm markov} therefore applies with $Z_k$ and $\chi_{k+1}$ in place of $Y_k$ and $\chi_k$, and with the Gaussian shift $\sqrt{h/2}\, \chi_k$ replaced by $0$; in \eqref{eq: drift condition} the drift bound in \eqref{eq: dissipative bounds} applies at $Z_k$ directly. For $q \Le 6$ that proof gives \eqref{eq: lyapunov} and the bound on $\mathbb E |Y_k|^{12}$ in \eqref{eq: moment bounds} for $Z_k$,
    \begin{equation*}
        \hat{\mathbb E}_k |Z_{k+1}|^{2q} \Le \Bigl( 1 - \frac{qm}{16} h \Bigr) |Z_k|^{2q} + C h , \qquad \mathbb E |Z_k|^{12} \Le C \bigl( 1 + |Z_0| \bigr)^{12} ,
    \end{equation*}
    and, as in Appendix~\ref{sec: invariant proof}, the drift condition of Harris's ergodic theorem with $W(x) = |x|^2$ and a level $R_W$ depending only on $m$, $M$, $R$ and $d$.

    The minorization is simpler than in Appendix~\ref{sec: invariant proof} and does not use the Hessian bound of Assumption~\ref{ass: sg}, since the step from $Z_0$ is Gaussian given $\xi_0$,
    \begin{equation*}
        Z_1 \sim \mathcal N\Bigl( Z_0 - h\, \nabla \hat U(Z_0;\xi_0) ,\, 2h I_d \Bigr) .
    \end{equation*}
    Let $h_0$ also satisfy $4 M h_0 \Le 1$, let $|Z_0|^2 \Le R_W$, and let
    \begin{equation*}
        E_0 := \Bigl\{ \xi : |\nabla \hat U(Z_0;\xi)| \Le 4M \bigl( 1 + |Z_0| \bigr) \Bigr\} ,
    \end{equation*}
    so that Markov's inequality and the moment bound of Assumption~\ref{ass: sg} give $\mathbb P_\xi(E_0^c) \Le 4^{-12} \Le \frac12$. For $\xi_0 \in E_0$ and $x$ in the closed unit ball $\mathbb B$, since $4Mh \Le 1$,
    \begin{equation*}
        \bigl| x - Z_0 + h\, \nabla \hat U(Z_0;\xi_0) \bigr| \Le 2 \bigl( 1 + R_W^{\frac12} \bigr) ,
    \end{equation*}
    and the Gaussian density of $Z_1$ given $\xi_0$, with $\nu$ the uniform distribution on $\mathbb B$ and $\mathbb P_\xi(E_0) \Ge \frac12$, gives for every Borel set $A$
    \begin{equation*}
        \mathcal P_h(Z_0, A) \Ge \frac12\, (4\pi h)^{-\frac d2} \exp\Bigl( -\frac{(1 + R_W^{\frac12})^2}{h} \Bigr) \mathrm{Leb}(\mathbb B)\, \nu(A) , \qquad |Z_0|^2 \Le R_W ,
    \end{equation*}
    which is the minorization, with $\alpha_h$ the minimum of $\frac12$ and the factor of $\nu(A)$.

    Harris's ergodic theorem then gives the unique invariant distribution $\bar\pi_h$ and \eqref{eq: harris} with $\bar\pi_h$ in place of $\mu_h$. The moment bound follows from the truncations $\min\bigl\{ |x|^{12}, \ell \bigr\}$, and the almost sure limit from the Harris recurrence and the strong law of large numbers \citep[Theorem~17.0.1(i)]{meyn1993markov}, which use only the drift condition and the minorization, applied to $g = f$, which is Lipschitz by Assumption~\ref{ass: f}, both as in Appendix~\ref{sec: invariant proof}. The passage from \eqref{eq: sglm markov} to the samples of \eqref{eq: sglm} is not needed, since $(Z_k)_{k \Ge 0}$ is itself the chain.
\end{proof}

\subsection{Proof of Theorem~\ref{thm: sgld bias}: numerical bias of SGLD}
\label{sec: sgld bias proof}

The proof uses the local weak error of \eqref{eq: sgld} against \eqref{eq: langevin}, expanded to the order $h^2$ with explicit coefficients. The following lemma is the counterpart of Lemma~\ref{lem: sg weak} for \eqref{eq: sglm markov}, with \eqref{eq: langevin} in place of the modified dynamics \eqref{eq: modified langevin}, and its proof follows the same steps.

\begin{lemma}
\label{lem: sgld weak}
    Let Assumptions~\ref{ass: U} and \ref{ass: sg} hold and let $G \in C^4(\mathbb R^d)$ have bounded derivatives of orders $1$ to $4$. Then, for every $h \in (0,1]$, the chain \eqref{eq: sgld} satisfies, for every $k \Ge 0$,
    \begin{equation}
    \begin{aligned}
        & \biggl| \hat{\mathbb E}_k\bigl[ G(Z_{k+1}) \bigr] - (e^{h\mathcal L} G)(Z_k) - \frac{h^2}{2}\, \nabla^2 G(Z_k) : \bigl( 2\nabla^2 U(Z_k) + \mathcal C(Z_k) \bigr) \\
        & + \frac{h^2}{2}\, \nabla G(Z_k) \cdot \bigl( \nabla^2 U(Z_k)\, \nabla U(Z_k) - \nabla \Delta U(Z_k) \bigr) \biggr| \Le C h^3 \sum_{j=1}^{4} \norm{\nabla^j G}_\infty \bigl( 1 + |Z_k| \bigr)^{6} .
    \end{aligned}
        \label{eq: sgld weak}
    \end{equation}
\end{lemma}

\begin{proof}
    By time homogeneity, it suffices to take $k = 0$. The proof expands $G(Z_1)$ about $Z_0 + D_0$ to the order $2$ in $h\eta_0$, replaces $\nabla^2 G(Z_0 + D_0)$ by $\nabla^2 G(Z_0)$ in the term of order $h^2$, and compares $G(Z_0 + D_0)$ with $G(X_h)$ through the moments of $D_0$ and $X_h - Z_0$. Let $\hat{\mathcal H}_0 := \sigma\bigl( \hat{\mathcal F}_0, (B_s)_{s \Le h} \bigr)$. Then $Z_0$, the increment $D_0$ of \eqref{eq: sgld step} and the solution $X_h$ are $\hat{\mathcal H}_0$-measurable, and $\xi_0$ is independent of $\hat{\mathcal H}_0$, so Assumption~\ref{ass: sg} and \eqref{eq: covariance} give
    \begin{equation*}
        \mathbb E\bigl[ \eta_0 \,\big|\, \hat{\mathcal H}_0 \bigr] = 0 , \qquad \mathbb E\bigl[ \eta_0 \otimes \eta_0 \,\big|\, \hat{\mathcal H}_0 \bigr] = \mathcal C(Z_0) .
    \end{equation*}
    Taylor's formula for $G$ at $Z_0 + D_0$, to the order $2$ in $h\eta_0$, as in the proof of Lemma~\ref{lem: sg weak}, and the bound $\hat{\mathbb E}_0 |\eta_0|^3 \Le C \bigl( 1 + |Z_0| \bigr)^3$ from Assumption~\ref{ass: sg} give
    \begin{equation}
        \Bigl| \hat{\mathbb E}_0\bigl[ G(Z_1) - G(Z_0 + D_0) \bigr] - \frac{h^2}{2}\, \hat{\mathbb E}_0\bigl[ \nabla^2 G(Z_0 + D_0) \bigr] : \mathcal C(Z_0) \Bigr| \Le C h^3 \norm{\nabla^3 G}_\infty \bigl( 1 + |Z_0| \bigr)^3 .
        \label{eq: sgld noise}
    \end{equation}
    Next, $\nabla^2 G(Z_0 + D_0)$ is replaced by $\nabla^2 G(Z_0)$ in \eqref{eq: sgld noise} at the cost $C h^3$. By Taylor's formula,
    \begin{equation*}
        \bigl\| \nabla^2 G(Z_0 + D_0) - \nabla^2 G(Z_0) - \nabla^3 G(Z_0)[D_0] \bigr\| \Le \frac12 \norm{\nabla^4 G}_\infty |D_0|^2 ,
    \end{equation*}
    and, since $\nabla^3 G(Z_0)$ and $\mathcal C(Z_0)$ are $\hat{\mathcal F}_0$-measurable while $B_h$ has mean $0$ and is independent of $\hat{\mathcal F}_0$,
    \begin{equation*}
        \hat{\mathbb E}_0\bigl[ \nabla^3 G(Z_0)[D_0] : \mathcal C(Z_0) \bigr] = -h\, \nabla^3 G(Z_0)[\nabla U(Z_0)] : \mathcal C(Z_0) .
    \end{equation*}
    With $|\nabla U(Z_0)| \Le M \bigl( 1 + |Z_0| \bigr)$ and $\hat{\mathbb E}_0 |D_0|^2 \Le C h \bigl( 1 + |Z_0| \bigr)^2$, \eqref{eq: trace bound} and $\Sigma \Le M^2$ give
    \begin{align}
        & \frac{h^2}{2} \Bigl| \hat{\mathbb E}_0\bigl[ \nabla^2 G(Z_0 + D_0) - \nabla^2 G(Z_0) \bigr] : \mathcal C(Z_0) \Bigr| \Le \frac{\Sigma h^2}{2} \Bigl( h \norm{\nabla^3 G}_\infty |\nabla U(Z_0)| + \frac12 \norm{\nabla^4 G}_\infty\, \hat{\mathbb E}_0 |D_0|^2 \Bigr) \notag \\
        & \quad \Le C h^3 \bigl( \norm{\nabla^3 G}_\infty + \norm{\nabla^4 G}_\infty \bigr) \bigl( 1 + |Z_0| \bigr)^2 .
        \label{eq: sgld hessian replacement}
    \end{align}
        

    The state $Z_0 + D_0$ is compared with $X_h$ as the state $Y_1$ after one step of \eqref{eq: lm markov} is compared with $\tilde Y_0(h)$ in the proof of Lemma~\ref{lem: weak}. Since $X_h = Z_0 + \sqrt2\, B_h - \int_0^h \nabla U(X_s)\, \D s$, the local error is
    \begin{equation*}
        Z_0 + D_0 - X_h = \int_0^h \bigl( \nabla U(X_s) - \nabla U(Z_0) \bigr) \D s ,
    \end{equation*}
    and Assumption~\ref{ass: U} with the bound on $X_s - Z_0$ in \eqref{eq: exact sizes} for \eqref{eq: langevin} give, with $\pnorm{\cdot}_4$ taken under $\hat{\mathbb E}_0$,
    \begin{equation}
        \pnorm{Z_0 + D_0 - X_h}_4 \Le C h^{\frac32} \bigl( 1 + |Z_0| \bigr) , \qquad \pnorm{D_0}_4 + \pnorm{X_h - Z_0}_4 \Le C \sqrt h \bigl( 1 + |Z_0| \bigr) .
        \label{eq: sgld sizes}
    \end{equation}
    With $Z_0$, $Z_0 + D_0$, $X_h$ and $X_h - Z_0$ in place of $Y_0$, $Y_1$, $\tilde Y_0(h)$ and $\tilde W_h$, the expansion \eqref{eq: G expansion full} reads
    \begin{equation}
        G(Z_0 + D_0) - G(X_h) = \sum_{j=1}^{3} \frac{1}{j!}\, \nabla^j G(Z_0) : \Bigl( D_0^{\otimes j} - (X_h - Z_0)^{\otimes j} \Bigr) + \varrho_4 ,
        \label{eq: sgld G expansion}
    \end{equation}
    with a remainder $\varrho_4$ that satisfies
    \begin{equation}
        |\varrho_4| \Le \frac16\, \norm{\nabla^4 G}_\infty\, \bigl| Z_0 + D_0 - X_h \bigr|\, \max\bigl\{ |D_0| ,\, |X_h - Z_0| \bigr\}^3 .
        \label{eq: sgld remainder}
    \end{equation}
    Taking $\hat{\mathbb E}_0$ in \eqref{eq: sgld G expansion}, H\"older's inequality with \eqref{eq: sgld remainder} and \eqref{eq: sgld sizes} gives $\hat{\mathbb E}_0 |\varrho_4| \Le C h^3 \norm{\nabla^4 G}_\infty \bigl( 1 + |Z_0| \bigr)^4$, so
    \begin{equation}
        \biggl| \hat{\mathbb E}_0\bigl[ G(Z_0 + D_0) - G(X_h) \bigr] - \sum_{j=1}^{3} \frac{1}{j!}\, \nabla^j G(Z_0) : \hat{\mathbb E}_0\Bigl[ D_0^{\otimes j} - (X_h - Z_0)^{\otimes j} \Bigr] \biggr| \Le C h^3 \norm{\nabla^4 G}_\infty \bigl( 1 + |Z_0| \bigr)^4 .
        \label{eq: sgld em expansion}
    \end{equation}

    The moments of $D_0$ in \eqref{eq: sgld em expansion} are explicit. With all derivatives of $U$ evaluated at $Z_0$ and $1 \Le i,j,l \Le d$, since $\sqrt2\, B_h \sim \mathcal N(0, 2h I_d)$ and $\nabla U(Z_0)$ is fixed,
    \begin{equation*}
    \begin{gathered}
        \hat{\mathbb E}_0[D_0] = -h\, \nabla U , \qquad \hat{\mathbb E}_0\bigl[ D_0^{\otimes 2} \bigr] = 2h\, I_d + h^2\, \nabla U^{\otimes 2} , \\
        \Bigl( \hat{\mathbb E}_0\bigl[ D_0^{\otimes 3} \bigr] \Bigr)_{ijl} = -2h^2 \bigl( \delta_{ij}\, \partial_l U + \delta_{il}\, \partial_j U + \delta_{jl}\, \partial_i U \bigr) - h^3\, \partial_i U\, \partial_j U\, \partial_l U .
    \end{gathered}
    \end{equation*}
    The moments of $X_h - Z_0$ are expanded in \eqref{eq: exact m1}--\eqref{eq: exact m3} of Lemma~\ref{lem: exact moments} for \eqref{eq: langevin}, with $U$, $Z_0$, $X_h$ and $\hat{\mathbb E}_0$ in place of $\tilde U_h$, $Y_0$, $\tilde Y_0(h)$ and $\mathbb E_0$. Next we subtract them from the moments of $D_0$, for $j = 1, 2, 3$.
    \begin{enumerate}
    \item For $j = 1$, the terms $-h\, \nabla U$ cancel, and the term of order $h^2$ in \eqref{eq: exact m1} remains:
    \begin{equation}
        \Bigl| \hat{\mathbb E}_0\bigl[ D_0 - (X_h - Z_0) \bigr] + \frac{h^2}{2} \bigl( \nabla^2 U\, \nabla U - \nabla \Delta U \bigr) \Bigr| \Le C h^3 \bigl( 1 + |Z_0| \bigr)^3 .
        \label{eq: sgld m1}
    \end{equation}
    \item For $j = 2$, the terms $2h\, I_d$ and $h^2\, \nabla U^{\otimes 2}$ cancel, and the term $-2h^2\, \nabla^2 U$ in \eqref{eq: exact m2} remains:
    \begin{equation}
        \Bigl\| \hat{\mathbb E}_0\bigl[ D_0^{\otimes 2} - (X_h - Z_0)^{\otimes 2} \bigr] - 2h^2\, \nabla^2 U \Bigr\| \Le C h^3 \bigl( 1 + |Z_0| \bigr)^4 .
        \label{eq: sgld m2}
    \end{equation}
    \item For $j = 3$, the terms of order $h^2$ cancel, and $h^3 |\nabla U(Z_0)|^3 \Le C h^3 \bigl( 1 + |Z_0| \bigr)^3$ bounds the term $-h^3\, \partial_i U\, \partial_j U\, \partial_l U$ of $\hat{\mathbb E}_0\bigl[ D_0^{\otimes 3} \bigr]$:
    \begin{equation}
        \Bigl\| \hat{\mathbb E}_0\bigl[ D_0^{\otimes 3} - (X_h - Z_0)^{\otimes 3} \bigr] \Bigr\| \Le C h^3 \bigl( 1 + |Z_0| \bigr)^6 .
        \label{eq: sgld m3}
    \end{equation}
    \end{enumerate}
    Inserting \eqref{eq: sgld m1}--\eqref{eq: sgld m3} in \eqref{eq: sgld em expansion} gives
    \begin{align}
        & \Bigl| \hat{\mathbb E}_0\bigl[ G(Z_0 + D_0) - G(X_h) \bigr] - h^2\, \nabla^2 G(Z_0) : \nabla^2 U(Z_0) + \frac{h^2}{2}\, \nabla G(Z_0) \cdot \bigl( \nabla^2 U(Z_0)\, \nabla U(Z_0) - \nabla \Delta U(Z_0) \bigr) \Bigr| \notag \\
        & \quad \Le C h^3 \sum_{j=1}^{4} \norm{\nabla^j G}_\infty \bigl( 1 + |Z_0| \bigr)^6 .
        \label{eq: sgld em weak}
    \end{align}
    Since $(e^{h\mathcal L} G)(Z_0) = \hat{\mathbb E}_0[G(X_h)]$, the sum of \eqref{eq: sgld noise}, \eqref{eq: sgld hessian replacement} and \eqref{eq: sgld em weak} proves \eqref{eq: sgld weak}.
\end{proof}

\begin{proof}[Proof of Theorem~\ref{thm: sgld bias}]
    Here $\phi_h$ is the solution \eqref{eq: discrete poisson solution} of the discrete Poisson equation \eqref{eq: discrete poisson} of \eqref{eq: langevin}, with $\mathcal L$, $\pi$ and $f$, and not the modified solution $\tilde\phi_h$ of \eqref{eq: modified poisson} used for \eqref{eq: sglm markov}. Let $h_0 \Le \lambda^{-1}$, with $\lambda$ the rate of Theorem~\ref{thm: u decay}, so that \eqref{eq: poisson bound} holds. Let $Z_0 \sim \bar\pi_h$, independent of $(B_t)_{t \Ge 0}$ and of $(\xi_k)_{k \Ge 0}$, so that $Z_1 \sim \bar\pi_h$, and add $Z_0$ to the generators of $\hat{\mathcal F}_k$ in \eqref{eq: sg filtration}, as in the proof of Theorem~\ref{thm: numerical bias}, so that $\hat{\mathcal F}_0 = \sigma(Z_0)$. The discrete Poisson equation \eqref{eq: discrete poisson} at $Z_0$, with $(e^{h\mathcal L} \phi_h)(Z_0) = \hat{\mathbb E}_0[\phi_h(X_h)]$, gives after taking expectations
    \begin{equation}
        \bar\pi_h(f) - \pi(f) = \frac1h\, \mathbb E\bigl[ \phi_h(Z_0) - \phi_h(X_h) \bigr] = \frac1h\, \mathbb E\bigl[ \phi_h(Z_1) - \phi_h(X_h) \bigr] ,
        \label{eq: sgld bias identity}
    \end{equation}
    where $\mathbb E\, \phi_h(Z_0) = \mathbb E\, \phi_h(Z_1)$ since $Z_1 \sim \bar\pi_h$. The expectations are finite, as in the proof of Theorem~\ref{thm: numerical bias}, since $\phi_h$ is Lipschitz with constant $C L$ by \eqref{eq: poisson bound}, $\bar\pi_h(|x|^{12}) \Le C$ by Lemma~\ref{lem: sgld invariant}, and $\mathbb E|X_h - Z_0| \Le C$ by the bound on $X_s - Z_0$ in \eqref{eq: exact sizes} for \eqref{eq: langevin}, averaged over $Z_0 \sim \bar\pi_h$. By \eqref{eq: poisson bound}, $G = \phi_h$ has derivatives of orders $1$ to $4$ bounded by $C L$, so \eqref{eq: sgld weak} of Lemma~\ref{lem: sgld weak} at $k = 0$, averaged over $Z_0 \sim \bar\pi_h$, gives
    \begin{equation*}
    \begin{aligned}
        & \biggl| \mathbb E\bigl[ \phi_h(Z_1) - \phi_h(X_h) \bigr] - \frac{h^2}{2}\, \mathbb E\Bigl[ \nabla^2 \phi_h(Z_0) : \bigl( 2\nabla^2 U(Z_0) + \mathcal C(Z_0) \bigr) \Bigr] \\
        & + \frac{h^2}{2}\, \mathbb E\Bigl[ \nabla \phi_h(Z_0) \cdot \bigl( \nabla^2 U(Z_0)\, \nabla U(Z_0) - \nabla \Delta U(Z_0) \bigr) \Bigr] \biggr| \Le C L h^3\, \mathbb E \bigl( 1 + |Z_0| \bigr)^6 \Le C L h^3 .
    \end{aligned}
    \end{equation*}
    Dividing by $h$ and inserting \eqref{eq: sgld bias identity} gives \eqref{eq: sgld bias}. Finally, \eqref{eq: trace bound}, $\Sigma \Le M^2$, \eqref{eq: poisson bound} and Assumption~\ref{ass: U} give
    \begin{equation*}
    \begin{aligned}
        & \biggl| \mathbb E\Bigl[ \nabla^2 \phi_h(Z_0) : \bigl( 2\nabla^2 U(Z_0) + \mathcal C(Z_0) \bigr) \Bigr] \biggr| \\
        & + \biggl| \mathbb E\Bigl[ \nabla \phi_h(Z_0) \cdot \bigl( \nabla^2 U(Z_0)\, \nabla U(Z_0) - \nabla \Delta U(Z_0) \bigr) \Bigr] \biggr| \Le C L\, \mathbb E \bigl( 1 + |Z_0| \bigr) \Le C L ,
    \end{aligned}
    \end{equation*}
    which with \eqref{eq: sgld bias} gives $|\bar\pi_h(f) - \pi(f)| \Le C L h$.
\end{proof}

\subsection{Proof of Corollary~\ref{cor: sgld coefficient}: coefficient of the numerical bias of SGLD}
\label{sec: sgld coefficient proof}

\begin{proof}[Proof of Corollary~\ref{cor: sgld coefficient}]
    Let
    \begin{equation}
    \begin{aligned}
        A_h(x) & := \nabla^2 \phi_h(x) : \bigl( 2\nabla^2 U(x) + \mathcal C(x) \bigr) - \nabla \phi_h(x) \cdot \bigl( \nabla^2 U(x)\, \nabla U(x) - \nabla \Delta U(x) \bigr) , \\
        A_0(x) & := \nabla^2 \phi_0(x) : \bigl( 2\nabla^2 U(x) + \mathcal C(x) \bigr) - \nabla \phi_0(x) \cdot \bigl( \nabla^2 U(x)\, \nabla U(x) - \nabla \Delta U(x) \bigr) ,
    \end{aligned}
        \label{eq: sgld coefficient integrand}
    \end{equation}
    so that dividing \eqref{eq: sgld bias} by $h$ gives
    \begin{equation*}
        \Bigl| \frac{\bar\pi_h(f) - \pi(f)}{h} - \frac12\, \bar\pi_h(A_h) \Bigr| \Le C L h .
    \end{equation*}
    It suffices to prove $\bar\pi_h(A_h) \to \pi(A_0)$ and to compute $\pi(A_0)$.

    The measures $\bar\pi_h$ converge weakly to $\pi$ as $h \to 0$. For $\varphi \in C_c^\infty(\mathbb R^d)$, which satisfies Assumption~\ref{ass: f} with some constant $L_\varphi$, Theorem~\ref{thm: sgld bias} gives $|\bar\pi_h(\varphi) - \pi(\varphi)| \Le C L_\varphi h$, and Lemma~\ref{lem: sgld invariant} gives, uniformly in $h$,
    \begin{equation}
        \bar\pi_h\Bigl( \bigl( 1 + |x| \bigr)\, \mathbf 1_{\{|x| > \rho\}} \Bigr) \Le \rho^{-11}\, \bar\pi_h\Bigl( \bigl( 1 + |x| \bigr)^{12} \Bigr) \Le C \rho^{-11} , \qquad \rho > 0 .
        \label{eq: sgld tail}
    \end{equation}
    The approximation of a bounded continuous $b$ by $\varphi \in C_c^\infty(\mathbb R^d)$ in Appendix~\ref{sec: bias coefficient proof}, with $\bar\pi_h$ and \eqref{eq: sgld tail} in place of $\mu_h$ and \eqref{eq: mu tail}, then gives $\bar\pi_h(b) \to \pi(b)$ for every bounded continuous $b$.

    The gradient and the Hessian of $\phi_h$ converge to those of $\phi_0$ uniformly on compact sets. The argument for the convergence of $\nabla^2 \tilde\phi_h$ in Appendix~\ref{sec: bias coefficient proof}, through \eqref{eq: modified dynkin} and \eqref{eq: g consistency}, applies with $U$, $\mathcal L$, $f$ and $\pi(f)$ in place of $\tilde U_h$, $\tilde{\mathcal L}_h$, $\tilde f_h$ and $\tilde\pi_h(\tilde f_h)$, with $\phi_h$, the solution $X_s$ of \eqref{eq: langevin} and its initial value $X_0$ in place of $\tilde\phi_h$, $\tilde Y_0(s)$ and $Y_0$, and with \eqref{eq: discrete poisson}, \eqref{eq: poisson bound} and Assumption~\ref{ass: U} in place of \eqref{eq: modified poisson}, \eqref{eq: modified poisson bound} and Lemma~\ref{lem: modified class}; the differences $\nabla \tilde U_h - \nabla U$, $\tilde f_h - f$ and $\tilde\pi_h(\tilde f_h) - \pi(f)$ estimated there are then $0$. It gives, for every sequence $h_n \to 0$, a subsequence along which the derivatives of $\phi_{h_n}$ of orders $1$ to $3$ converge uniformly on compact sets to those of a function that differs from $\phi_0$ by a constant, hence
    \begin{equation*}
        \sup_{|x| \Le \rho} \Bigl( |\nabla \phi_h(x) - \nabla \phi_0(x)| + \norm{\nabla^2 \phi_h(x) - \nabla^2 \phi_0(x)} \Bigr) \to 0 , \qquad h \to 0 , \quad \rho > 0 .
    \end{equation*}
    By \eqref{eq: sgld coefficient integrand}, \eqref{eq: trace bound}, $\Sigma \Le M^2$ and Assumption~\ref{ass: U},
    \begin{equation*}
        |A_h(x) - A_0(x)| \Le C \norm{\nabla^2 \phi_h(x) - \nabla^2 \phi_0(x)} + C \bigl( 1 + |x| \bigr) |\nabla \phi_h(x) - \nabla \phi_0(x)| ,
    \end{equation*}
    so $A_h \to A_0$ uniformly on compact sets, and \eqref{eq: poisson bound} and \eqref{eq: continuous poisson bound} give
    \begin{equation*}
        |A_h(x)| + |A_0(x)| \Le C L \bigl( 1 + |x| \bigr) , \qquad x \in \mathbb R^d .
    \end{equation*}
    The function $A_0$ is continuous, since $\nabla \phi_0$, $\nabla^2 \phi_0$ and the derivatives of $U$ are continuous and $\mathcal C(x)$ is Lipschitz by \eqref{eq: covariance lipschitz}. For every $\rho > 0$,
    \begin{equation*}
        \bigl| \bar\pi_h(A_h) - \pi(A_0) \bigr| \Le \sup_{|x| \Le \rho} |A_h(x) - A_0(x)| + C L\, \bar\pi_h\Bigl( \bigl( 1 + |x| \bigr)\, \mathbf 1_{\{|x| > \rho\}} \Bigr) + \bigl| \bar\pi_h(A_0) - \pi(A_0) \bigr| .
    \end{equation*}
    The weak convergence of $\bar\pi_h$ to $\pi$, with the uniform integrability of $1 + |x|$ under $\bar\pi_h$ that \eqref{eq: sgld tail} provides, gives $\bar\pi_h(A_0) \to \pi(A_0)$, and \eqref{eq: sgld tail} bounds the term carrying $\mathbf 1_{\{|x| > \rho\}}$, so
    \begin{equation*}
        \limsup_{h \to 0}\, \bigl| \bar\pi_h(A_h) - \pi(A_0) \bigr| \Le C L \rho^{-11} , \qquad \rho > 0 ,
    \end{equation*}
    and letting $\rho \to \infty$ gives $\bar\pi_h(A_h) \to \pi(A_0)$.

    It remains to compute $\pi(A_0)$ for $A_0$ of \eqref{eq: sgld coefficient integrand}. Since $\phi_0$ has bounded derivatives of orders $1$ to $3$ by \eqref{eq: continuous poisson bound}, \eqref{eq: sgld integration by parts} with $\psi = \phi_0$ gives
    \begin{equation*}
        \pi(A_0) = \pi\Bigl( \nabla^2 \phi_0(x) : \bigl( 2\nabla^2 U(x) + \mathcal C(x) \bigr) \Bigr) - \pi\bigl( \nabla^2 \phi_0(x) : \nabla^2 U(x) \bigr) = \pi\Bigl( \nabla^2 \phi_0(x) : \bigl( \nabla^2 U(x) + \mathcal C(x) \bigr) \Bigr) ,
    \end{equation*}
    which with $\bar\pi_h(A_h) \to \pi(A_0)$ proves \eqref{eq: sgld coefficient}.
\end{proof}


\section{Proofs for the explicit coefficient in Wasserstein-1 error}
\label{app: w1}

Throughout this appendix, the constant $C$ depends only on $m$, $M$, $R$ and $d$, and
\begin{equation}
    \mathcal D := \bigl\{ \omega \in C^4(\mathbb R^d) : \norm{\nabla^j \omega}_\infty < \infty ,\ j = 1, 2, 3, 4 \bigr\} , \qquad L^2_0(\pi) := \bigl\{ g \in L^2(\pi) : \pi(g) = 0 \bigr\} ,
    \label{eq: test class}
\end{equation}
so $\mathcal D$ is the set of the test functions $\omega$ of \eqref{eq: riesz}, and $L^2_0(\pi)$ is a Hilbert space with the inner product of $L^2(\pi)$.

\subsection{Proof of Theorem~\ref{thm: riesz}: well-posedness of the Riesz representor}
\label{sec: riesz well-posedness}

Theorem~\ref{thm: riesz} constructs the Riesz representor $\psi_{\mathcal A}$ of \eqref{eq: riesz} in $L^2_0(\pi)$ from a bound of $\pi(|\nabla \omega|^2)$ by $\pi((\mathcal L \omega)^2)$ for $\omega \in \mathcal D$, which follows from the Poincar\'e inequality \eqref{eq: poincare} of Lemma~\ref{lem: poincare}. The constant $C_P$ of \eqref{eq: poincare} therefore enters the bound on $\pi(\psi_{\mathcal A}^2)$ in Theorem~\ref{thm: riesz} and the constant of Theorem~\ref{thm: kr expansion}.

By \eqref{eq: dissipative bounds}, $U$ is bounded from below, and $\langle x, \nabla U(x) \rangle \Ge \frac m4 |x|^2 - C \Ge |x|$ for $|x| \Ge \rho$, with $\rho$ depending only on $m$, $M$, $R$ and $d$. Hence $U$ satisfies the drift condition of \citep[Corollary~1.6]{bakry2008simple}, and \citep[Theorem~1.4]{bakry2008simple} gives \eqref{eq: poincare} for smooth $g$. The constant $C_P$ depends only on the Lyapunov function of \citep[Corollary~1.6]{bakry2008simple} and on the Poincar\'e constant of $\pi$ restricted to the ball $\{|x| < \rho\}$, hence only on $m$, $M$, $R$ and $d$. Mollification and dominated convergence, with the moments of Lemma~\ref{lem: dissipative}, extend \eqref{eq: poincare} to every $g \in C^1(\mathbb R^d)$ with a bounded gradient. The next lemma is the identity of Bochner under $\pi$, used in the proofs of Theorems~\ref{thm: riesz} and \ref{thm: batch threshold}.

\begin{lemma}
\label{lem: bochner}
    Let Assumption~\ref{ass: U} hold. For every $\omega \in \mathcal D$,
    \begin{equation}
        \pi\bigl( (\mathcal L \omega)^2 \bigr) = \pi\bigl( \nabla^2 \omega : \nabla^2 \omega \bigr) + \pi\bigl( \nabla \omega \cdot \nabla^2 U\, \nabla \omega \bigr) .
        \label{eq: bochner}
    \end{equation}
\end{lemma}

\begin{proof}
    Differentiating \eqref{eq: generator} gives $\partial_i \mathcal L \omega = \mathcal L \partial_i \omega - \sum_{j=1}^d \partial_i \partial_j U\, \partial_j \omega$ for each component $i$. The functions $\omega$, $\mathcal L \omega$ and $\partial_i \omega$ and their derivatives of orders $1$ and $2$ grow at most linearly, so \eqref{eq: generator ibp} with $g_1 = \mathcal L \omega$, $g_2 = \omega$, and then with $g_1 = g_2 = \partial_i \omega$, gives
    \begin{equation*}
    \begin{aligned}
        \pi\bigl( (\mathcal L \omega)^2 \bigr) & = -\sum_{i=1}^d \pi\bigl( \partial_i \omega\, \partial_i \mathcal L \omega \bigr) = -\sum_{i=1}^d \pi\bigl( \partial_i \omega\, \mathcal L \partial_i \omega \bigr) + \pi\bigl( \nabla \omega \cdot \nabla^2 U\, \nabla \omega \bigr) \\
        & = \sum_{i=1}^d \pi\bigl( |\nabla \partial_i \omega|^2 \bigr) + \pi\bigl( \nabla \omega \cdot \nabla^2 U\, \nabla \omega \bigr) ,
    \end{aligned}
    \end{equation*}
    which is \eqref{eq: bochner}.
\end{proof}

\begin{proof}[Proof of Theorem~\ref{thm: riesz}]
    The proof bounds $|\pi(\nabla^2 \omega : \mathcal A)|$ by the norm of $\mathcal L \omega$ in $L^2(\pi)$, through \eqref{eq: poincare} and \eqref{eq: bochner}, and applies the Riesz representation theorem on the dense set of the functions $-\mathcal L \omega$, $\omega \in \mathcal D$, in $L^2_0(\pi)$. Let $\omega \in \mathcal D$. Since $\omega$ grows at most linearly, \eqref{eq: generator ibp} and \eqref{eq: poincare} give
    \begin{equation*}
        \pi\bigl( |\nabla \omega|^2 \bigr) = \pi\bigl( (\omega - \pi(\omega)) (-\mathcal L \omega) \bigr) \Le C_P^{\frac12}\, \pi\bigl( |\nabla \omega|^2 \bigr)^{\frac12}\, \pi\bigl( (\mathcal L \omega)^2 \bigr)^{\frac12} ,
    \end{equation*}
    so $\pi(|\nabla \omega|^2) \Le C_P\, \pi((\mathcal L \omega)^2)$, and \eqref{eq: bochner} with $\nabla^2 U \succeq -M I_d$ gives
    \begin{equation}
        \pi\bigl( \nabla^2 \omega : \nabla^2 \omega \bigr) \Le \pi\bigl( (\mathcal L \omega)^2 \bigr) + M\, \pi\bigl( |\nabla \omega|^2 \bigr) \Le (1 + M C_P)\, \pi\bigl( (\mathcal L \omega)^2 \bigr) .
        \label{eq: hessian l2}
    \end{equation}
    For every symmetric matrix $B$, the Cauchy--Schwarz inequality for the inner product $B : \mathcal A(x)$ and $\mathcal A(x) : \mathcal A(x) \Le d\, \norm{\mathcal A(x)}^2$ give $|B : \mathcal A(x)| \Le \sqrt d\, \norm{\mathcal A(x)}\, (B : B)^{\frac12}$. With $B = \nabla^2 \omega(x)$, this is the pointwise bound
    \begin{equation*}
        \bigl| \nabla^2 \omega(x) : \mathcal A(x) \bigr| \Le \sqrt d\, \norm{\mathcal A}_\infty \bigl( \nabla^2 \omega(x) : \nabla^2 \omega(x) \bigr)^{\frac12} , \qquad x \in \mathbb R^d .
    \end{equation*}
    Integrating it against $\pi$, the Cauchy--Schwarz inequality under $\pi$ gives
    \begin{equation*}
        \bigl| \pi\bigl( \nabla^2 \omega : \mathcal A \bigr) \bigr| \Le \sqrt d\, \norm{\mathcal A}_\infty\, \pi\Bigl( \bigl( \nabla^2 \omega : \nabla^2 \omega \bigr)^{\frac12} \Bigr) \Le \sqrt d\, \norm{\mathcal A}_\infty\, \pi\bigl( \nabla^2 \omega : \nabla^2 \omega \bigr)^{\frac12} ,
    \end{equation*}
    and inserting the bound on $\pi(\nabla^2 \omega : \nabla^2 \omega)$ in \eqref{eq: hessian l2} gives
    \begin{equation}
        \bigl| \pi\bigl( \nabla^2 \omega : \mathcal A \bigr) \bigr| \Le \sqrt{d\, (1 + M C_P)}\; \norm{\mathcal A}_\infty\; \pi\bigl( (\mathcal L \omega)^2 \bigr)^{\frac12} , \qquad \omega \in \mathcal D .
        \label{eq: riesz functional bound}
    \end{equation}

    The functions $-\mathcal L \omega$, $\omega \in \mathcal D$, which lie in $L^2_0(\pi)$ by their linear growth and \eqref{eq: generator ibp} with $g_1 = 1$, are dense in $L^2_0(\pi)$. Every $g \in C_c^\infty(\mathbb R^d)$ satisfies Assumption~\ref{ass: f}, so the solution $\phi_0$ of the continuous Poisson equation \eqref{eq: continuous poisson} with $g$ in place of $f$ belongs to $\mathcal D$ by \eqref{eq: continuous poisson bound}, and $g - \pi(g) = -\mathcal L \phi_0$. The functions $g - \pi(g)$, $g \in C_c^\infty(\mathbb R^d)$, are dense in $L^2_0(\pi)$. By \eqref{eq: riesz functional bound}, $-\mathcal L \omega \mapsto \pi(\nabla^2 \omega : \mathcal A)$ is a well-defined linear functional on this dense set, bounded in the norm of $L^2_0(\pi)$, and linear in $\mathcal A$. It extends to $L^2_0(\pi)$, and the Riesz representation theorem gives the unique $\psi_{\mathcal A} \in L^2_0(\pi)$ that satisfies \eqref{eq: riesz} for every $\omega \in \mathcal D$, and \eqref{eq: riesz functional bound} gives \eqref{eq: riesz bound}.
\end{proof}


\subsection{Auxiliary estimates for Theorem~\ref{thm: kr expansion}}
\label{sec: kr auxiliary}

Throughout this subsection and the next, $h \Le h_0$ with $h_0 \Le \lambda^{-1}$, $\lambda$ the rate of Lemma~\ref{lem: smoothing}. For a signed measure $\nu$ on $\mathbb R^d$ with $\nu(\mathbb R^d) = 0$ and $\int (1 + |x|)\, |\nu|(\D x) < \infty$, such as $\psi\pi$ with $\psi \in L^2_0(\pi)$, the Kantorovich--Rubinstein norm of Section~\ref{subsec: notation} satisfies
\begin{equation}
    \norm{\nu}_{\mathrm{KR}} = \sup\biggl\{ \int_{\mathbb R^d} f\, \D\nu : f \text{ satisfies Assumption~\ref{ass: f}},\ \norm{\nabla f}_\infty \Le 1 \biggr\} .
    \label{eq: kr test functions}
\end{equation}
The supremum on the right side of \eqref{eq: kr test functions} is at most $\norm{\nu}_{\mathrm{KR}}$, since its functions $f$ are Lipschitz with $\norm{\nabla f}_\infty \Le 1$, and it equals $\norm{\nu}_{\mathrm{KR}}$, since the mollifications of a Lipschitz function $g$ with $\norm{\nabla g}_\infty \Le 1$ satisfy Assumption~\ref{ass: f} with the same bound on the gradient and converge to $g$ uniformly.

In the proofs of Theorems~\ref{thm: kr expansion} and \ref{thm: batch threshold}, the test function $f$ satisfies Assumption~\ref{ass: f} and $\norm{\nabla f}_\infty \Le 1$, and the estimates do not depend on the constant $L$ of Assumption~\ref{ass: f}. The modified test function \eqref{eq: heat shift} is $\tilde f_h = f * \mathcal N(0, \frac h2 I_d)$, where we informally use $\mathcal N(0, \frac h2 I_d)$ also for the density of this Gaussian distribution, so
\begin{equation*}
    \nabla^j \tilde f_h = (\nabla f) * \nabla^{j-1}\mathcal N\Bigl( 0, \frac h2 I_d \Bigr) , \qquad \int_{\mathbb R^d} \Bigl\| \nabla^{j-1}\mathcal N\Bigl( 0, \frac h2 I_d \Bigr)(y) \Bigr\| \,\D y \Le C h^{-\frac{j-1}{2}} , \qquad j = 1, 2, 3, 4 ,
\end{equation*}
and
\begin{equation}
    \norm{\nabla \tilde f_h}_\infty \Le 1 , \qquad
    \norm{\nabla^j \tilde f_h}_\infty \Le C h^{-\frac{j-1}{2}} , \qquad j = 2, 3, 4 .
    \label{eq: heat derivatives}
\end{equation}

The semigroup estimates of Appendix~\ref{sec: semigroup} and Lemma~\ref{lem: fourth from lipschitz} combine into one bound for the exact and the modified dynamics. The proof technique is essentially that of \citep[Lemma~10]{ye2026mse}, which splits the times into $0 < t \Le 1$ and $t > 1$.

\begin{lemma}
\label{lem: smoothing}
    Let Assumption~\ref{ass: U} hold. There is $\lambda > 0$, depending only on $m$, $M$, $R$ and $d$, such that, for every $h \in (0,h_0]$, every $g \in \mathcal D$ of \eqref{eq: test class} and every $t > 0$,
    \begin{equation}
        \norm{\nabla^j (e^{t\mathcal L} g)}_\infty + \norm{\nabla^j (e^{t\tilde{\mathcal L}_h} g)}_\infty \Le C \min(1,t)^{-\frac{j-1}{2}}\, e^{-\lambda t}\, \norm{\nabla g}_\infty , \qquad j = 1, 2, 3, 4 .
        \label{eq: smoothing}
    \end{equation}
\end{lemma}

\begin{proof}
    For $j = 1$, the bound on $e^{t\mathcal L} g$ is \eqref{eq: grad decay}. For $j = 2, 3, 4$ and $0 < t \Le 1$, \eqref{eq: smoothing 23} and \eqref{eq: fourth from lipschitz} give
    \begin{equation*}
        \norm{\nabla^j (e^{t\mathcal L} g)}_\infty \Le C\, t^{-\frac{j-1}{2}}\, \norm{\nabla g}_\infty \Le C e^{\lambda} \min(1,t)^{-\frac{j-1}{2}}\, e^{-\lambda t}\, \norm{\nabla g}_\infty ,
    \end{equation*}
    since $\min(1,t) = t$ and $e^{-\lambda t} \Ge e^{-\lambda}$, which is \eqref{eq: smoothing} for $e^{t\mathcal L}$ with $C e^{\lambda}$ in place of $C$. For $j = 2, 3, 4$ and $t > 1$, the semigroup property gives $e^{t\mathcal L} g = e^{\mathcal L} \bigl( e^{(t-1)\mathcal L} g \bigr)$, and $e^{(t-1)\mathcal L} g \in \mathcal D$, since \eqref{eq: short time} at time $t - 1$ bounds its derivatives of orders $1$ to $4$. Hence \eqref{eq: smoothing 23} and \eqref{eq: fourth from lipschitz} at time $1$, with $e^{(t-1)\mathcal L} g$ in place of $g$, bound $\norm{\nabla^j (e^{t\mathcal L} g)}_\infty$ by $C \norm{\nabla (e^{(t-1)\mathcal L} g)}_\infty$, and \eqref{eq: grad decay} at time $t - 1$ bounds this gradient, so that
    \begin{equation*}
        \norm{\nabla^j (e^{t\mathcal L} g)}_\infty \Le C \bigl\| \nabla (e^{(t-1)\mathcal L} g) \bigr\|_\infty \Le C e^{-\lambda (t-1)} \norm{\nabla g}_\infty = C e^{\lambda} e^{-\lambda t} \norm{\nabla g}_\infty .
    \end{equation*}
    Since $\min(1,t) = 1$ for $t > 1$, this is \eqref{eq: smoothing} for $e^{t\mathcal L}$ with $C e^{\lambda}$ in place of $C$.

    The same argument applies to $e^{t\tilde{\mathcal L}_h}$, since $\tilde U_h$ satisfies the conditions of Assumption~\ref{ass: U} that this argument uses uniformly in $h$. Lemmas~\ref{lem: semigroup}, \ref{lem: short time} and \ref{lem: fourth from lipschitz} use only the bounds of Assumption~\ref{ass: U} up to the order $5$ and the convexity outside the ball of radius $R$, and by Lemma~\ref{lem: modified class}, for every $h \in (0, h_0]$, $\tilde U_h$ satisfies them with $2M$ and $\frac m2$ in place of $M$ and $m$. These lemmas therefore hold for $e^{t\tilde{\mathcal L}_h}$ with a rate and a constant that depend only on $m$, $M$, $R$ and $d$, as in Corollary~\ref{cor: modified poisson}. Since the left side of \eqref{eq: smoothing} is the sum of the norms for $e^{t\mathcal L}$ and for $e^{t\tilde{\mathcal L}_h}$, \eqref{eq: smoothing} holds with $\lambda$ the smaller of the rates for $U$ and for $\tilde U_h$, and $C$ twice the larger of the constants.
\end{proof}

In the proofs of Theorem~\ref{thm: kr expansion}, the bounds \eqref{eq: smoothing} at the times $t = ih$ are summed over $i \Ge 0$ through the following elementary bounds.



\begin{lemma}
\label{lem: time sums}
    For every $h \in (0,h_0]$,
    \begin{equation}
        h \sum_{i=0}^{\infty} \min\bigl( 1, (i+1)h \bigr)^{-\frac{j-1}{2}} e^{-\lambda i h} \Le
        \begin{cases}
            C , & j = 1, 2 , \\
            C \bigl( 1 + \log\frac1h \bigr) , & j = 3 , \\
            C h^{-\frac12} , & j = 4 .
        \end{cases}
        \label{eq: time sums}
    \end{equation}
\end{lemma}

\begin{proof}
    Split the sum at $(i+1) h = 1$. The terms with $(i+1)h > 1$, that is, $i \Ge \lfloor 1/h \rfloor$, have the minimum $1$, and since $\lambda h \Le 1$,
    \begin{equation*}
        h \sum_{i = \lfloor 1/h \rfloor}^{\infty} e^{-\lambda i h} \Le \frac{h}{1 - e^{-\lambda h}} \Le \frac{2}{\lambda} .
    \end{equation*}
    For the terms with $(i+1)h \Le 1$, put $n = i+1$ and bound $e^{-\lambda ih}$ by $1$:
    \begin{equation*}
        h \sum_{n=1}^{\lfloor 1/h \rfloor} (nh)^{-\frac{j-1}{2}} = h^{\frac{3-j}{2}} \sum_{n=1}^{\lfloor 1/h \rfloor} n^{-\frac{j-1}{2}} \Le
        \begin{cases}
            1 , & j = 1 , \\
            2 , & j = 2 , \\
            1 + \log\frac1h , & j = 3 , \\
            3 h^{-\frac12} , & j = 4 .
        \end{cases}
    \end{equation*}
    Adding the bound $2/\lambda$ gives \eqref{eq: time sums}.
\end{proof}

For a Lipschitz test function, a step of \eqref{eq: sgld} and a step of \eqref{eq: langevin} from the same state $Z_0$ give expectations that differ by at most $C h^{\frac32} (1 + |Z_0|)$, although the gradient noise enters the step with the factor $h$: the Gaussian increment of the step smooths the test function, and the noise has mean $0$.

\begin{lemma}
\label{lem: sgld lipschitz step}
    Let Assumptions~\ref{ass: U} and \ref{ass: sg} hold. Then, for every Lipschitz function $f$ with $\norm{\nabla f}_\infty \Le 1$ and every $h \in (0,1]$, the chain \eqref{eq: sgld} satisfies
    \begin{equation}
        \bigl| \hat{\mathbb E}_0\bigl[ f(Z_1) \bigr] - (e^{h\mathcal L} f)(Z_0) \bigr| \Le C h^{\frac32} \bigl( 1 + |Z_0| \bigr) .
        \label{eq: sgld lipschitz step}
    \end{equation}
\end{lemma}

\begin{proof}
    Write the step from $Z_0$ as in \eqref{eq: sgld step}, with $D_0$, $\eta_0$ and the solution $X_s$ of \eqref{eq: langevin} from $X_0 = Z_0$ as in Appendix~\ref{app: sgld}. Then
    \begin{equation}
        \hat{\mathbb E}_0\bigl[ f(Z_1) \bigr] - (e^{h\mathcal L} f)(Z_0) = \hat{\mathbb E}_0\bigl[ f(Z_0 + D_0 - h\eta_0) - f(Z_0 + D_0) \bigr] + \hat{\mathbb E}_0\bigl[ f(Z_0 + D_0) - f(X_h) \bigr] .
        \label{eq: sgld step split}
    \end{equation}

    Since $\norm{\nabla f}_\infty \Le 1$, the bound on $Z_0 + D_0 - X_h$ in \eqref{eq: sgld sizes} gives
    \begin{equation*}
        \Bigl| \hat{\mathbb E}_0\bigl[ f(Z_0 + D_0) - f(X_h) \bigr] \Bigr| \Le \hat{\mathbb E}_0 |Z_0 + D_0 - X_h| \Le C h^{\frac32} \bigl( 1 + |Z_0| \bigr) .
    \end{equation*}
    For the expectation of $f(Z_0 + D_0 - h\eta_0) - f(Z_0 + D_0)$ in \eqref{eq: sgld step split}, condition on $\hat{\mathcal F}_0$ and $\eta_0$. Since $\sqrt2\, B_h \sim \mathcal N(0, 2h I_d)$ is independent of $\sigma(\hat{\mathcal F}_0, \eta_0)$,
    \begin{equation*}
        \mathbb E\bigl[ f(Z_0 + D_0 - h\eta_0) - f(Z_0 + D_0) \,\big|\, \hat{\mathcal F}_0, \eta_0 \bigr] = (e^{h\Delta} f)\bigl( Z_0 - h\, \nabla U(Z_0) - h\eta_0 \bigr) - (e^{h\Delta} f)\bigl( Z_0 - h\, \nabla U(Z_0) \bigr) ,
    \end{equation*}
    with $e^{h\Delta} f = f * \mathcal N(0, 2h I_d)$, for which $\norm{\nabla^2 e^{h\Delta} f}_\infty \Le C h^{-\frac12}$ as in \eqref{eq: heat derivatives}. Taylor's formula in $h\eta_0$ gives
    \begin{equation*}
        \Bigl| \mathbb E\bigl[ f(Z_0 + D_0 - h\eta_0) - f(Z_0 + D_0) \,\big|\, \hat{\mathcal F}_0, \eta_0 \bigr] + h\, \nabla (e^{h\Delta} f)\bigl( Z_0 - h\, \nabla U(Z_0) \bigr) \cdot \eta_0 \Bigr| \Le C h^{\frac32}\, |\eta_0|^2 .
    \end{equation*}
    Since $\nabla (e^{h\Delta} f)\bigl( Z_0 - h\, \nabla U(Z_0) \bigr)$ is $\hat{\mathcal F}_0$-measurable and $\hat{\mathbb E}_0[\eta_0] = 0$, taking $\hat{\mathbb E}_0$ removes the term linear in $\eta_0$, and $\hat{\mathbb E}_0 |\eta_0|^2 \Le \Sigma \Le M^2$ by Assumption~\ref{ass: sg} gives
    \begin{equation*}
        \Bigl| \hat{\mathbb E}_0\bigl[ f(Z_0 + D_0 - h\eta_0) - f(Z_0 + D_0) \bigr] \Bigr| \Le C h^{\frac32}\, \hat{\mathbb E}_0 |\eta_0|^2 \Le C h^{\frac32} .
    \end{equation*}
    With the bound $C h^{\frac32} (1 + |Z_0|)$ on the expectation of $f(Z_0 + D_0) - f(X_h)$, \eqref{eq: sgld step split} then gives \eqref{eq: sgld lipschitz step}.
\end{proof}

In the proof of \eqref{eq: kr sgld}, the integrand of the sum in \eqref{eq: sgld leading} carries $\nabla^2 U\, \nabla U - \nabla \Delta U$, whose gradient grows linearly, so the integrand satisfies only the bound \eqref{eq: local lipschitz} below, and $\bar\pi_h$ is replaced by $\pi$ in the coefficient of $\frac{h^2}{2}$ in \eqref{eq: sgld leading} through the following truncation.

\begin{lemma}
\label{lem: local lipschitz}
    Let $\mu$ and $\nu$ be probability measures on $\mathbb R^d$ with $\mu(|x|^{12}) + \nu(|x|^{12}) < \infty$ and $\mathcal W_1(\mu, \nu) \Le 1$, and let $F : \mathbb R^d \to \mathbb R$ satisfy, for some $K \Ge 0$,
    \begin{equation}
        |F(x) - F(x')| \Le K \bigl( 1 + |x| + |x'| \bigr) |x - x'| , \qquad x, x' \in \mathbb R^d .
        \label{eq: local lipschitz}
    \end{equation}
    Then $|\mu(F) - \nu(F)| \Le 3 K \bigl( 1 + \mu(|x|^{12}) + \nu(|x|^{12}) \bigr)\, \mathcal W_1(\mu, \nu)^{\frac{10}{11}}$.
\end{lemma}
Without the factor $1 + |x| + |x'|$ in \eqref{eq: local lipschitz}, $F$ is Lipschitz and $|\mu(F) - \nu(F)| \Le K\, \mathcal W_1(\mu, \nu)$ directly; with this factor, the moments of order $12$ of $\mu$ and $\nu$ compensate its growth, at the cost of the reduced exponent $\frac{10}{11}$.

\begin{proof}
    If $\mathcal W_1(\mu,\nu) = 0$, then $\mu = \nu$. Otherwise let $\rho := \mathcal W_1(\mu,\nu)^{-\frac1{11}} \Ge 1$, and let $T_\rho(x) := x$ for $|x| \Le \rho$ and $T_\rho(x) := \rho x / |x|$ for $|x| > \rho$, the retraction onto the closed ball of radius $\rho$, which is $1$-Lipschitz. By \eqref{eq: local lipschitz}, $F \circ T_\rho$ is Lipschitz with constant $K(1 + 2\rho) \Le 3K\rho$, so the Kantorovich--Rubinstein duality gives
    \begin{equation*}
        \bigl| \mu(F \circ T_\rho) - \nu(F \circ T_\rho) \bigr| \Le 3 K \rho\, \mathcal W_1(\mu,\nu) .
    \end{equation*}
    For $|x| > \rho$, $|T_\rho(x)| = \rho < |x|$ and $|x - T_\rho(x)| \Le |x|$, so \eqref{eq: local lipschitz} and $|x| > \rho \Ge 1$ give
    \begin{equation*}
        |F(x) - F(T_\rho(x))| \Le K \bigl( 1 + 2|x| \bigr) |x| \Le 3 K |x|^2 \Le 3 K \rho^{-10} |x|^{12} ,
    \end{equation*}
    while $F = F \circ T_\rho$ on $\{|x| \Le \rho\}$. Hence
    \begin{equation*}
        |\mu(F) - \nu(F)| \Le 3 K \rho\, \mathcal W_1(\mu,\nu) + 3 K \rho^{-10} \bigl( \mu(|x|^{12}) + \nu(|x|^{12}) \bigr) = 3 K \bigl( 1 + \mu(|x|^{12}) + \nu(|x|^{12}) \bigr)\, \mathcal W_1(\mu,\nu)^{\frac{10}{11}} ,
    \end{equation*}
    with $F$ integrable under $\mu$ and $\nu$ since $|F(x)| \Le |F(0)| + K(1 + |x|)|x|$.
\end{proof}

\subsection{Proof of Theorem~\ref{thm: kr expansion}: expansions in the Kantorovich--Rubinstein norm}
\label{sec: kr expansion proof}

Since $\mathcal C(x)$ is positive semidefinite, $\norm{\mathcal C}_\infty \Le \Sigma \Le M^2$ by \eqref{eq: covariance}, and $\norm{\nabla^2 U}_\infty \Le M$. The bound \eqref{eq: riesz bound} of Theorem~\ref{thm: riesz}, the Cauchy--Schwarz inequality and the moments of Lemma~\ref{lem: dissipative} then give
\begin{equation}
    \pi\bigl( |\psi_{\mathcal C}|\, (1 + |x|)^2 \bigr) + \pi\bigl( |\psi_{\nabla^2 U} + \psi_{\mathcal C}|\, (1 + |x|)^2 \bigr) \Le C .
    \label{eq: psi moments}
\end{equation}

\begin{proof}[Proof of \eqref{eq: kr sglm} in Theorem~\ref{thm: kr expansion}]
    Recall that $\mu_h$ is the invariant distribution of \eqref{eq: sglm markov} and $\tilde\pi_h \propto e^{-\tilde U_h}$ is the modified target, the invariant distribution of \eqref{eq: modified langevin}. The proof writes $\mu_h(\tilde f_h) - \tilde\pi_h(\tilde f_h)$, by the telescoping formula \eqref{eq: sglm telescope}, as a series of local errors of Lemma~\ref{lem: sg weak} on the functions
    \begin{equation}
        \tilde{\mathcal V}_i := e^{ih\tilde{\mathcal L}_h} \tilde f_h , \qquad i = 0, 1, 2, \dots ,
        \label{eq: propagated test functions}
    \end{equation}
    moves the leading term of each from the law of $Y_0$ to $\pi$, and sums the series through \eqref{eq: riesz} into $\frac h2\, \pi(f \psi_{\mathcal C})$.
    
    Let $f$ satisfy Assumption~\ref{ass: f} and $\norm{\nabla f}_\infty \Le 1$.
    Let $Y_0 \sim \mu_h$, independent of $(B_t)_{t \Ge 0}$ and $(\xi_k)_{k \Ge 0}$, so that $Y_1 \sim \mu_h$, and add $Y_0$ to the generators of $\hat{\mathcal F}_k$ in \eqref{eq: sg filtration}, as in the proof of Theorem~\ref{thm: numerical bias}, so that $\hat{\mathcal F}_0 = \sigma(Y_0)$, while $\chi_0$ and $\xi_0$ remain independent of $\hat{\mathcal F}_0$; then $Z_0 = Y_0 + \sqrt{h/2}\, \chi_0 \sim \pi_h$. By \eqref{eq: heat shift invariant}, $\pi_h(f) = \mu_h(\tilde f_h)$, so
    \begin{equation}
        \pi_h(f) - \pi(f) = \bigl( \mu_h(\tilde f_h) - \tilde\pi_h(\tilde f_h) \bigr) + \bigl( \tilde\pi_h(\tilde f_h) - \pi(f) \bigr) ,
        \label{eq: sglm split}
    \end{equation}
    and \eqref{eq: shifted bias} bounds the bracket $\tilde\pi_h(\tilde f_h) - \pi(f)$ of \eqref{eq: sglm split} by $C h^2$.

    For $i = 0$, $\tilde{\mathcal V}_0 = \tilde f_h$ and \eqref{eq: heat derivatives} bounds its derivatives, and for $i \Ge 1$, \eqref{eq: smoothing} at $t = ih$ with $\min(1, ih) \Ge \frac12 \min(1, (i+1)h)$ bounds those of $\tilde{\mathcal V}_i$, so that, for $i \Ge 0$ and $j = 1, 2, 3, 4$,
    \begin{equation}
        \norm{\nabla^j \tilde{\mathcal V}_i}_\infty \Le C \min\bigl( 1, (i+1)h \bigr)^{-\frac{j-1}{2}} e^{-\lambda i h} ,
        \label{eq: propagated derivatives}
    \end{equation}
    and $\tilde{\mathcal V}_i \in \mathcal D$. Since $\tilde Y_0(h)$ solves the modified dynamics \eqref{eq: modified langevin} from $Y_0$, driven by Brownian increments independent of $\hat{\mathcal F}_0$, $\hat{\mathbb E}_0[\tilde{\mathcal V}_i(\tilde Y_0(h))] = \tilde{\mathcal V}_{i+1}(Y_0)$, and since $Y_0$ and $Y_1$ have the law $\mu_h$,
    \begin{equation*}
        \mathbb E\bigl[ \tilde{\mathcal V}_i(Y_1) - \tilde{\mathcal V}_i(\tilde Y_0(h)) \bigr] = \mu_h(\tilde{\mathcal V}_i) - \mu_h(\tilde{\mathcal V}_{i+1}) .
    \end{equation*}
    The modified target $\tilde\pi_h$ is invariant under $e^{t\tilde{\mathcal L}_h}$ and has a moment of order $1$ bounded uniformly in $h$, as in the proof of Lemma~\ref{lem: shifted bias}, so \eqref{eq: propagated derivatives} at $j = 1$ gives
    \begin{equation}
        \bigl| \tilde{\mathcal V}_N(y) - \tilde\pi_h(\tilde f_h) \bigr| \Le \norm{\nabla \tilde{\mathcal V}_N}_\infty \int_{\mathbb R^d} |y - y'|\, \tilde\pi_h(\D y') \Le C e^{-\lambda N h} \bigl( 1 + |y| \bigr) ,
        \label{eq: sglm relaxation}
    \end{equation}
    and $\mu_h(\tilde{\mathcal V}_N) \to \tilde\pi_h(\tilde f_h)$ as $N \to \infty$. Summing over $i$,
    \begin{equation}
        \mu_h(\tilde f_h) - \tilde\pi_h(\tilde f_h) = \sum_{i=0}^{\infty} \mathbb E\bigl[ \tilde{\mathcal V}_i(Y_1) - \tilde{\mathcal V}_i(\tilde Y_0(h)) \bigr] .
        \label{eq: sglm telescope}
    \end{equation}

    By \eqref{eq: sg weak} at $k = 0$ with $G = \tilde{\mathcal V}_i$, the tower property for the $\hat{\mathcal F}_0$-measurable factor $\nabla^2 \tilde{\mathcal V}_i(Y_0)$, and the moment bound $\mu_h(|y|^{12}) \Le C$ of Lemma~\ref{lem: invariant},
    \begin{equation}
        \Bigl| \mathbb E\bigl[ \tilde{\mathcal V}_i(Y_1) - \tilde{\mathcal V}_i(\tilde Y_0(h)) \bigr] - \frac{h^2}{2}\, \mathbb E\bigl[ \nabla^2 \tilde{\mathcal V}_i(Y_0) : \mathcal C(Z_0) \bigr] \Bigr| \Le C h^3 \sum_{j=1}^{4} \norm{\nabla^j \tilde{\mathcal V}_i}_\infty .
        \label{eq: sglm local}
    \end{equation}
    By \eqref{eq: propagated derivatives} and \eqref{eq: time sums}, the remainders of \eqref{eq: sglm local} satisfy
    \begin{equation*}
        C h^3 \sum_{i=0}^{\infty} \sum_{j=1}^{4} \norm{\nabla^j \tilde{\mathcal V}_i}_\infty \Le C h^2 \Bigl( 1 + \log\frac1h + h^{-\frac12} \Bigr) \Le C h^{\frac32} .
    \end{equation*}
    With \eqref{eq: sglm split}, \eqref{eq: shifted bias} and \eqref{eq: sglm telescope},
    \begin{equation}
        \biggl| \pi_h(f) - \pi(f) - \frac{h^2}{2} \sum_{i=0}^{\infty} \mathbb E\bigl[ \nabla^2 \tilde{\mathcal V}_i(Y_0) : \mathcal C(Z_0) \bigr] \biggr| \Le C h^{\frac32} .
        \label{eq: sglm leading}
    \end{equation}
    By \eqref{eq: trace bound}, \eqref{eq: propagated derivatives} and \eqref{eq: time sums} at $j = 2$, the sum in \eqref{eq: sglm leading} satisfies
    \begin{equation*}
        \frac{h^2}{2} \sum_{i=0}^{\infty} \Bigl| \mathbb E\bigl[ \nabla^2 \tilde{\mathcal V}_i(Y_0) : \mathcal C(Z_0) \bigr] \Bigr| \Le \frac{\Sigma h^2}{2} \sum_{i=0}^{\infty} \norm{\nabla^2 \tilde{\mathcal V}_i}_\infty \Le C \Sigma h .
    \end{equation*}
    Hence $|\pi_h(f) - \pi(f)| \Le C(\Sigma h + h^{\frac32})$ for every $f$ satisfying Assumption~\ref{ass: f} and $\norm{\nabla f}_\infty \Le 1$, and \eqref{eq: kr test functions} gives
    \begin{equation}
        \mathcal W_1(\pi_h, \pi) \Le C \bigl( \Sigma h + h^{\frac32} \bigr) \Le C h .
        \label{eq: w1 rough}
    \end{equation}

    In the coefficient of $\frac{h^2}{2}$ in \eqref{eq: sglm leading}, the expectation can be replaced by an integral against $\pi$. Since $Z_0 - Y_0 = \sqrt{h/2}\, \chi_0$, \eqref{eq: trace bound} gives
    \begin{equation}
        \Bigl| \mathbb E\bigl[ \nabla^2 \tilde{\mathcal V}_i(Y_0) : \mathcal C(Z_0) \bigr] - \pi_h\bigl( \nabla^2 \tilde{\mathcal V}_i : \mathcal C \bigr) \Bigr| = \Bigl| \mathbb E\bigl[ \bigl( \nabla^2 \tilde{\mathcal V}_i(Y_0) - \nabla^2 \tilde{\mathcal V}_i(Z_0) \bigr) : \mathcal C(Z_0) \bigr] \Bigr| \Le C \sqrt h\, \norm{\nabla^3 \tilde{\mathcal V}_i}_\infty .
        \label{eq: sglm point replacement}
    \end{equation}
    By \eqref{eq: trace bound}, \eqref{eq: covariance lipschitz} and $|A : B| \Le d \norm{A} \norm{B}$ for $A, B \in \mathbb R^{d \times d}$, the function $\nabla^2 \tilde{\mathcal V}_i : \mathcal C$ is Lipschitz,
    \begin{equation*}
        \bigl| \nabla^2 \tilde{\mathcal V}_i(x) : \mathcal C(x) - \nabla^2 \tilde{\mathcal V}_i(x') : \mathcal C(x') \bigr| \Le \bigl( \Sigma \norm{\nabla^3 \tilde{\mathcal V}_i}_\infty + 4 d M^2 \norm{\nabla^2 \tilde{\mathcal V}_i}_\infty \bigr) |x - x'| ,
    \end{equation*}
    so the Kantorovich--Rubinstein duality and \eqref{eq: w1 rough} give
    \begin{equation}
        \bigl| \pi_h(\nabla^2 \tilde{\mathcal V}_i : \mathcal C) - \pi(\nabla^2 \tilde{\mathcal V}_i : \mathcal C) \bigr| \Le C h \bigl( \norm{\nabla^2 \tilde{\mathcal V}_i}_\infty + \norm{\nabla^3 \tilde{\mathcal V}_i}_\infty \bigr) .
        \label{eq: sglm measure replacement}
    \end{equation}
    By \eqref{eq: sglm point replacement}, \eqref{eq: sglm measure replacement}, \eqref{eq: propagated derivatives} and \eqref{eq: time sums} at $j = 2, 3$, the replacements of $Y_0$ by $Z_0$ and of $\pi_h$ by $\pi$ in the coefficient of $\frac{h^2}{2}$ in \eqref{eq: sglm leading} introduce an additional error of at most
    \begin{equation*}
    \begin{aligned}
        & \frac{h^2}{2} \sum_{i=0}^{\infty} \Bigl( C \sqrt h\, \norm{\nabla^3 \tilde{\mathcal V}_i}_\infty + C h \bigl( \norm{\nabla^2 \tilde{\mathcal V}_i}_\infty + \norm{\nabla^3 \tilde{\mathcal V}_i}_\infty \bigr) \Bigr) \\
        & \quad \Le C h^{\frac32} \Bigl( 1 + \log\frac1h \Bigr) + C h^2 \Bigl( 1 + \log\frac1h \Bigr) \Le C h^{\frac32} \Bigl( 1 + \log\frac1h \Bigr) ,
    \end{aligned}
    \end{equation*}
    and \eqref{eq: sglm leading} gives
    \begin{equation}
        \biggl| \pi_h(f) - \pi(f) - \frac{h^2}{2} \sum_{i=0}^{\infty} \pi\bigl( \nabla^2 \tilde{\mathcal V}_i : \mathcal C \bigr) \biggr| \Le C h^{\frac32} \Bigl( 1 + \log\frac1h \Bigr) .
        \label{eq: sglm under pi}
    \end{equation}

    Since $\tilde{\mathcal V}_i \in \mathcal D$, \eqref{eq: riesz} with $\mathcal A = \mathcal C$ and $\omega = \tilde{\mathcal V}_i$ gives
    \begin{equation}
        \pi\bigl( \nabla^2 \tilde{\mathcal V}_i : \mathcal C \bigr) = \pi\bigl( \psi_{\mathcal C}\, (-\mathcal L \tilde{\mathcal V}_i) \bigr) .
        \label{eq: sglm riesz}
    \end{equation}
    The generators \eqref{eq: generator} and \eqref{eq: modified generator} differ only in the drift, so
    \begin{equation}
        -\mathcal L \tilde{\mathcal V}_i = -\tilde{\mathcal L}_h \tilde{\mathcal V}_i - \bigl( \nabla \tilde U_h - \nabla U \bigr) \cdot \nabla \tilde{\mathcal V}_i ,
        \label{eq: generator difference}
    \end{equation}
    and \eqref{eq: modified class} bounds the drift term of \eqref{eq: generator difference}:
    \begin{equation*}
        \bigl| \bigl( \nabla \tilde U_h(x) - \nabla U(x) \bigr) \cdot \nabla \tilde{\mathcal V}_i(x) \bigr| \Le C h \bigl( 1 + |x| \bigr) \norm{\nabla \tilde{\mathcal V}_i}_\infty , \qquad x \in \mathbb R^d .
    \end{equation*}
    Inserting \eqref{eq: generator difference} into \eqref{eq: sglm riesz} and using the bound on $\pi\bigl( |\psi_{\mathcal C}|\, (1 + |x|)^2 \bigr)$ in \eqref{eq: psi moments} gives
    \begin{equation*}
        \Bigl| \pi\bigl( \nabla^2 \tilde{\mathcal V}_i : \mathcal C \bigr) - \pi\bigl( \psi_{\mathcal C}\, (-\tilde{\mathcal L}_h \tilde{\mathcal V}_i) \bigr) \Bigr| \Le C h\, \pi\bigl( |\psi_{\mathcal C}|\, (1 + |x|) \bigr) \norm{\nabla \tilde{\mathcal V}_i}_\infty \Le C h \norm{\nabla \tilde{\mathcal V}_i}_\infty .
    \end{equation*}
    By \eqref{eq: propagated derivatives} and \eqref{eq: time sums} at $j = 1$, $h \sum_{i=0}^{\infty} \norm{\nabla \tilde{\mathcal V}_i}_\infty \Le C$, so
    \begin{equation*}
        \frac{h^2}{2} \sum_{i=0}^{\infty} \Bigl| \pi\bigl( \nabla^2 \tilde{\mathcal V}_i : \mathcal C \bigr) - \pi\bigl( \psi_{\mathcal C}\, (-\tilde{\mathcal L}_h \tilde{\mathcal V}_i) \bigr) \Bigr| \Le C h^3 \sum_{i=0}^{\infty} \norm{\nabla \tilde{\mathcal V}_i}_\infty \Le C h^2 ,
    \end{equation*}
    and \eqref{eq: sglm under pi} gives
    \begin{equation}
        \biggl| \pi_h(f) - \pi(f) - \frac{h^2}{2} \sum_{i=0}^{\infty} \pi\bigl( \psi_{\mathcal C}\, (-\tilde{\mathcal L}_h \tilde{\mathcal V}_i) \bigr) \biggr| \Le C h^{\frac32} \Bigl( 1 + \log\frac1h \Bigr) .
        \label{eq: sglm modified generator}
    \end{equation}

    It remains to sum $h (-\tilde{\mathcal L}_h \tilde{\mathcal V}_i)$ over $i$. For this sum, fix a deterministic initial point $Y_0 \in \mathbb R^d$. By \eqref{eq: modified dynkin}, with $\tilde{\mathcal V}_i$ and $\tilde{\mathcal L}_h \tilde{\mathcal V}_i$ in place of $\tilde\phi_h$ and $\tilde g_h$, which applies since $\nabla \tilde{\mathcal V}_i$ is bounded, and $e^{h\tilde{\mathcal L}_h} \tilde{\mathcal V}_i = \tilde{\mathcal V}_{i+1}$,
    \begin{equation*}
        \tilde{\mathcal V}_i(Y_0) - \tilde{\mathcal V}_{i+1}(Y_0) = \int_0^h \mathbb E\bigl[ (-\tilde{\mathcal L}_h \tilde{\mathcal V}_i)(\tilde Y_0(s)) \bigr] \D s ,
    \end{equation*}
    hence
    \begin{equation}
        h\, (-\tilde{\mathcal L}_h \tilde{\mathcal V}_i)(Y_0) - \bigl( \tilde{\mathcal V}_i(Y_0) - \tilde{\mathcal V}_{i+1}(Y_0) \bigr) = \int_0^h \mathbb E\bigl[ (\tilde{\mathcal L}_h \tilde{\mathcal V}_i)(\tilde Y_0(s)) - (\tilde{\mathcal L}_h \tilde{\mathcal V}_i)(Y_0) \bigr] \D s .
        \label{eq: dynkin error}
    \end{equation}
    Differentiating \eqref{eq: modified generator} gives
    \begin{equation*}
        \nabla (\tilde{\mathcal L}_h \tilde{\mathcal V}_i) = \nabla \Delta \tilde{\mathcal V}_i - \nabla^2 \tilde U_h\, \nabla \tilde{\mathcal V}_i - \nabla^2 \tilde{\mathcal V}_i\, \nabla \tilde U_h ,
    \end{equation*}
    and \eqref{eq: modified class bounds} gives
    \begin{equation*}
        \norm{\nabla^2 \tilde U_h(z)} \Le 2M , \qquad |\nabla \tilde U_h(z)| \Le 2M (1 + |z|) , \qquad z \in \mathbb R^d .
    \end{equation*}
    Hence, for $z \in \mathbb R^d$ and $0 \Le s \Le h$,
    \begin{equation*}
    \begin{aligned}
        & \bigl| \nabla (\tilde{\mathcal L}_h \tilde{\mathcal V}_i)(z) \bigr| \Le C (1 + |z|) \sum_{j=1}^{3} \norm{\nabla^j \tilde{\mathcal V}_i}_\infty , \\[-8pt]
        & \bigl| (\tilde{\mathcal L}_h \tilde{\mathcal V}_i)(\tilde Y_0(s)) - (\tilde{\mathcal L}_h \tilde{\mathcal V}_i)(Y_0) \bigr| \Le C \bigl( 1 + |Y_0| + |\tilde Y_0(s) - Y_0| \bigr) |\tilde Y_0(s) - Y_0| \sum_{j=1}^{3} \norm{\nabla^j \tilde{\mathcal V}_i}_\infty .
    \end{aligned}
    \end{equation*}
    The bound on $\tilde Y_0(s) - Y_0$ in \eqref{eq: exact sizes} with $p = 2$ gives
    \begin{equation*}
        \mathbb E\bigl[ (1 + |Y_0| + |\tilde Y_0(s) - Y_0|)\, |\tilde Y_0(s) - Y_0| \bigr] \Le C \sqrt h\, (1+|Y_0|)^2 , \qquad 0 \Le s \Le h ,
    \end{equation*}
    so integrating over $s \in [0,h]$ in \eqref{eq: dynkin error} gives, for every $Y_0 \in \mathbb R^d$,
    \begin{equation*}
        \Bigl| h\, (-\tilde{\mathcal L}_h \tilde{\mathcal V}_i)(Y_0) - \bigl( \tilde{\mathcal V}_i(Y_0) - \tilde{\mathcal V}_{i+1}(Y_0) \bigr) \Bigr| \Le C h^{\frac32} \bigl( 1 + |Y_0| \bigr)^2 \sum_{j=1}^{3} \norm{\nabla^j \tilde{\mathcal V}_i}_\infty .
    \end{equation*}
    By \eqref{eq: psi moments}, \eqref{eq: propagated derivatives} and \eqref{eq: time sums} at $j = 1, 2, 3$,
    \begin{equation*}
        \frac h2 \sum_{i=0}^{\infty} \Bigl| \pi\Bigl( \psi_{\mathcal C}\, \bigl( h\, (-\tilde{\mathcal L}_h \tilde{\mathcal V}_i) - \tilde{\mathcal V}_i + \tilde{\mathcal V}_{i+1} \bigr) \Bigr) \Bigr| \Le C h^{\frac32} \Bigl( 1 + \log\frac1h \Bigr) .
    \end{equation*}
    The partial sums of the telescoping series are
    \begin{equation*}
        \sum_{i=0}^{N-1} \pi\bigl( \psi_{\mathcal C}\, (\tilde{\mathcal V}_i - \tilde{\mathcal V}_{i+1}) \bigr) = \pi\bigl( \psi_{\mathcal C} \tilde f_h \bigr) - \pi\bigl( \psi_{\mathcal C} \tilde{\mathcal V}_N \bigr) ,
    \end{equation*}
    and $\pi(\psi_{\mathcal C}) = 0$, \eqref{eq: sglm relaxation} and \eqref{eq: psi moments} give $\pi(\psi_{\mathcal C} \tilde{\mathcal V}_N) \to 0$ as $N \to \infty$. The series of $\pi(\psi_{\mathcal C}\, (-\tilde{\mathcal L}_h \tilde{\mathcal V}_i))$ converges absolutely: by $|\tilde{\mathcal L}_h \tilde{\mathcal V}_i(x)| \Le C (1+|x|) (\norm{\nabla \tilde{\mathcal V}_i}_\infty + \norm{\nabla^2 \tilde{\mathcal V}_i}_\infty)$, \eqref{eq: psi moments} and \eqref{eq: propagated derivatives},
    \begin{equation*}
        \sum_{i=0}^{\infty} \Bigl| \pi\bigl( \psi_{\mathcal C}\, (-\tilde{\mathcal L}_h \tilde{\mathcal V}_i) \bigr) \Bigr| \Le C \sum_{i=0}^{\infty} \bigl( \norm{\nabla \tilde{\mathcal V}_i}_\infty + \norm{\nabla^2 \tilde{\mathcal V}_i}_\infty \bigr) \Le C h^{-\frac12} \sum_{i=0}^{\infty} e^{-\lambda i h} < \infty .
    \end{equation*}
    Therefore
    \begin{equation}
        \biggl| \frac{h^2}{2} \sum_{i=0}^{\infty} \pi\bigl( \psi_{\mathcal C}\, (-\tilde{\mathcal L}_h \tilde{\mathcal V}_i) \bigr) - \frac h2\, \pi\bigl( \psi_{\mathcal C} \tilde f_h \bigr) \biggr| \Le C h^{\frac32} \Bigl( 1 + \log\frac1h \Bigr) .
        \label{eq: sglm telescope psi}
    \end{equation}
    Finally, \eqref{eq: heat shift} and $\norm{\nabla f}_\infty \Le 1$ give $|\tilde f_h - f| \Le \sqrt{h/2}\; \mathbb E|\chi| \Le \sqrt{d h / 2}$, so $\frac h2\, |\pi(\psi_{\mathcal C}\, (\tilde f_h - f))| \Le C h^{\frac32}$ by \eqref{eq: psi moments}. With \eqref{eq: sglm modified generator} and \eqref{eq: sglm telescope psi},
    \begin{equation*}
        \Bigl| \pi_h(f) - \pi(f) - \frac h2\, \pi( f \psi_{\mathcal C} ) \Bigr| \Le C h^{\frac32} \Bigl( 1 + \log\frac1h \Bigr) .
    \end{equation*}
    Since $\pi_h(f) - \pi(f) - \frac h2\, \pi(f \psi_{\mathcal C}) = \int f\, \D(\pi_h - \pi - \frac h2\, \psi_{\mathcal C}\, \pi)$, the supremum over the functions $f$ satisfying Assumption~\ref{ass: f} and $\norm{\nabla f}_\infty \Le 1$, with \eqref{eq: kr test functions}, proves \eqref{eq: kr sglm}.
\end{proof}

\begin{proof}[Proof of \eqref{eq: kr sgld} in Theorem~\ref{thm: kr expansion}]
    The proof follows that of \eqref{eq: kr sglm}, with the local error \eqref{eq: sgld weak} in place of \eqref{eq: sg weak}, Lemma~\ref{lem: sgld lipschitz step} for the term $i = 0$, and Lemma~\ref{lem: local lipschitz} for the move from $\bar\pi_h$ to $\pi$. Let $f$ satisfy Assumption~\ref{ass: f} and $\norm{\nabla f}_\infty \Le 1$, and let
    \begin{equation}
        \mathcal V_i := e^{ih\mathcal L} f , \qquad i = 0, 1, 2, \dots ,
        \label{eq: sgld test functions}
    \end{equation}
    so that $\mathcal V_0 = f$ and $\mathcal V_i \in \mathcal D$ by \eqref{eq: short time}. By \eqref{eq: smoothing} at $t = ih$, as for \eqref{eq: propagated derivatives}, for $i \Ge 1$ and $j = 1, 2, 3, 4$,
    \begin{equation}
        \norm{\nabla^j \mathcal V_i}_\infty \Le C \min\bigl( 1, (i+1)h \bigr)^{-\frac{j-1}{2}} e^{-\lambda i h} .
        \label{eq: sgld propagated}
    \end{equation}
    With the transition kernel $\mathcal P_h$ of \eqref{eq: sgld}, as in Appendix~\ref{app: sgld}, $(\mathcal P_h g)(Z_0) = \hat{\mathbb E}_0\bigl[ g(Z_1) \bigr]$ for a given $Z_0$. Since $\bar\pi_h$ is invariant under $\mathcal P_h$ and $e^{h\mathcal L} \mathcal V_i = \mathcal V_{i+1}$,
    \begin{equation*}
        \bar\pi_h(\mathcal V_i) - \bar\pi_h(\mathcal V_{i+1}) = \bar\pi_h\bigl( \mathcal P_h \mathcal V_i - e^{h\mathcal L} \mathcal V_i \bigr) ,
    \end{equation*}
    and $|\mathcal V_N(x) - \pi(f)| \Le C e^{-\lambda N h} (1 + |x|)$ by \eqref{eq: smoothing} at $j = 1$ and the invariance of $\pi$ under $e^{t\mathcal L}$, as in \eqref{eq: sglm relaxation}. Summing over $i$,
    \begin{equation}
        \bar\pi_h(f) - \pi(f) = \sum_{i=0}^{\infty} \bar\pi_h\bigl( \mathcal P_h \mathcal V_i - e^{h\mathcal L} \mathcal V_i \bigr) .
        \label{eq: sgld telescope}
    \end{equation}
    For the term $i = 0$ of \eqref{eq: sgld telescope}, averaging \eqref{eq: sgld lipschitz step} over $Z_0 \sim \bar\pi_h$, with $\bar\pi_h(|x|) \Le C$ from Lemma~\ref{lem: sgld invariant}, gives
    \begin{equation*}
        \bigl| \bar\pi_h\bigl( \mathcal P_h f - e^{h\mathcal L} f \bigr) \bigr| \Le C h^{\frac32} \bigl( 1 + \bar\pi_h(|x|) \bigr) \Le C h^{\frac32} .
    \end{equation*}
    For $i \Ge 1$, averaging \eqref{eq: sgld weak} at $k = 0$, with $G = \mathcal V_i$, over $Z_0 \sim \bar\pi_h$, with $\bar\pi_h(|x|^{12}) \Le C$, gives
    \begin{equation*}
        \biggl| \bar\pi_h\bigl( \mathcal P_h \mathcal V_i - e^{h\mathcal L} \mathcal V_i \bigr) - \frac{h^2}{2}\, \bar\pi_h\Bigl( \nabla^2 \mathcal V_i : \bigl( 2 \nabla^2 U + \mathcal C \bigr) - \nabla \mathcal V_i \cdot \bigl( \nabla^2 U\, \nabla U - \nabla \Delta U \bigr) \Bigr) \biggr| \Le C h^3 \sum_{j=1}^{4} \norm{\nabla^j \mathcal V_i}_\infty ,
    \end{equation*}
    and, by \eqref{eq: sgld propagated} and \eqref{eq: time sums},
    \begin{equation*}
        C h^3 \sum_{i=1}^{\infty} \sum_{j=1}^{4} \norm{\nabla^j \mathcal V_i}_\infty \Le C h^2 \Bigl( 1 + \log\frac1h + h^{-\frac12} \Bigr) \Le C h^{\frac32} .
    \end{equation*}
    Hence
    \begin{equation}
        \biggl| \bar\pi_h(f) - \pi(f) - \frac{h^2}{2} \sum_{i=1}^{\infty} \bar\pi_h\Bigl( \nabla^2 \mathcal V_i : \bigl( 2 \nabla^2 U + \mathcal C \bigr) - \nabla \mathcal V_i \cdot \bigl( \nabla^2 U\, \nabla U - \nabla \Delta U \bigr) \Bigr) \biggr| \Le C h^{\frac32} .
        \label{eq: sgld leading}
    \end{equation}
    By Assumption~\ref{ass: U}, \eqref{eq: trace bound} and $|A : B| \Le d \norm{A} \norm{B}$ for $A, B \in \mathbb R^{d \times d}$, for every $x \in \mathbb R^d$ and every symmetric matrix $A$,
    \begin{equation*}
    \begin{aligned}
        & \bigl| \nabla^2 U(x)\, \nabla U(x) - \nabla \Delta U(x) \bigr| \Le M^2 (1 + |x|) + d M \Le C (1 + |x|) , \\
        & \bigl| A : \bigl( 2 \nabla^2 U(x) + \mathcal C(x) \bigr) \bigr| \Le (2 d M + \Sigma) \norm{A} .
    \end{aligned}
    \end{equation*}
    Hence the integrand in \eqref{eq: sgld leading} satisfies
    \begin{equation*}
        \Bigl| \nabla^2 \mathcal V_i : \bigl( 2 \nabla^2 U + \mathcal C \bigr) - \nabla \mathcal V_i \cdot \bigl( \nabla^2 U\, \nabla U - \nabla \Delta U \bigr) \Bigr| \Le C \bigl( 1 + |x| \bigr) \bigl( \norm{\nabla \mathcal V_i}_\infty + \norm{\nabla^2 \mathcal V_i}_\infty \bigr) ,
    \end{equation*}
    and the moment bound $\bar\pi_h(|x|) \Le C$ of Lemma~\ref{lem: sgld invariant}, \eqref{eq: sgld propagated} and \eqref{eq: time sums} at $j = 1, 2$ give
    \begin{equation*}
        \frac{h^2}{2} \sum_{i=1}^{\infty} \biggl| \bar\pi_h\Bigl( \nabla^2 \mathcal V_i : \bigl( 2 \nabla^2 U + \mathcal C \bigr) - \nabla \mathcal V_i \cdot \bigl( \nabla^2 U\, \nabla U - \nabla \Delta U \bigr) \Bigr) \biggr| \Le C h^2 \sum_{i=1}^{\infty} \bigl( \norm{\nabla \mathcal V_i}_\infty + \norm{\nabla^2 \mathcal V_i}_\infty \bigr) \Le C h .
    \end{equation*}
    With \eqref{eq: sgld leading}, $|\bar\pi_h(f) - \pi(f)| \Le C (h + h^{\frac32}) \Le C h$ for every $f$ satisfying Assumption~\ref{ass: f} and $\norm{\nabla f}_\infty \Le 1$, and \eqref{eq: kr test functions} gives
    \begin{equation}
        \mathcal W_1(\bar\pi_h, \pi) \Le C h .
        \label{eq: w1 rough sgld}
    \end{equation}

    In the coefficient of $\frac{h^2}{2}$ in \eqref{eq: sgld leading}, the measure $\bar\pi_h$ can be replaced by $\pi$. By Assumption~\ref{ass: U}, the derivatives of $U$ of orders $2$ to $7$ are bounded, so $\nabla^2 U$ and $\nabla \Delta U$ are bounded and Lipschitz, and
    \begin{equation}
        \bigl| \bigl( \nabla^2 U\, \nabla U - \nabla \Delta U \bigr)(x) - \bigl( \nabla^2 U\, \nabla U - \nabla \Delta U \bigr)(x') \bigr| \Le C \bigl( 1 + |x| + |x'| \bigr) |x - x'| , \qquad x, x' \in \mathbb R^d ,
        \label{eq: drift correction lipschitz}
    \end{equation}
    where the factor $1 + |x| + |x'|$ comes entirely from the linear growth of $\nabla U$. With \eqref{eq: drift correction lipschitz}, \eqref{eq: trace bound}, \eqref{eq: covariance lipschitz} and $|A : B| \Le d \norm{A} \norm{B}$, for each $i$, the integrand in the $\bar\pi_h$ bracket of \eqref{eq: sgld leading} satisfies \eqref{eq: local lipschitz} with $K = C \sum_{j=1}^{3} \norm{\nabla^j \mathcal V_i}_\infty$. Lemma~\ref{lem: local lipschitz}, with $\mu = \bar\pi_h$, $\nu = \pi$, the moments of Lemmas~\ref{lem: sgld invariant} and \ref{lem: dissipative}, and $\mathcal W_1(\bar\pi_h, \pi) \Le C h \Le 1$ by \eqref{eq: w1 rough sgld} with $C h_0 \Le 1$, gives, for $i \Ge 1$,
    \begin{equation}
    \begin{aligned}
        & \Bigl| (\bar\pi_h - \pi)\Bigl( \nabla^2 \mathcal V_i : \bigl( 2 \nabla^2 U + \mathcal C \bigr) - \nabla \mathcal V_i \cdot \bigl( \nabla^2 U\, \nabla U - \nabla \Delta U \bigr) \Bigr) \Bigr| \\
        & \quad \Le C\, \mathcal W_1(\bar\pi_h, \pi)^{\frac{10}{11}} \sum_{j=1}^{3} \norm{\nabla^j \mathcal V_i}_\infty \Le C h^{\frac{10}{11}} \sum_{j=1}^{3} \norm{\nabla^j \mathcal V_i}_\infty .
    \end{aligned}
    \label{eq: sgld measure replacement}
    \end{equation}
    By \eqref{eq: sgld propagated} and \eqref{eq: time sums} at $j = 1, 2, 3$, $h \sum_{i=1}^{\infty} \sum_{j=1}^{3} \norm{\nabla^j \mathcal V_i}_\infty \Le C (1 + \log\frac1h)$, so, by \eqref{eq: sgld measure replacement}, replacing $\bar\pi_h$ by $\pi$ in the coefficient of $\frac{h^2}{2}$ in \eqref{eq: sgld leading} introduces an additional error of at most
    \begin{equation*}
        \frac{h^2}{2} \sum_{i=1}^{\infty} C h^{\frac{10}{11}} \sum_{j=1}^{3} \norm{\nabla^j \mathcal V_i}_\infty \Le C h^{1 + \frac{10}{11}} \Bigl( 1 + \log\frac1h \Bigr) \Le C h^{\frac32} .
    \end{equation*}
    Under $\pi$, \eqref{eq: sgld integration by parts} with $\mathcal V_i$ in place of $\psi$ and \eqref{eq: riesz} give
    \begin{equation*}
        \pi\Bigl( \nabla^2 \mathcal V_i : \bigl( 2 \nabla^2 U + \mathcal C \bigr) - \nabla \mathcal V_i \cdot \bigl( \nabla^2 U\, \nabla U - \nabla \Delta U \bigr) \Bigr) = \pi\bigl( \nabla^2 \mathcal V_i : ( \nabla^2 U + \mathcal C ) \bigr) = \pi\bigl( ( \psi_{\nabla^2 U} + \psi_{\mathcal C} )\, (-\mathcal L \mathcal V_i) \bigr) ,
    \end{equation*}
    so we obtain
    \begin{equation}
        \biggl| \bar\pi_h(f) - \pi(f) - \frac{h^2}{2} \sum_{i=1}^{\infty} \pi\bigl( ( \psi_{\nabla^2 U} + \psi_{\mathcal C} )\, (-\mathcal L \mathcal V_i) \bigr) \biggr| \Le C h^{\frac32} .
        \label{eq: sgld under pi}
    \end{equation}

    The sum in \eqref{eq: sgld under pi} telescopes as in the proof of \eqref{eq: kr sglm}, with \eqref{eq: langevin}, $\mathcal L$, $\mathcal V_i = e^{ih\mathcal L} f$, the solution $X_s$ of \eqref{eq: langevin} and its initial value $X_0$ in place of \eqref{eq: modified langevin}, $\tilde{\mathcal L}_h$, $\tilde{\mathcal V}_i = e^{ih\tilde{\mathcal L}_h} \tilde f_h$, $\tilde Y_0(s)$ and $Y_0$, with $\psi_{\nabla^2 U} + \psi_{\mathcal C}$ in place of $\psi_{\mathcal C}$, and with \eqref{eq: sgld propagated} and Assumption~\ref{ass: U} in place of \eqref{eq: propagated derivatives} and \eqref{eq: modified class bounds}. The identity \eqref{eq: dynkin error} and the bound on its integral hold in this form, and since the sum starts at $i = 1$, the partial sums of the telescoping series are $\pi((\psi_{\nabla^2 U} + \psi_{\mathcal C})\, \mathcal V_1) - \pi((\psi_{\nabla^2 U} + \psi_{\mathcal C})\, \mathcal V_N)$, where $\pi((\psi_{\nabla^2 U} + \psi_{\mathcal C})\, \mathcal V_N) \to 0$ as $N \to \infty$ by $\pi(\psi_{\nabla^2 U} + \psi_{\mathcal C}) = 0$, the bound $|\mathcal V_N(x) - \pi(f)| \Le C e^{-\lambda N h} (1 + |x|)$ and \eqref{eq: psi moments}. Hence
    \begin{equation*}
        \biggl| \frac{h^2}{2} \sum_{i=1}^{\infty} \pi\bigl( ( \psi_{\nabla^2 U} + \psi_{\mathcal C} )\, (-\mathcal L \mathcal V_i) \bigr) - \frac h2\, \pi\bigl( ( \psi_{\nabla^2 U} + \psi_{\mathcal C} )\, \mathcal V_1 \bigr) \biggr| \Le C h^{\frac32} \Bigl( 1 + \log\frac1h \Bigr) .
    \end{equation*}
    Finally, the bound on $X_s - X_0$ in \eqref{eq: exact sizes} for \eqref{eq: langevin} gives $|\mathcal V_1(x) - f(x)| = |\mathbb E^x[f(X_h) - f(X_0)]| \Le \mathbb E^x |X_h - X_0| \Le C \sqrt h\, (1 + |x|)$, so $\frac h2\, |\pi((\psi_{\nabla^2 U} + \psi_{\mathcal C})\, (\mathcal V_1 - f))| \Le C h^{\frac32}$ by \eqref{eq: psi moments}. With \eqref{eq: sgld under pi},
    \begin{equation*}
        \Bigl| \bar\pi_h(f) - \pi(f) - \frac h2\, \pi\bigl( f\, ( \psi_{\nabla^2 U} + \psi_{\mathcal C} ) \bigr) \Bigr| \Le C h^{\frac32} \Bigl( 1 + \log\frac1h \Bigr) ,
    \end{equation*}
    and the supremum over the functions $f$ satisfying Assumption~\ref{ass: f} and $\norm{\nabla f}_\infty \Le 1$, with \eqref{eq: kr test functions}, proves \eqref{eq: kr sgld}.
\end{proof}

\subsection{Proof of Corollary~\ref{cor: w1 limits}: limits of the Wasserstein-1 errors}
\label{sec: w1 limits proof}

\begin{proof}[Proof of Corollary~\ref{cor: w1 limits}]
    Both limits of Corollary~\ref{cor: w1 limits} follow from \eqref{eq: w1 limits} as in Section~\ref{subsec: w1}, since $\pi_h$, $\bar\pi_h$ and $\pi$ have finite first moments by Lemmas~\ref{lem: invariant}, \ref{lem: sgld invariant} and \ref{lem: dissipative}.

    For the positivity, a function $\psi \in L^2_0(\pi)$ with $\norm{\psi\, \pi}_{\mathrm{KR}} = 0$ has $\pi(g \psi) = 0$ for every $g \in C_c^\infty(\mathbb R^d)$, hence $\psi = 0$ $\pi$-almost everywhere. By \eqref{eq: psi hessian}, $\psi_{\nabla^2 U}$ is continuous and unbounded, since $|\Delta U| \Le d M$ while $|\nabla U(x)| \Ge \langle x, \nabla U(x) \rangle / |x| \to \infty$ as $|x| \to \infty$ by \eqref{eq: dissipative bounds}. Hence $\psi_{\nabla^2 U}$ is nonzero on a nonempty open set, which has positive $\pi$-measure since $\pi(x) > 0$ for every $x \in \mathbb R^d$, and $\norm{\psi_{\nabla^2 U}\, \pi}_{\mathrm{KR}} > 0$. For $\psi_{\mathcal C}$, let $r > 0$ and
    \begin{equation}
        \varphi_r(x) := \bigl( r^2 + |x|^2 \bigr)^{\frac12} , \qquad x \in \mathbb R^d ,
        \label{eq: phi r}
    \end{equation}
    so that $\varphi_r \in \mathcal D$ and
    \begin{equation}
    \begin{aligned}
        & |\nabla \varphi_r| \Le 1 , \qquad \nabla^2 \varphi_r = \frac{\varphi_r^2 I_d - x \otimes x}{\varphi_r^3} \succeq \frac{r^2 I_d}{\varphi_r^3} , \qquad \norm{\nabla^2 \varphi_r} \Le \frac{1}{\varphi_r} , \\
        & \Delta \varphi_r = \frac{d-1}{\varphi_r} + \frac{r^2}{\varphi_r^3} , \qquad |\nabla \Delta \varphi_r| = \frac{(d-1)\, |x|}{\varphi_r^3} + \frac{3 r^2 |x|}{\varphi_r^5} \Le \frac{d+1}{2 r^2} .
    \end{aligned}
        \label{eq: phi r bounds}
    \end{equation}
    Since $|\nabla U(x)| \Le M(1 + |x|) \Le \sqrt2\, M \varphi_1(x)$, the gradient of $\mathcal L \varphi_1 = \Delta \varphi_1 - \nabla U \cdot \nabla \varphi_1$ satisfies, by \eqref{eq: phi r bounds}, $|\nabla (\mathcal L \varphi_1)| \Le \frac{d+1}{2} + M + \sqrt2\, M \Le d + 3M$. With \eqref{eq: riesz} for $\omega = \varphi_1$ and $\mathcal C(x)$ positive semidefinite,
    \begin{equation*}
        \norm{\psi_{\mathcal C}\, \pi}_{\mathrm{KR}} \Ge \frac{\pi\bigl( \psi_{\mathcal C}\, (-\mathcal L \varphi_1) \bigr)}{d + 3M} = \frac{\pi\bigl( \nabla^2 \varphi_1 : \mathcal C \bigr)}{d + 3M} \Ge \frac{1}{d + 3M}\, \pi\biggl( \frac{\mathrm{Tr}\bigl[ \mathcal C(x) \bigr]}{(1 + |x|^2)^{\frac32}} \biggr) .
    \end{equation*}
    The trace $\mathrm{Tr}\bigl[ \mathcal C(x) \bigr]$ is nonnegative and continuous by \eqref{eq: covariance lipschitz}; if $\mathcal C(x)$ does not vanish identically, the trace is positive on a nonempty open set, and $\pi\Bigl( \mathrm{Tr}\bigl[ \mathcal C(x) \bigr]\, (1 + |x|^2)^{-\frac32} \Bigr) > 0$, so $\norm{\psi_{\mathcal C}\, \pi}_{\mathrm{KR}} > 0$.
\end{proof}


\subsection{Proof of Lemma~\ref{lem: w1 counterexample}: failure of the comparison}
\label{sec: w1 counterexample proof}

\begin{proof}[Proof of Lemma~\ref{lem: w1 counterexample}]
    The proof checks the assumptions, writes the Kantorovich--Rubinstein norm in dimension $1$ through primitives in \eqref{eq: kr one dimension}, computes $\psi_{\mathcal C}$ explicitly, and evaluates both norms of Lemma~\ref{lem: w1 counterexample} by quadrature. Assumption~\ref{ass: U} holds with $m = 1$, $R = 0$ and every $M \Ge 1$. The estimator is unbiased, $|\nabla \hat U(x;\xi)| \Le |x| + 10$, its Hessian $1 + 10\, \xi\, a'(x)$ is bounded by $1 + 10 \norm{a'}_\infty$, and its variance is $100\, a(x)^2 \Le 100$, so Assumption~\ref{ass: sg} holds for $M$ large enough.

    Here $\pi$ is the standard Gaussian distribution, with $\pi'(x) = -x\, \pi(x)$. For $\psi \in L^2_0(\pi)$ and $s \in \mathbb R$, the Cauchy--Schwarz inequality on $(-\infty, s]$ for $s \Le 0$, and on $[s, \infty)$ for $s > 0$ since $\pi(\psi) = 0$, gives
    \begin{equation*}
        \biggl| \int_{-\infty}^s \psi(y)\, \pi(y)\, \D y \biggr| \Le \pi(\psi^2)^{\frac12} e^{-\frac14 s^2} ,
    \end{equation*}
    so for $f$ Lipschitz with $|f'| \Le 1$ an integration by parts gives
    \begin{equation*}
        \pi(f \psi) = -\int_{\mathbb R} f'(s) \int_{-\infty}^s \psi(y)\, \pi(y)\, \D y \,\D s \Le \int_{\mathbb R} \biggl| \int_{-\infty}^s \psi(y)\, \pi(y)\, \D y \biggr| \,\D s ,
    \end{equation*}
    with equality for $f(s) = -\int_0^s \mathrm{sign}\bigl( \int_{-\infty}^{s'} \psi(y)\, \pi(y)\, \D y \bigr) \D s'$. Hence
    \begin{equation}
        \norm{\psi\, \pi}_{\mathrm{KR}} = \int_{\mathbb R} \biggl| \int_{-\infty}^s \psi(y)\, \pi(y)\, \D y \biggr| \,\D s .
        \label{eq: kr one dimension}
    \end{equation}

    By \eqref{eq: psi hessian}, $\psi_{\nabla^2 U}(x) = \frac12 (x^2 - 1)$, and
    \begin{equation}
        \int_{-\infty}^s \psi_{\nabla^2 U}(y)\, \pi(y)\, \D y = -\frac12\, s\, \pi(s) , \qquad s \in \mathbb R .
        \label{eq: G hessian}
    \end{equation}
    For $\mathcal A = a^2$, let $J(x) := \int_{-\infty}^x y\, a(y)^2\, \D y$ and $\zeta := \pi(a^2) - \pi(J)$. The function $J - a^2 + \zeta$ is smooth and bounded, satisfies $\pi(J - a^2 + \zeta) = 0$, and has the derivative $x\, a^2 - (a^2)'$, so for $\omega \in \mathcal D$, \eqref{eq: generator ibp} gives
    \begin{equation*}
        \pi\bigl( (J - a^2 + \zeta) (-\mathcal L \omega) \bigr) = \pi\bigl( (x\, a^2 - (a^2)')\, \omega' \bigr) = -\int_{\mathbb R} \omega'(x)\, \bigl( a(x)^2 \pi(x) \bigr)'\, \D x = \pi\bigl( \omega''\, a^2 \bigr) ,
    \end{equation*}
    which is \eqref{eq: riesz} for $\mathcal A = a^2$. By the uniqueness in Theorem~\ref{thm: riesz},
    \begin{equation}
        \psi_{a^2} = J - a^2 + \zeta .
        \label{eq: psi shells}
    \end{equation}
    Since $\mathcal C = 100\, a^2$, the linearity in Theorem~\ref{thm: riesz} gives $\psi_{\mathcal C} = 100\, \psi_{a^2}$. Quadrature in \eqref{eq: kr one dimension}, with \eqref{eq: G hessian} and \eqref{eq: psi shells}, then gives the values in Lemma~\ref{lem: w1 counterexample}, so \eqref{eq: w1 comparison} fails.
\end{proof}


\subsection{Proof of Theorem~\ref{thm: batch threshold}: threshold on the batch size}
\label{sec: batch threshold proof}

\begin{proof}[Proof of Theorem~\ref{thm: batch threshold}]
    The proof bounds $\norm{( \psi_{\nabla^2 U} + \psi_{\mathcal C} )\, \pi}_{\mathrm{KR}}$ from below by testing \eqref{eq: riesz} with a function of the distance to the minimizer of $U$, bounds $\norm{\psi_{\mathcal C}\, \pi}_{\mathrm{KR}}$ from above through \eqref{eq: bochner}, and compares the $2$ bounds. Let $x_*$ be the minimizer of $U$, and let $\varphi_r$ be the function of \eqref{eq: phi r} with $r^2 = d/m$ and $x - x_*$ in place of $x$. The gradient of $\mathcal L \varphi_r = \Delta \varphi_r - \nabla U \cdot \nabla \varphi_r$ satisfies
    \begin{equation}
        |\nabla (\mathcal L \varphi_r)| \Le 2\, \norm{\nabla^2 U}_\infty + \frac{d+1}{2r^2} \Le 3\, \norm{\nabla^2 U}_\infty ,
        \label{eq: lipschitz phi R}
    \end{equation}
    since \eqref{eq: phi r bounds} bounds $|\nabla \Delta \varphi_r|$ by $\frac{d+1}{2r^2} = \frac{(d+1)\, m}{2d} \Le m \Le \norm{\nabla^2 U}_\infty$, and, by $|\nabla \varphi_r| \Le 1$, $\norm{\nabla^2 \varphi_r} \Le \varphi_r^{-1}$ and $|\nabla U(x)| = |\nabla U(x) - \nabla U(x_*)| \Le \norm{\nabla^2 U}_\infty\, \varphi_r(x)$,
    \begin{equation*}
        |\nabla(\nabla U \cdot \nabla \varphi_r)| \Le \norm{\nabla^2 U}\, |\nabla \varphi_r| + \norm{\nabla^2 \varphi_r}\, |\nabla U| \Le 2 \norm{\nabla^2 U}_\infty .
    \end{equation*}

    By \eqref{eq: lipschitz phi R}, the function $-\mathcal L \varphi_r / (3\, \norm{\nabla^2 U}_\infty)$ is $1$-Lipschitz, and \eqref{eq: riesz} with $\omega = \varphi_r$ gives $\pi((-\mathcal L \varphi_r)\, \psi_{\mathcal A}) = \pi(\nabla^2 \varphi_r : \mathcal A)$ for $\mathcal A = \nabla^2 U$ and $\mathcal A = \mathcal C$. As $\nabla^2 \varphi_r$ and $\mathcal C(x)$ are positive semidefinite, $\pi(\nabla^2 \varphi_r : \mathcal C) \Ge 0$, and as $\nabla^2 U \succeq m I_d$, $\nabla^2 \varphi_r : \nabla^2 U \Ge m\, \Delta \varphi_r$. Hence
    \begin{equation}
        \norm{( \psi_{\nabla^2 U} + \psi_{\mathcal C} )\, \pi}_{\mathrm{KR}} \Ge \frac{\pi\bigl( (-\mathcal L \varphi_r)\, ( \psi_{\nabla^2 U} + \psi_{\mathcal C} ) \bigr)}{3\, \norm{\nabla^2 U}_\infty} \Ge \frac{m\, \pi( \Delta \varphi_r )}{3\, \norm{\nabla^2 U}_\infty} .
        \label{eq: threshold lower}
    \end{equation}
    The identity \eqref{eq: generator ibp} with $g_1 = 1$ and $g_2 = \frac12 |x - x_*|^2$ gives $\pi(\nabla U \cdot (x - x_*)) = d$, and $\nabla U(x) \cdot (x - x_*) \Ge m |x - x_*|^2$, so $\pi(|x - x_*|^2) \Le d/m$. As functions of $|x - x_*|^2$, $(r^2 + |x - x_*|^2)^{-\frac12}$ and $(r^2 + |x - x_*|^2)^{-\frac32}$ are convex and decreasing, so Jensen's inequality and $r^2 = d/m$ give
    \begin{equation*}
        \pi( \Delta \varphi_r ) \Ge \frac{d-1}{(2d/m)^{\frac12}} + \frac{d/m}{(2d/m)^{\frac32}} = \Bigl( \frac{m}{2d} \Bigr)^{\frac12} \Bigl( d - \frac12 \Bigr) \Ge \Bigl( \frac{m d}{8} \Bigr)^{\frac12} ,
    \end{equation*}
    and \eqref{eq: threshold lower} gives
    \begin{equation}
        \norm{( \psi_{\nabla^2 U} + \psi_{\mathcal C} )\, \pi}_{\mathrm{KR}} \Ge \frac{m^{\frac32} \sqrt d}{6 \sqrt2\, \norm{\nabla^2 U}_\infty} \Ge \frac{m^{\frac32} \sqrt d}{9\, \norm{\nabla^2 U}_\infty} .
        \label{eq: threshold lower explicit}
    \end{equation}

    For $\norm{\psi_{\mathcal C}\, \pi}_{\mathrm{KR}}$ in \eqref{eq: w1 comparison}, let $f$ satisfy Assumption~\ref{ass: f} and $\norm{\nabla f}_\infty \Le 1$, and let $\phi_0$ solve \eqref{eq: continuous poisson}. By \eqref{eq: bochner} with $\omega = \phi_0 \in \mathcal D$ and $\nabla^2 U \succeq m I_d$, and \eqref{eq: generator ibp},
    \begin{equation}
        \pi\bigl( \nabla^2 \phi_0 : \nabla^2 \phi_0 \bigr) + m\, \pi\bigl( |\nabla \phi_0|^2 \bigr) \Le \pi\bigl( (f - \pi(f))^2 \bigr) = \pi\bigl( \nabla f \cdot \nabla \phi_0 \bigr) \Le \pi\bigl( |\nabla \phi_0|^2 \bigr)^{\frac12} .
        \label{eq: bochner phi0}
    \end{equation}
    By \eqref{eq: bochner phi0}, $\pi(|\nabla \phi_0|^2)^{\frac12} \Le m^{-1}$ and then $\pi(\nabla^2 \phi_0 : \nabla^2 \phi_0) \Le m^{-1}$. As for \eqref{eq: riesz functional bound}, \eqref{eq: riesz} with $\omega = \phi_0$ and the Cauchy--Schwarz inequality give $|\pi(f \psi_{\mathcal C})| = |\pi(\nabla^2 \phi_0 : \mathcal C)| \Le \sqrt d\, \norm{\mathcal C}_\infty\, m^{-\frac12}$, and the supremum over $f$ with \eqref{eq: kr test functions} gives
    \begin{equation}
        \norm{\psi_{\mathcal C}\, \pi}_{\mathrm{KR}} \Le \frac{\sqrt d\, \norm{\mathcal C}_\infty}{m^{\frac12}} .
        \label{eq: threshold upper}
    \end{equation}
    Hence, whenever $\norm{\mathcal C}_\infty \Le m^2 / (9\, \norm{\nabla^2 U}_\infty)$, \eqref{eq: threshold upper} and \eqref{eq: threshold lower explicit} give
    \begin{equation*}
        \norm{\psi_{\mathcal C}\, \pi}_{\mathrm{KR}} \Le \frac{\sqrt d\, \norm{\mathcal C}_\infty}{m^{\frac12}} \Le \frac{m^{\frac32} \sqrt d}{9\, \norm{\nabla^2 U}_\infty} \Le \norm{( \psi_{\nabla^2 U} + \psi_{\mathcal C} )\, \pi}_{\mathrm{KR}} ,
    \end{equation*}
    which is \eqref{eq: w1 comparison}. For SGLM and SGLD with a batch of $P$ independent stochastic gradients, the averaged estimator satisfies Assumption~\ref{ass: sg} with $\Sigma/P$ in place of $\Sigma$ (Section~\ref{subsec: mse analysis}), so Corollary~\ref{cor: w1 limits} applies with its covariance $\mathcal C(x)/P$, and $\norm{\mathcal C}_\infty / P \Le m^2 / (9\, \norm{\nabla^2 U}_\infty)$ exactly when $P$ satisfies \eqref{eq: batch threshold}. The representor of $\mathcal C(x)/P$ is $\psi_{\mathcal C}/P$ by the linearity in Theorem~\ref{thm: riesz}, and \eqref{eq: w1 comparison} with $\psi_{\mathcal C}/P$ in place of $\psi_{\mathcal C}$ is \eqref{eq: w1 comparison batch}.
\end{proof}
